% !TEX program = lualatex
% ---------------------------------------------------------------------------------------------
%  Readback note 23: Chains of definite theta lifts and the Gauss-sum operator at the level prime
%
%  Build:  python3 tools/rb.py note <slug>
%
%  One unit of a Lean artifact, rewritten as ordinary mathematics. The prose is hand-written. Lean names,
%  line ranges and counts come from data.tex; quoted Lean statements are copied from the artifact by the build.
%
%  Rules of the text: standard terminology; every symbol is explained where it first appears and listed again
%  in the notation at the end; no letter has two meanings; a proof is complete or says what it takes for granted;
%  Lean names appear only in the LEAN lines and in the concordance.
%
%  The signature rule (the build enforces it): a function, relation or predicate that the artifact defines is
%  written only through its notation macro, which the build makes from the spec (`notation`) with one slot for
%  each explicit argument of the Lean definition. No argument is fixed and dropped, also not "for brevity" and
%  not for one section. Every argument and every hypothesis of a stated declaration stands inside its
%  statement, or in a Setting that the statement cites by number; only implicit and instance arguments (the
%  ambient objects) may stand in the conventions.
% ---------------------------------------------------------------------------------------------
\documentclass[11pt,a4paper]{article}
\ifdefined\pdfvariable\pdfvariable suppressoptionalinfo \numexpr 1+2+4+8\relax\fi
\usepackage[a4paper,hmargin=35mm,top=22mm,bottom=27mm,footskip=12mm,nohead]{geometry}

\usepackage{fontspec}
\directlua{
  luaotfload.add_fallback("gfallback", {"STIXTwoMath-Regular.otf:mode=node;", "DejaVuSans.ttf:mode=node;"})
  luaotfload.add_fallback("gmonofallback", {"DejaVuSansMono.ttf:mode=node;", "STIXTwoMath-Regular.otf:mode=node;"})
}
\setmainfont{NewCM10}[Extension=.otf,UprightFont=*-Book,ItalicFont=*-BookItalic,
  BoldFont=*-Bold,BoldItalicFont=*-BoldItalic,SmallCapsFont=*-Book,SmallCapsFeatures={Letters=SmallCaps},
  RawFeature={fallback=gfallback}]
\setsansfont{NewCMSans10}[Extension=.otf,UprightFont=*-Book,ItalicFont=*-BookOblique,
  BoldFont=*-Bold,BoldItalicFont=*-BoldOblique,RawFeature={fallback=gfallback}]
\setmonofont{JetBrainsMono}[Extension=.otf,UprightFont=*-Regular,ItalicFont=*-Italic,
  BoldFont=*-SemiBold,BoldItalicFont=*-SemiBoldItalic,Scale=MatchLowercase,
  RawFeature={-calt;-liga;fallback=gmonofallback}]
\usepackage{amsmath}
\usepackage{amsthm}
\usepackage[math-style=TeX]{unicode-math}
\setmathfont{NewCMMath-Book.otf}
\emergencystretch=2.5em\relax

\usepackage{xcolor}
\usepackage{tikz}
\usetikzlibrary{calc,arrows.meta,positioning,backgrounds}
\usepackage[protrusion=true,expansion=true]{microtype}
\usepackage{enumitem}
\usepackage{needspace}
\usepackage{booktabs}
\usepackage{array}
\usepackage{longtable}
\usepackage{fancyhdr}
\usepackage{lastpage}
\usepackage{multicol}
\usepackage[hidelinks,pdfusetitle]{hyperref}

\definecolor{ink}{HTML}{0F172A}
\definecolor{inksoft}{HTML}{475569}
\definecolor{inkfaint}{HTML}{94A3B8}
\definecolor{rulegrey}{HTML}{CBD5E1}
\definecolor{arith}{HTML}{C2410C}
\definecolor{est}{HTML}{7C3AED}
\definecolor{ana}{HTML}{0891B2}
\definecolor{ok}{HTML}{047857}
\definecolor{leancol}{HTML}{3730A3}

\title{Chains of definite theta lifts and the Gauss-sum operator at the level prime}
\author{grh-compact}
\date{}

% ------------------------------------------------------------------------------------- macros
\newcommand{\leanfont}{\ttfamily\fontsize{8.6}{10.5}\selectfont\color{leancol}}
\newcommand{\lean}[1]{\csname LN@#1\endcsname}
\newcommand{\fillin}[1]{{\color{arith}\itshape[#1]}}  % a place of the template that is not written yet; the build reports every one that is left
\newcommand{\leanquote}[1]{\csname LQ@#1\endcsname}
% generated by rb.py note; do not edit
\newcommand{\ConcDecls}{78}
\newcommand{\Pin}{72fead10}
\newcommand{\PinText}{that of commit {\leanfont 72fead10}}
\newcommand{\ArtifactName}{grh-compact}
\expandafter\newcommand\csname LN@restrictToYofdiag\endcsname{\href{https://github.com/lsernau/grh-lean/blob/72fead102f663d5c7a407964a13dc0cf9d977207/grh-compact/src/New/Seesaw/LiftCharacters.lean\#L20-L33}{{\leanfont Data.\allowbreak{}GaugeLe.\allowbreak{}restrictToY\_of\_diag}}}
\expandafter\newcommand\csname LN@restrictToY\endcsname{\href{https://github.com/lsernau/grh-lean/blob/72fead102f663d5c7a407964a13dc0cf9d977207/grh-compact/src/New/Seesaw/LiftCharacters.lean\#L37-L39}{{\leanfont Data.\allowbreak{}GaugeLe.\allowbreak{}restrictToY}}}
\expandafter\newcommand\csname LN@weylNormSqrestrictToYle\endcsname{\href{https://github.com/lsernau/grh-lean/blob/72fead102f663d5c7a407964a13dc0cf9d977207/grh-compact/src/New/Seesaw/LiftCharacters.lean\#L41-L51}{{\leanfont Data.\allowbreak{}weylNormSq\_restrictToY\_le}}}
\expandafter\newcommand\csname LN@nonInducedPlaces\endcsname{\href{https://github.com/lsernau/grh-lean/blob/72fead102f663d5c7a407964a13dc0cf9d977207/grh-compact/src/New/Seesaw/LiftCharacters.lean\#L55-L55}{{\leanfont Data.\allowbreak{}nonInducedPlaces}}}
\expandafter\newcommand\csname LN@isSplitLevelnonInducedPlaces\endcsname{\href{https://github.com/lsernau/grh-lean/blob/72fead102f663d5c7a407964a13dc0cf9d977207/grh-compact/src/New/Seesaw/LiftCharacters.lean\#L57-L61}{{\leanfont Data.\allowbreak{}isSplitLevel\_nonInducedPlaces}}}
\expandafter\newcommand\csname LN@diagPointN\endcsname{\href{https://github.com/lsernau/grh-lean/blob/72fead102f663d5c7a407964a13dc0cf9d977207/grh-compact/src/New/Seesaw/LiftCharacters.lean\#L67-L67}{{\leanfont Data.\allowbreak{}diagPointN}}}
\expandafter\newcommand\csname LN@conjEmbedGL\endcsname{\href{https://github.com/lsernau/grh-lean/blob/72fead102f663d5c7a407964a13dc0cf9d977207/grh-compact/src/New/Seesaw/LiftCharacters.lean\#L71-L71}{{\leanfont Data.\allowbreak{}conjEmbedGL}}}
\expandafter\newcommand\csname LN@nonInducedRep\endcsname{\href{https://github.com/lsernau/grh-lean/blob/72fead102f663d5c7a407964a13dc0cf9d977207/grh-compact/src/New/Seesaw/LiftCharacters.lean\#L75-L76}{{\leanfont Data.\allowbreak{}nonInducedRep}}}
\expandafter\newcommand\csname LN@adeleExtensionmapdiagN\endcsname{\href{https://github.com/lsernau/grh-lean/blob/72fead102f663d5c7a407964a13dc0cf9d977207/grh-compact/src/New/Seesaw/LiftCharacters.lean\#L78-L82}{{\leanfont Data.\allowbreak{}adeleExtension\_map\_diagN}}}
\expandafter\newcommand\csname LN@lastColDatumconjnezeroiff\endcsname{\href{https://github.com/lsernau/grh-lean/blob/72fead102f663d5c7a407964a13dc0cf9d977207/grh-compact/src/New/Seesaw/LiftCharacters.lean\#L84-L95}{{\leanfont Data.\allowbreak{}lastColDatum\_conj\_ne\_zero\_iff}}}
\expandafter\newcommand\csname LN@atPlacesconjmemmirabolic\endcsname{\href{https://github.com/lsernau/grh-lean/blob/72fead102f663d5c7a407964a13dc0cf9d977207/grh-compact/src/New/Seesaw/LiftCharacters.lean\#L97-L101}{{\leanfont Data.\allowbreak{}atPlaces\_conj\_mem\_mirabolic}}}
\expandafter\newcommand\csname LN@memramifiedofmemramifiedColumns\endcsname{\href{https://github.com/lsernau/grh-lean/blob/72fead102f663d5c7a407964a13dc0cf9d977207/grh-compact/src/New/Seesaw/LiftCharacters.lean\#L103-L108}{{\leanfont Data.\allowbreak{}mem\_ramified\_of\_mem\_ramifiedColumns}}}
\expandafter\newcommand\csname LN@fourierOpmemcolumn\endcsname{\href{https://github.com/lsernau/grh-lean/blob/72fead102f663d5c7a407964a13dc0cf9d977207/grh-compact/src/New/Seesaw/LiftCharacters.lean\#L112-L123}{{\leanfont Data.\allowbreak{}fourierOp\_mem\_column}}}
\expandafter\newcommand\csname LN@existsindexDvdfourierOppow\endcsname{\href{https://github.com/lsernau/grh-lean/blob/72fead102f663d5c7a407964a13dc0cf9d977207/grh-compact/src/New/Seesaw/LiftCharacters.lean\#L128-L169}{{\leanfont Data.\allowbreak{}exists\_indexDvd\_fourierOp\_pow}}}
\expandafter\newcommand\csname LN@charUnramifiedAtrankTwistofNmem\endcsname{\href{https://github.com/lsernau/grh-lean/blob/72fead102f663d5c7a407964a13dc0cf9d977207/grh-compact/src/New/Seesaw/LiftCharacters.lean\#L171-L174}{{\leanfont Data.\allowbreak{}charUnramifiedAt\_rankTwist\_of\_N\_mem}}}
\expandafter\newcommand\csname LN@spaceCharrankTwist\endcsname{\href{https://github.com/lsernau/grh-lean/blob/72fead102f663d5c7a407964a13dc0cf9d977207/grh-compact/src/New/Seesaw/LiftCharacters.lean\#L178-L180}{{\leanfont Data.\allowbreak{}spaceChar\_rankTwist}}}
\expandafter\newcommand\csname LN@IsSplitLevel\endcsname{\href{https://github.com/lsernau/grh-lean/blob/72fead102f663d5c7a407964a13dc0cf9d977207/grh-compact/src/New/Seesaw/NamesAtSplitLevel.lean\#L25-L29}{{\leanfont IsSplitLevel}}}
\expandafter\newcommand\csname LQ@IsSplitLevel\endcsname{\leanquoteline{0}{class IsSplitLevel (P : Finset (FinPlace E)) (N : outParam ℕ) where}\leanquoteline{2}{star\_notMem : ∀ w ∈ P, star w ∉ P}\leanquoteline{2}{q\_eq : ∀ w ∈ P, w.q = N}\leanquoteline{2}{protected v\_natCast : ∀ w ∈ P, v (N : LocField w) = exp (-1)}\leanquoteline{2}{natCast\_mem\_iff w : ↑N ∈ w.1 ↔ w ∈ P ∨ star w ∈ P}}
\expandafter\newcommand\csname LN@residueOverN\endcsname{\href{https://github.com/lsernau/grh-lean/blob/72fead102f663d5c7a407964a13dc0cf9d977207/grh-compact/src/New/Seesaw/NamesAtSplitLevel.lean\#L33-L34}{{\leanfont residueOverN}}}
\expandafter\newcommand\csname LN@residueOverNapply\endcsname{\href{https://github.com/lsernau/grh-lean/blob/72fead102f663d5c7a407964a13dc0cf9d977207/grh-compact/src/New/Seesaw/NamesAtSplitLevel.lean\#L37-L38}{{\leanfont residueOverN\_apply}}}
\expandafter\newcommand\csname LN@idemmulresidueOverN\endcsname{\href{https://github.com/lsernau/grh-lean/blob/72fead102f663d5c7a407964a13dc0cf9d977207/grh-compact/src/New/Seesaw/NamesAtSplitLevel.lean\#L41-L41}{{\leanfont idem\_mul\_residueOverN}}}
\expandafter\newcommand\csname LN@diagN\endcsname{\href{https://github.com/lsernau/grh-lean/blob/72fead102f663d5c7a407964a13dc0cf9d977207/grh-compact/src/New/Seesaw/NamesAtSplitLevel.lean\#L48-L49}{{\leanfont IsSplitLevel.\allowbreak{}diagN}}}
\expandafter\newcommand\csname LN@residueEquiv\endcsname{\href{https://github.com/lsernau/grh-lean/blob/72fead102f663d5c7a407964a13dc0cf9d977207/grh-compact/src/New/Seesaw/NamesAtSplitLevel.lean\#L56-L57}{{\leanfont IsSplitLevel.\allowbreak{}residueEquiv}}}
\expandafter\newcommand\csname LN@reduction\endcsname{\href{https://github.com/lsernau/grh-lean/blob/72fead102f663d5c7a407964a13dc0cf9d977207/grh-compact/src/New/Seesaw/NamesAtSplitLevel.lean\#L61-L61}{{\leanfont IsSplitLevel.\allowbreak{}reduction}}}
\expandafter\newcommand\csname LN@reductioneqiff\endcsname{\href{https://github.com/lsernau/grh-lean/blob/72fead102f663d5c7a407964a13dc0cf9d977207/grh-compact/src/New/Seesaw/NamesAtSplitLevel.lean\#L64-L67}{{\leanfont IsSplitLevel.\allowbreak{}reduction\_eq\_iff}}}
\expandafter\newcommand\csname LN@inclGL\endcsname{\href{https://github.com/lsernau/grh-lean/blob/72fead102f663d5c7a407964a13dc0cf9d977207/grh-compact/src/New/Seesaw/NamesAtSplitLevel.lean\#L72-L73}{{\leanfont IsSplitLevel.\allowbreak{}inclGL}}}
\expandafter\newcommand\csname LN@reductionGL\endcsname{\href{https://github.com/lsernau/grh-lean/blob/72fead102f663d5c7a407964a13dc0cf9d977207/grh-compact/src/New/Seesaw/NamesAtSplitLevel.lean\#L77-L77}{{\leanfont IsSplitLevel.\allowbreak{}reductionGL}}}
\expandafter\newcommand\csname LN@liftGL\endcsname{\href{https://github.com/lsernau/grh-lean/blob/72fead102f663d5c7a407964a13dc0cf9d977207/grh-compact/src/New/Seesaw/NamesAtSplitLevel.lean\#L80-L81}{{\leanfont IsSplitLevel.\allowbreak{}liftGL}}}
\expandafter\newcommand\csname LN@diagNInv\endcsname{\href{https://github.com/lsernau/grh-lean/blob/72fead102f663d5c7a407964a13dc0cf9d977207/grh-compact/src/New/Seesaw/NamesAtSplitLevel.lean\#L87-L87}{{\leanfont IsSplitLevel.\allowbreak{}diagNInv}}}
\expandafter\newcommand\csname LN@psiN\endcsname{\href{https://github.com/lsernau/grh-lean/blob/72fead102f663d5c7a407964a13dc0cf9d977207/grh-compact/src/New/Seesaw/NamesAtSplitLevel.lean\#L91-L91}{{\leanfont IsSplitLevel.\allowbreak{}psiN}}}
\expandafter\newcommand\csname LN@lastBasisVector\endcsname{\href{https://github.com/lsernau/grh-lean/blob/72fead102f663d5c7a407964a13dc0cf9d977207/grh-compact/src/New/Seesaw/NamesAtSplitLevel.lean\#L94-L94}{{\leanfont IsSplitLevel.\allowbreak{}lastBasisVector}}}
\expandafter\newcommand\csname LN@existsGLlastcolandrow\endcsname{\href{https://github.com/lsernau/grh-lean/blob/72fead102f663d5c7a407964a13dc0cf9d977207/grh-compact/src/New/Seesaw/NamesAtSplitLevel.lean\#L98-L102}{{\leanfont IsSplitLevel.\allowbreak{}exists\_GL\_last\_col\_and\_row}}}
\expandafter\newcommand\csname LN@fourierOp\endcsname{\href{https://github.com/lsernau/grh-lean/blob/72fead102f663d5c7a407964a13dc0cf9d977207/grh-compact/src/New/Seesaw/NamesAtSplitLevel.lean\#L106-L107}{{\leanfont IsSplitLevel.\allowbreak{}fourierOp}}}
\expandafter\newcommand\csname LQ@fourierOp\endcsname{\leanquoteline{0}{def fourierOp := π.ρ diagNInv * ∑ y with IsUnit (y (.inr ())),}\leanquoteline{2}{psiN (y (.inr ())) • π.ρ (inclGL (liftGL (exists\_GL\_last\_col\_and\_row.choose y)))}}
\expandafter\newcommand\csname LN@mirabolic\endcsname{\href{https://github.com/lsernau/grh-lean/blob/72fead102f663d5c7a407964a13dc0cf9d977207/grh-compact/src/New/Seesaw/NamesAtSplitLevel.lean\#L111-L112}{{\leanfont IsSplitLevel.\allowbreak{}mirabolic}}}
\expandafter\newcommand\csname LN@ramified\endcsname{\href{https://github.com/lsernau/grh-lean/blob/72fead102f663d5c7a407964a13dc0cf9d977207/grh-compact/src/New/Seesaw/NamesAtSplitLevel.lean\#L116-L117}{{\leanfont IsSplitLevel.\allowbreak{}ramified}}}
\expandafter\newcommand\csname LN@liftFn\endcsname{\href{https://github.com/lsernau/grh-lean/blob/72fead102f663d5c7a407964a13dc0cf9d977207/grh-compact/src/New/Seesaw/NamesAtSplitLevel.lean\#L121-L123}{{\leanfont IsSplitLevel.\allowbreak{}liftFn}}}
\expandafter\newcommand\csname LN@lastRowDatum\endcsname{\href{https://github.com/lsernau/grh-lean/blob/72fead102f663d5c7a407964a13dc0cf9d977207/grh-compact/src/New/Seesaw/NamesAtSplitLevel.lean\#L127-L128}{{\leanfont IsSplitLevel.\allowbreak{}lastRowDatum}}}
\expandafter\newcommand\csname LN@lastColDatum\endcsname{\href{https://github.com/lsernau/grh-lean/blob/72fead102f663d5c7a407964a13dc0cf9d977207/grh-compact/src/New/Seesaw/NamesAtSplitLevel.lean\#L132-L132}{{\leanfont IsSplitLevel.\allowbreak{}lastColDatum}}}
\expandafter\newcommand\csname LN@charDatum\endcsname{\href{https://github.com/lsernau/grh-lean/blob/72fead102f663d5c7a407964a13dc0cf9d977207/grh-compact/src/New/Seesaw/NamesAtSplitLevel.lean\#L136-L136}{{\leanfont IsSplitLevel.\allowbreak{}charDatum}}}
\expandafter\newcommand\csname LN@LiftLaw\endcsname{\href{https://github.com/lsernau/grh-lean/blob/72fead102f663d5c7a407964a13dc0cf9d977207/grh-compact/src/New/Seesaw/NamesAtSplitLevel.lean\#L141-L146}{{\leanfont IsSplitLevel.\allowbreak{}LiftLaw}}}
\expandafter\newcommand\csname LN@lastRowDatumconju\endcsname{\href{https://github.com/lsernau/grh-lean/blob/72fead102f663d5c7a407964a13dc0cf9d977207/grh-compact/src/New/Seesaw/NamesAtSplitLevel.lean\#L161-L165}{{\leanfont IsSplitLevel.\allowbreak{}lastRowDatum\_conj\_u₀}}}
\expandafter\newcommand\csname LN@existsinclGLeq\endcsname{\href{https://github.com/lsernau/grh-lean/blob/72fead102f663d5c7a407964a13dc0cf9d977207/grh-compact/src/New/Seesaw/NamesAtSplitLevel.lean\#L167-L168}{{\leanfont IsSplitLevel.\allowbreak{}exists\_inclGL\_eq}}}
\expandafter\newcommand\csname LN@lastColDatumnezero\endcsname{\href{https://github.com/lsernau/grh-lean/blob/72fead102f663d5c7a407964a13dc0cf9d977207/grh-compact/src/New/Seesaw/NamesAtSplitLevel.lean\#L170-L171}{{\leanfont IsSplitLevel.\allowbreak{}lastColDatum\_ne\_zero}}}
\expandafter\newcommand\csname LN@mirabolicfixfourierOp\endcsname{\href{https://github.com/lsernau/grh-lean/blob/72fead102f663d5c7a407964a13dc0cf9d977207/grh-compact/src/New/Seesaw/NamesAtSplitLevel.lean\#L178-L188}{{\leanfont IsSplitLevel.\allowbreak{}mirabolic\_fix\_fourierOp}}}
\expandafter\newcommand\csname LN@liftFncompeqsumsmul\endcsname{\href{https://github.com/lsernau/grh-lean/blob/72fead102f663d5c7a407964a13dc0cf9d977207/grh-compact/src/New/Seesaw/NamesAtSplitLevel.lean\#L191-L193}{{\leanfont IsSplitLevel.\allowbreak{}liftFn\_comp\_eq\_sum\_smul}}}
\expandafter\newcommand\csname LN@eqzeroofsmulinvariant\endcsname{\href{https://github.com/lsernau/grh-lean/blob/72fead102f663d5c7a407964a13dc0cf9d977207/grh-compact/src/New/Seesaw/NamesAtSplitLevel.lean\#L198-L202}{{\leanfont IsSplitLevel.\allowbreak{}LiftLaw.\allowbreak{}eq\_zero\_of\_smul\_invariant}}}
\expandafter\newcommand\csname LN@inclGLlastColDatum\endcsname{\href{https://github.com/lsernau/grh-lean/blob/72fead102f663d5c7a407964a13dc0cf9d977207/grh-compact/src/New/Seesaw/NamesAtSplitLevel.lean\#L205-L207}{{\leanfont IsSplitLevel.\allowbreak{}LiftLaw.\allowbreak{}inclGL\_lastColDatum}}}
\expandafter\newcommand\csname LN@ucharDatumeqfourierOp\endcsname{\href{https://github.com/lsernau/grh-lean/blob/72fead102f663d5c7a407964a13dc0cf9d977207/grh-compact/src/New/Seesaw/NamesAtSplitLevel.lean\#L212-L225}{{\leanfont IsSplitLevel.\allowbreak{}LiftLaw.\allowbreak{}u₀\_charDatum\_eq\_fourierOp}}}
\expandafter\newcommand\csname LN@compUpTopairseqsmul\endcsname{\href{https://github.com/lsernau/grh-lean/blob/72fead102f663d5c7a407964a13dc0cf9d977207/grh-compact/src/New/Seesaw/NamesAtSplitLevel.lean\#L233-L255}{{\leanfont IsSplitLevel.\allowbreak{}compUpTo\_pairs\_eq\_smul}}}
\expandafter\newcommand\csname LN@chain\endcsname{\href{https://github.com/lsernau/grh-lean/blob/72fead102f663d5c7a407964a13dc0cf9d977207/grh-compact/src/New/Seesaw/NamesAtSplitLevel.lean\#L269-L269}{{\leanfont chain}}}
\expandafter\newcommand\csname LN@chainequivariant\endcsname{\href{https://github.com/lsernau/grh-lean/blob/72fead102f663d5c7a407964a13dc0cf9d977207/grh-compact/src/New/Seesaw/NamesAtSplitLevel.lean\#L275-L277}{{\leanfont chain\_equivariant}}}
\expandafter\newcommand\csname LN@chaina\endcsname{\href{https://github.com/lsernau/grh-lean/blob/72fead102f663d5c7a407964a13dc0cf9d977207/grh-compact/src/New/Seesaw/NamesAtSplitLevel.lean\#L281-L281}{{\leanfont chain₂}}}
\expandafter\newcommand\csname LN@twoKernelAction\endcsname{\href{https://github.com/lsernau/grh-lean/blob/72fead102f663d5c7a407964a13dc0cf9d977207/grh-compact/src/New/Seesaw/NamesAtSplitLevel.lean\#L285-L285}{{\leanfont twoKernelAction}}}
\expandafter\newcommand\csname LN@chainequivarianta\endcsname{\href{https://github.com/lsernau/grh-lean/blob/72fead102f663d5c7a407964a13dc0cf9d977207/grh-compact/src/New/Seesaw/NamesAtSplitLevel.lean\#L290-L293}{{\leanfont chain₂\_equivariant}}}
\expandafter\newcommand\csname LN@existslinearchaineqpure\endcsname{\href{https://github.com/lsernau/grh-lean/blob/72fead102f663d5c7a407964a13dc0cf9d977207/grh-compact/src/New/Seesaw/NamesAtSplitLevel.lean\#L295-L296}{{\leanfont exists\_linear\_chain\_eq\_pure}}}
\expandafter\newcommand\csname LN@IsJointlyInvariant\endcsname{\href{https://github.com/lsernau/grh-lean/blob/72fead102f663d5c7a407964a13dc0cf9d977207/grh-compact/src/New/Seesaw/NamesAtSplitLevel.lean\#L300-L302}{{\leanfont IsJointlyInvariant}}}
\expandafter\newcommand\csname LN@sumchainmemnsmul\endcsname{\href{https://github.com/lsernau/grh-lean/blob/72fead102f663d5c7a407964a13dc0cf9d977207/grh-compact/src/New/Seesaw/NamesAtSplitLevel.lean\#L307-L323}{{\leanfont IsJointlyInvariant.\allowbreak{}sum\_chain\_mem\_nsmul}}}
\expandafter\newcommand\csname LN@scalarFinatPlace\endcsname{\href{https://github.com/lsernau/grh-lean/blob/72fead102f663d5c7a407964a13dc0cf9d977207/grh-compact/src/New/Seesaw/NamesAtSplitLevel.lean\#L334-L335}{{\leanfont scalarFin\_atPlace}}}
\expandafter\newcommand\csname LN@pairWeilFinitescalarFin\endcsname{\href{https://github.com/lsernau/grh-lean/blob/72fead102f663d5c7a407964a13dc0cf9d977207/grh-compact/src/New/Seesaw/NamesAtSplitLevel.lean\#L337-L340}{{\leanfont pairWeilFinite\_scalarFin}}}
\expandafter\newcommand\csname LN@translateSum\endcsname{\href{https://github.com/lsernau/grh-lean/blob/72fead102f663d5c7a407964a13dc0cf9d977207/grh-compact/src/New/Seesaw/NamesAtSplitLevel.lean\#L345-L347}{{\leanfont translateSum}}}
\expandafter\newcommand\csname LN@projOnFinPart\endcsname{\href{https://github.com/lsernau/grh-lean/blob/72fead102f663d5c7a407964a13dc0cf9d977207/grh-compact/src/New/Seesaw/NamesAtSplitLevel.lean\#L351-L352}{{\leanfont Space.\allowbreak{}projOnFinPart}}}
\expandafter\newcommand\csname LN@detCharAt\endcsname{\href{https://github.com/lsernau/grh-lean/blob/72fead102f663d5c7a407964a13dc0cf9d977207/grh-compact/src/New/Seesaw/NamesAtSplitLevel.lean\#L358-L358}{{\leanfont IsSplitLevel.\allowbreak{}detCharAt}}}
\expandafter\newcommand\csname LN@embedGL\endcsname{\href{https://github.com/lsernau/grh-lean/blob/72fead102f663d5c7a407964a13dc0cf9d977207/grh-compact/src/New/Seesaw/NamesAtSplitLevel.lean\#L366-L367}{{\leanfont IsSplitLevel.\allowbreak{}embedGL}}}
\expandafter\newcommand\csname LN@atPlacesembedGL\endcsname{\href{https://github.com/lsernau/grh-lean/blob/72fead102f663d5c7a407964a13dc0cf9d977207/grh-compact/src/New/Seesaw/NamesAtSplitLevel.lean\#L370-L371}{{\leanfont IsSplitLevel.\allowbreak{}atPlaces\_embedGL}}}
\expandafter\newcommand\csname LN@embedGLmemsupportedOn\endcsname{\href{https://github.com/lsernau/grh-lean/blob/72fead102f663d5c7a407964a13dc0cf9d977207/grh-compact/src/New/Seesaw/NamesAtSplitLevel.lean\#L374-L375}{{\leanfont IsSplitLevel.\allowbreak{}embedGL\_mem\_supportedOn}}}
\expandafter\newcommand\csname LN@disjointstar\endcsname{\href{https://github.com/lsernau/grh-lean/blob/72fead102f663d5c7a407964a13dc0cf9d977207/grh-compact/src/New/Seesaw/NamesAtSplitLevel.lean\#L383-L383}{{\leanfont IsSplitLevel.\allowbreak{}disjoint\_star}}}
\expandafter\newcommand\csname LN@natCastmem\endcsname{\href{https://github.com/lsernau/grh-lean/blob/72fead102f663d5c7a407964a13dc0cf9d977207/grh-compact/src/New/Seesaw/NamesAtSplitLevel.lean\#L385-L385}{{\leanfont IsSplitLevel.\allowbreak{}natCast\_mem}}}
\expandafter\newcommand\csname LN@memintegralAboveofinclGLeq\endcsname{\href{https://github.com/lsernau/grh-lean/blob/72fead102f663d5c7a407964a13dc0cf9d977207/grh-compact/src/New/Seesaw/NamesAtSplitLevel.lean\#L387-L390}{{\leanfont IsSplitLevel.\allowbreak{}mem\_integralAbove\_of\_inclGL\_eq}}}
\expandafter\newcommand\csname LN@reductioneqiffidemmulmem\endcsname{\href{https://github.com/lsernau/grh-lean/blob/72fead102f663d5c7a407964a13dc0cf9d977207/grh-compact/src/New/Seesaw/NamesAtSplitLevel.lean\#L395-L396}{{\leanfont IsSplitLevel.\allowbreak{}reduction\_eq\_iff\_idem\_mul\_mem}}}
\expandafter\newcommand\csname LN@existsembedGLscalareqscalarFin\endcsname{\href{https://github.com/lsernau/grh-lean/blob/72fead102f663d5c7a407964a13dc0cf9d977207/grh-compact/src/New/Seesaw/NamesAtSplitLevel.lean\#L402-L408}{{\leanfont IsSplitLevel.\allowbreak{}exists\_embedGL\_scalar\_eq\_scalarFin}}}
\expandafter\newcommand\csname LN@name\endcsname{\href{https://github.com/lsernau/grh-lean/blob/72fead102f663d5c7a407964a13dc0cf9d977207/grh-compact/src/New/Seesaw/NamesAtSplitLevel.lean\#L414-L416}{{\leanfont IsSplitLevel.\allowbreak{}name}}}
\expandafter\newcommand\csname LN@liftFneqresidueDatum\endcsname{\href{https://github.com/lsernau/grh-lean/blob/72fead102f663d5c7a407964a13dc0cf9d977207/grh-compact/src/New/Seesaw/NamesAtSplitLevel.lean\#L419-L421}{{\leanfont IsSplitLevel.\allowbreak{}liftFn\_eq\_residueDatum}}}
\expandafter\newcommand\csname LN@isSchwartzBruhatliftFn\endcsname{\href{https://github.com/lsernau/grh-lean/blob/72fead102f663d5c7a407964a13dc0cf9d977207/grh-compact/src/New/Seesaw/NamesAtSplitLevel.lean\#L423-L425}{{\leanfont IsSplitLevel.\allowbreak{}isSchwartzBruhat\_liftFn}}}
\expandafter\newcommand\csname LN@pairWeilFinitename\endcsname{\href{https://github.com/lsernau/grh-lean/blob/72fead102f663d5c7a407964a13dc0cf9d977207/grh-compact/src/New/Seesaw/NamesAtSplitLevel.lean\#L430-L436}{{\leanfont IsSplitLevel.\allowbreak{}pairWeilFinite\_name}}}
\expandafter\newcommand\csname LN@detCharAtinclGL\endcsname{\href{https://github.com/lsernau/grh-lean/blob/72fead102f663d5c7a407964a13dc0cf9d977207/grh-compact/src/New/Seesaw/NamesAtSplitLevel.lean\#L441-L445}{{\leanfont IsSplitLevel.\allowbreak{}detCharAt\_inclGL}}}
\expandafter\newcommand\csname LN@normdetCharAtdiagNanti\endcsname{\href{https://github.com/lsernau/grh-lean/blob/72fead102f663d5c7a407964a13dc0cf9d977207/grh-compact/src/New/Seesaw/NamesAtSplitLevel.lean\#L451-L455}{{\leanfont IsSplitLevel.\allowbreak{}norm\_detCharAt\_diagN\_anti}}}
\expandafter\newcommand\csname LN@diagNeqsdiffmul\endcsname{\href{https://github.com/lsernau/grh-lean/blob/72fead102f663d5c7a407964a13dc0cf9d977207/grh-compact/src/New/Seesaw/NamesAtSplitLevel.lean\#L460-L461}{{\leanfont IsSplitLevel.\allowbreak{}diagN\_eq\_sdiff\_mul}}}
\expandafter\newcommand\csname LN@existssubresidueOverNmemiff\endcsname{\href{https://github.com/lsernau/grh-lean/blob/72fead102f663d5c7a407964a13dc0cf9d977207/grh-compact/src/New/Seesaw/NamesAtSplitLevel.lean\#L467-L472}{{\leanfont IsSplitLevel.\allowbreak{}exists\_sub\_residueOverN\_mem\_iff}}}
\expandafter\newcommand\csname LN@translateSumeqsumname\endcsname{\href{https://github.com/lsernau/grh-lean/blob/72fead102f663d5c7a407964a13dc0cf9d977207/grh-compact/src/New/Seesaw/NamesAtSplitLevel.lean\#L479-L521}{{\leanfont IsSplitLevel.\allowbreak{}translateSum\_eq\_sum\_name}}}
\newcommand{\ConcordanceRows}{\crow{restrictToYofdiag}{\href{https://github.com/lsernau/grh-lean/blob/72fead102f663d5c7a407964a13dc0cf9d977207/grh-compact/src/New/Seesaw/LiftCharacters.lean\#L20-L33}{{\leanfont Data.\allowbreak{}GaugeLe.\allowbreak{}restrictToY\_of\_diag}}}{LiftCharacters.lean}{20--33}
\crow{restrictToY}{\href{https://github.com/lsernau/grh-lean/blob/72fead102f663d5c7a407964a13dc0cf9d977207/grh-compact/src/New/Seesaw/LiftCharacters.lean\#L37-L39}{{\leanfont Data.\allowbreak{}GaugeLe.\allowbreak{}restrictToY}}}{LiftCharacters.lean}{37--39}
\crow{weylNormSqrestrictToYle}{\href{https://github.com/lsernau/grh-lean/blob/72fead102f663d5c7a407964a13dc0cf9d977207/grh-compact/src/New/Seesaw/LiftCharacters.lean\#L41-L51}{{\leanfont Data.\allowbreak{}weylNormSq\_restrictToY\_le}}}{LiftCharacters.lean}{41--51}
\crow{nonInducedPlaces}{\href{https://github.com/lsernau/grh-lean/blob/72fead102f663d5c7a407964a13dc0cf9d977207/grh-compact/src/New/Seesaw/LiftCharacters.lean\#L55-L55}{{\leanfont Data.\allowbreak{}nonInducedPlaces}}}{LiftCharacters.lean}{55--55}
\crow{isSplitLevelnonInducedPlaces}{\href{https://github.com/lsernau/grh-lean/blob/72fead102f663d5c7a407964a13dc0cf9d977207/grh-compact/src/New/Seesaw/LiftCharacters.lean\#L57-L61}{{\leanfont Data.\allowbreak{}isSplitLevel\_nonInducedPlaces}}}{LiftCharacters.lean}{57--61}
\crow{diagPointN}{\href{https://github.com/lsernau/grh-lean/blob/72fead102f663d5c7a407964a13dc0cf9d977207/grh-compact/src/New/Seesaw/LiftCharacters.lean\#L67-L67}{{\leanfont Data.\allowbreak{}diagPointN}}}{LiftCharacters.lean}{67--67}
\crow{conjEmbedGL}{\href{https://github.com/lsernau/grh-lean/blob/72fead102f663d5c7a407964a13dc0cf9d977207/grh-compact/src/New/Seesaw/LiftCharacters.lean\#L71-L71}{{\leanfont Data.\allowbreak{}conjEmbedGL}}}{LiftCharacters.lean}{71--71}
\crow{nonInducedRep}{\href{https://github.com/lsernau/grh-lean/blob/72fead102f663d5c7a407964a13dc0cf9d977207/grh-compact/src/New/Seesaw/LiftCharacters.lean\#L75-L76}{{\leanfont Data.\allowbreak{}nonInducedRep}}}{LiftCharacters.lean}{75--76}
\crow{adeleExtensionmapdiagN}{\href{https://github.com/lsernau/grh-lean/blob/72fead102f663d5c7a407964a13dc0cf9d977207/grh-compact/src/New/Seesaw/LiftCharacters.lean\#L78-L82}{{\leanfont Data.\allowbreak{}adeleExtension\_map\_diagN}}}{LiftCharacters.lean}{78--82}
\crow{lastColDatumconjnezeroiff}{\href{https://github.com/lsernau/grh-lean/blob/72fead102f663d5c7a407964a13dc0cf9d977207/grh-compact/src/New/Seesaw/LiftCharacters.lean\#L84-L95}{{\leanfont Data.\allowbreak{}lastColDatum\_conj\_ne\_zero\_iff}}}{LiftCharacters.lean}{84--95}
\crow{atPlacesconjmemmirabolic}{\href{https://github.com/lsernau/grh-lean/blob/72fead102f663d5c7a407964a13dc0cf9d977207/grh-compact/src/New/Seesaw/LiftCharacters.lean\#L97-L101}{{\leanfont Data.\allowbreak{}atPlaces\_conj\_mem\_mirabolic}}}{LiftCharacters.lean}{97--101}
\crow{memramifiedofmemramifiedColumns}{\href{https://github.com/lsernau/grh-lean/blob/72fead102f663d5c7a407964a13dc0cf9d977207/grh-compact/src/New/Seesaw/LiftCharacters.lean\#L103-L108}{{\leanfont Data.\allowbreak{}mem\_ramified\_of\_mem\_ramifiedColumns}}}{LiftCharacters.lean}{103--108}
\crow{fourierOpmemcolumn}{\href{https://github.com/lsernau/grh-lean/blob/72fead102f663d5c7a407964a13dc0cf9d977207/grh-compact/src/New/Seesaw/LiftCharacters.lean\#L112-L123}{{\leanfont Data.\allowbreak{}fourierOp\_mem\_column}}}{LiftCharacters.lean}{112--123}
\crow{existsindexDvdfourierOppow}{\href{https://github.com/lsernau/grh-lean/blob/72fead102f663d5c7a407964a13dc0cf9d977207/grh-compact/src/New/Seesaw/LiftCharacters.lean\#L128-L169}{{\leanfont Data.\allowbreak{}exists\_indexDvd\_fourierOp\_pow}}}{LiftCharacters.lean}{128--169}
\crow{charUnramifiedAtrankTwistofNmem}{\href{https://github.com/lsernau/grh-lean/blob/72fead102f663d5c7a407964a13dc0cf9d977207/grh-compact/src/New/Seesaw/LiftCharacters.lean\#L171-L174}{{\leanfont Data.\allowbreak{}charUnramifiedAt\_rankTwist\_of\_N\_mem}}}{LiftCharacters.lean}{171--174}
\crow{spaceCharrankTwist}{\href{https://github.com/lsernau/grh-lean/blob/72fead102f663d5c7a407964a13dc0cf9d977207/grh-compact/src/New/Seesaw/LiftCharacters.lean\#L178-L180}{{\leanfont Data.\allowbreak{}spaceChar\_rankTwist}}}{LiftCharacters.lean}{178--180}
\crow{IsSplitLevel}{\href{https://github.com/lsernau/grh-lean/blob/72fead102f663d5c7a407964a13dc0cf9d977207/grh-compact/src/New/Seesaw/NamesAtSplitLevel.lean\#L25-L29}{{\leanfont IsSplitLevel}}}{NamesAtSplitLevel.lean}{25--29}
\crow{residueOverN}{\href{https://github.com/lsernau/grh-lean/blob/72fead102f663d5c7a407964a13dc0cf9d977207/grh-compact/src/New/Seesaw/NamesAtSplitLevel.lean\#L33-L34}{{\leanfont residueOverN}}}{NamesAtSplitLevel.lean}{33--34}
\crow{residueOverNapply}{\href{https://github.com/lsernau/grh-lean/blob/72fead102f663d5c7a407964a13dc0cf9d977207/grh-compact/src/New/Seesaw/NamesAtSplitLevel.lean\#L37-L38}{{\leanfont residueOverN\_apply}}}{NamesAtSplitLevel.lean}{37--38}
\crow{idemmulresidueOverN}{\href{https://github.com/lsernau/grh-lean/blob/72fead102f663d5c7a407964a13dc0cf9d977207/grh-compact/src/New/Seesaw/NamesAtSplitLevel.lean\#L41-L41}{{\leanfont idem\_mul\_residueOverN}}}{NamesAtSplitLevel.lean}{41--41}
\crow{diagN}{\href{https://github.com/lsernau/grh-lean/blob/72fead102f663d5c7a407964a13dc0cf9d977207/grh-compact/src/New/Seesaw/NamesAtSplitLevel.lean\#L48-L49}{{\leanfont IsSplitLevel.\allowbreak{}diagN}}}{NamesAtSplitLevel.lean}{48--49}
\crow{residueEquiv}{\href{https://github.com/lsernau/grh-lean/blob/72fead102f663d5c7a407964a13dc0cf9d977207/grh-compact/src/New/Seesaw/NamesAtSplitLevel.lean\#L56-L57}{{\leanfont IsSplitLevel.\allowbreak{}residueEquiv}}}{NamesAtSplitLevel.lean}{56--57}
\crow{reduction}{\href{https://github.com/lsernau/grh-lean/blob/72fead102f663d5c7a407964a13dc0cf9d977207/grh-compact/src/New/Seesaw/NamesAtSplitLevel.lean\#L61-L61}{{\leanfont IsSplitLevel.\allowbreak{}reduction}}}{NamesAtSplitLevel.lean}{61--61}
\crow{reductioneqiff}{\href{https://github.com/lsernau/grh-lean/blob/72fead102f663d5c7a407964a13dc0cf9d977207/grh-compact/src/New/Seesaw/NamesAtSplitLevel.lean\#L64-L67}{{\leanfont IsSplitLevel.\allowbreak{}reduction\_eq\_iff}}}{NamesAtSplitLevel.lean}{64--67}
\crow{inclGL}{\href{https://github.com/lsernau/grh-lean/blob/72fead102f663d5c7a407964a13dc0cf9d977207/grh-compact/src/New/Seesaw/NamesAtSplitLevel.lean\#L72-L73}{{\leanfont IsSplitLevel.\allowbreak{}inclGL}}}{NamesAtSplitLevel.lean}{72--73}
\crow{reductionGL}{\href{https://github.com/lsernau/grh-lean/blob/72fead102f663d5c7a407964a13dc0cf9d977207/grh-compact/src/New/Seesaw/NamesAtSplitLevel.lean\#L77-L77}{{\leanfont IsSplitLevel.\allowbreak{}reductionGL}}}{NamesAtSplitLevel.lean}{77--77}
\crow{liftGL}{\href{https://github.com/lsernau/grh-lean/blob/72fead102f663d5c7a407964a13dc0cf9d977207/grh-compact/src/New/Seesaw/NamesAtSplitLevel.lean\#L80-L81}{{\leanfont IsSplitLevel.\allowbreak{}liftGL}}}{NamesAtSplitLevel.lean}{80--81}
\crow{diagNInv}{\href{https://github.com/lsernau/grh-lean/blob/72fead102f663d5c7a407964a13dc0cf9d977207/grh-compact/src/New/Seesaw/NamesAtSplitLevel.lean\#L87-L87}{{\leanfont IsSplitLevel.\allowbreak{}diagNInv}}}{NamesAtSplitLevel.lean}{87--87}
\crow{psiN}{\href{https://github.com/lsernau/grh-lean/blob/72fead102f663d5c7a407964a13dc0cf9d977207/grh-compact/src/New/Seesaw/NamesAtSplitLevel.lean\#L91-L91}{{\leanfont IsSplitLevel.\allowbreak{}psiN}}}{NamesAtSplitLevel.lean}{91--91}
\crow{lastBasisVector}{\href{https://github.com/lsernau/grh-lean/blob/72fead102f663d5c7a407964a13dc0cf9d977207/grh-compact/src/New/Seesaw/NamesAtSplitLevel.lean\#L94-L94}{{\leanfont IsSplitLevel.\allowbreak{}lastBasisVector}}}{NamesAtSplitLevel.lean}{94--94}
\crow{existsGLlastcolandrow}{\href{https://github.com/lsernau/grh-lean/blob/72fead102f663d5c7a407964a13dc0cf9d977207/grh-compact/src/New/Seesaw/NamesAtSplitLevel.lean\#L98-L102}{{\leanfont IsSplitLevel.\allowbreak{}exists\_GL\_last\_col\_and\_row}}}{NamesAtSplitLevel.lean}{98--102}
\crow{fourierOp}{\href{https://github.com/lsernau/grh-lean/blob/72fead102f663d5c7a407964a13dc0cf9d977207/grh-compact/src/New/Seesaw/NamesAtSplitLevel.lean\#L106-L107}{{\leanfont IsSplitLevel.\allowbreak{}fourierOp}}}{NamesAtSplitLevel.lean}{106--107}
\crow{mirabolic}{\href{https://github.com/lsernau/grh-lean/blob/72fead102f663d5c7a407964a13dc0cf9d977207/grh-compact/src/New/Seesaw/NamesAtSplitLevel.lean\#L111-L112}{{\leanfont IsSplitLevel.\allowbreak{}mirabolic}}}{NamesAtSplitLevel.lean}{111--112}
\crow{ramified}{\href{https://github.com/lsernau/grh-lean/blob/72fead102f663d5c7a407964a13dc0cf9d977207/grh-compact/src/New/Seesaw/NamesAtSplitLevel.lean\#L116-L117}{{\leanfont IsSplitLevel.\allowbreak{}ramified}}}{NamesAtSplitLevel.lean}{116--117}
\crow{liftFn}{\href{https://github.com/lsernau/grh-lean/blob/72fead102f663d5c7a407964a13dc0cf9d977207/grh-compact/src/New/Seesaw/NamesAtSplitLevel.lean\#L121-L123}{{\leanfont IsSplitLevel.\allowbreak{}liftFn}}}{NamesAtSplitLevel.lean}{121--123}
\crow{lastRowDatum}{\href{https://github.com/lsernau/grh-lean/blob/72fead102f663d5c7a407964a13dc0cf9d977207/grh-compact/src/New/Seesaw/NamesAtSplitLevel.lean\#L127-L128}{{\leanfont IsSplitLevel.\allowbreak{}lastRowDatum}}}{NamesAtSplitLevel.lean}{127--128}
\crow{lastColDatum}{\href{https://github.com/lsernau/grh-lean/blob/72fead102f663d5c7a407964a13dc0cf9d977207/grh-compact/src/New/Seesaw/NamesAtSplitLevel.lean\#L132-L132}{{\leanfont IsSplitLevel.\allowbreak{}lastColDatum}}}{NamesAtSplitLevel.lean}{132--132}
\crow{charDatum}{\href{https://github.com/lsernau/grh-lean/blob/72fead102f663d5c7a407964a13dc0cf9d977207/grh-compact/src/New/Seesaw/NamesAtSplitLevel.lean\#L136-L136}{{\leanfont IsSplitLevel.\allowbreak{}charDatum}}}{NamesAtSplitLevel.lean}{136--136}
\crow{LiftLaw}{\href{https://github.com/lsernau/grh-lean/blob/72fead102f663d5c7a407964a13dc0cf9d977207/grh-compact/src/New/Seesaw/NamesAtSplitLevel.lean\#L141-L146}{{\leanfont IsSplitLevel.\allowbreak{}LiftLaw}}}{NamesAtSplitLevel.lean}{141--146}
\crow{lastRowDatumconju}{\href{https://github.com/lsernau/grh-lean/blob/72fead102f663d5c7a407964a13dc0cf9d977207/grh-compact/src/New/Seesaw/NamesAtSplitLevel.lean\#L161-L165}{{\leanfont IsSplitLevel.\allowbreak{}lastRowDatum\_conj\_u₀}}}{NamesAtSplitLevel.lean}{161--165}
\crow{existsinclGLeq}{\href{https://github.com/lsernau/grh-lean/blob/72fead102f663d5c7a407964a13dc0cf9d977207/grh-compact/src/New/Seesaw/NamesAtSplitLevel.lean\#L167-L168}{{\leanfont IsSplitLevel.\allowbreak{}exists\_inclGL\_eq}}}{NamesAtSplitLevel.lean}{167--168}
\crow{lastColDatumnezero}{\href{https://github.com/lsernau/grh-lean/blob/72fead102f663d5c7a407964a13dc0cf9d977207/grh-compact/src/New/Seesaw/NamesAtSplitLevel.lean\#L170-L171}{{\leanfont IsSplitLevel.\allowbreak{}lastColDatum\_ne\_zero}}}{NamesAtSplitLevel.lean}{170--171}
\crow{mirabolicfixfourierOp}{\href{https://github.com/lsernau/grh-lean/blob/72fead102f663d5c7a407964a13dc0cf9d977207/grh-compact/src/New/Seesaw/NamesAtSplitLevel.lean\#L178-L188}{{\leanfont IsSplitLevel.\allowbreak{}mirabolic\_fix\_fourierOp}}}{NamesAtSplitLevel.lean}{178--188}
\crow{liftFncompeqsumsmul}{\href{https://github.com/lsernau/grh-lean/blob/72fead102f663d5c7a407964a13dc0cf9d977207/grh-compact/src/New/Seesaw/NamesAtSplitLevel.lean\#L191-L193}{{\leanfont IsSplitLevel.\allowbreak{}liftFn\_comp\_eq\_sum\_smul}}}{NamesAtSplitLevel.lean}{191--193}
\crow{eqzeroofsmulinvariant}{\href{https://github.com/lsernau/grh-lean/blob/72fead102f663d5c7a407964a13dc0cf9d977207/grh-compact/src/New/Seesaw/NamesAtSplitLevel.lean\#L198-L202}{{\leanfont IsSplitLevel.\allowbreak{}LiftLaw.\allowbreak{}eq\_zero\_of\_smul\_invariant}}}{NamesAtSplitLevel.lean}{198--202}
\crow{inclGLlastColDatum}{\href{https://github.com/lsernau/grh-lean/blob/72fead102f663d5c7a407964a13dc0cf9d977207/grh-compact/src/New/Seesaw/NamesAtSplitLevel.lean\#L205-L207}{{\leanfont IsSplitLevel.\allowbreak{}LiftLaw.\allowbreak{}inclGL\_lastColDatum}}}{NamesAtSplitLevel.lean}{205--207}
\crow{ucharDatumeqfourierOp}{\href{https://github.com/lsernau/grh-lean/blob/72fead102f663d5c7a407964a13dc0cf9d977207/grh-compact/src/New/Seesaw/NamesAtSplitLevel.lean\#L212-L225}{{\leanfont IsSplitLevel.\allowbreak{}LiftLaw.\allowbreak{}u₀\_charDatum\_eq\_fourierOp}}}{NamesAtSplitLevel.lean}{212--225}
\crow{compUpTopairseqsmul}{\href{https://github.com/lsernau/grh-lean/blob/72fead102f663d5c7a407964a13dc0cf9d977207/grh-compact/src/New/Seesaw/NamesAtSplitLevel.lean\#L233-L255}{{\leanfont IsSplitLevel.\allowbreak{}compUpTo\_pairs\_eq\_smul}}}{NamesAtSplitLevel.lean}{233--255}
\crow{chain}{\href{https://github.com/lsernau/grh-lean/blob/72fead102f663d5c7a407964a13dc0cf9d977207/grh-compact/src/New/Seesaw/NamesAtSplitLevel.lean\#L269-L269}{{\leanfont chain}}}{NamesAtSplitLevel.lean}{269--269}
\crow{chainequivariant}{\href{https://github.com/lsernau/grh-lean/blob/72fead102f663d5c7a407964a13dc0cf9d977207/grh-compact/src/New/Seesaw/NamesAtSplitLevel.lean\#L275-L277}{{\leanfont chain\_equivariant}}}{NamesAtSplitLevel.lean}{275--277}
\crow{chaina}{\href{https://github.com/lsernau/grh-lean/blob/72fead102f663d5c7a407964a13dc0cf9d977207/grh-compact/src/New/Seesaw/NamesAtSplitLevel.lean\#L281-L281}{{\leanfont chain₂}}}{NamesAtSplitLevel.lean}{281--281}
\crow{twoKernelAction}{\href{https://github.com/lsernau/grh-lean/blob/72fead102f663d5c7a407964a13dc0cf9d977207/grh-compact/src/New/Seesaw/NamesAtSplitLevel.lean\#L285-L285}{{\leanfont twoKernelAction}}}{NamesAtSplitLevel.lean}{285--285}
\crow{chainequivarianta}{\href{https://github.com/lsernau/grh-lean/blob/72fead102f663d5c7a407964a13dc0cf9d977207/grh-compact/src/New/Seesaw/NamesAtSplitLevel.lean\#L290-L293}{{\leanfont chain₂\_equivariant}}}{NamesAtSplitLevel.lean}{290--293}
\crow{existslinearchaineqpure}{\href{https://github.com/lsernau/grh-lean/blob/72fead102f663d5c7a407964a13dc0cf9d977207/grh-compact/src/New/Seesaw/NamesAtSplitLevel.lean\#L295-L296}{{\leanfont exists\_linear\_chain\_eq\_pure}}}{NamesAtSplitLevel.lean}{295--296}
\crow{IsJointlyInvariant}{\href{https://github.com/lsernau/grh-lean/blob/72fead102f663d5c7a407964a13dc0cf9d977207/grh-compact/src/New/Seesaw/NamesAtSplitLevel.lean\#L300-L302}{{\leanfont IsJointlyInvariant}}}{NamesAtSplitLevel.lean}{300--302}
\crow{sumchainmemnsmul}{\href{https://github.com/lsernau/grh-lean/blob/72fead102f663d5c7a407964a13dc0cf9d977207/grh-compact/src/New/Seesaw/NamesAtSplitLevel.lean\#L307-L323}{{\leanfont IsJointlyInvariant.\allowbreak{}sum\_chain\_mem\_nsmul}}}{NamesAtSplitLevel.lean}{307--323}
\crow{scalarFinatPlace}{\href{https://github.com/lsernau/grh-lean/blob/72fead102f663d5c7a407964a13dc0cf9d977207/grh-compact/src/New/Seesaw/NamesAtSplitLevel.lean\#L334-L335}{{\leanfont scalarFin\_atPlace}}}{NamesAtSplitLevel.lean}{334--335}
\crow{pairWeilFinitescalarFin}{\href{https://github.com/lsernau/grh-lean/blob/72fead102f663d5c7a407964a13dc0cf9d977207/grh-compact/src/New/Seesaw/NamesAtSplitLevel.lean\#L337-L340}{{\leanfont pairWeilFinite\_scalarFin}}}{NamesAtSplitLevel.lean}{337--340}
\crow{translateSum}{\href{https://github.com/lsernau/grh-lean/blob/72fead102f663d5c7a407964a13dc0cf9d977207/grh-compact/src/New/Seesaw/NamesAtSplitLevel.lean\#L345-L347}{{\leanfont translateSum}}}{NamesAtSplitLevel.lean}{345--347}
\crow{projOnFinPart}{\href{https://github.com/lsernau/grh-lean/blob/72fead102f663d5c7a407964a13dc0cf9d977207/grh-compact/src/New/Seesaw/NamesAtSplitLevel.lean\#L351-L352}{{\leanfont Space.\allowbreak{}projOnFinPart}}}{NamesAtSplitLevel.lean}{351--352}
\crow{detCharAt}{\href{https://github.com/lsernau/grh-lean/blob/72fead102f663d5c7a407964a13dc0cf9d977207/grh-compact/src/New/Seesaw/NamesAtSplitLevel.lean\#L358-L358}{{\leanfont IsSplitLevel.\allowbreak{}detCharAt}}}{NamesAtSplitLevel.lean}{358--358}
\crow{embedGL}{\href{https://github.com/lsernau/grh-lean/blob/72fead102f663d5c7a407964a13dc0cf9d977207/grh-compact/src/New/Seesaw/NamesAtSplitLevel.lean\#L366-L367}{{\leanfont IsSplitLevel.\allowbreak{}embedGL}}}{NamesAtSplitLevel.lean}{366--367}
\crow{atPlacesembedGL}{\href{https://github.com/lsernau/grh-lean/blob/72fead102f663d5c7a407964a13dc0cf9d977207/grh-compact/src/New/Seesaw/NamesAtSplitLevel.lean\#L370-L371}{{\leanfont IsSplitLevel.\allowbreak{}atPlaces\_embedGL}}}{NamesAtSplitLevel.lean}{370--371}
\crow{embedGLmemsupportedOn}{\href{https://github.com/lsernau/grh-lean/blob/72fead102f663d5c7a407964a13dc0cf9d977207/grh-compact/src/New/Seesaw/NamesAtSplitLevel.lean\#L374-L375}{{\leanfont IsSplitLevel.\allowbreak{}embedGL\_mem\_supportedOn}}}{NamesAtSplitLevel.lean}{374--375}
\crow{disjointstar}{\href{https://github.com/lsernau/grh-lean/blob/72fead102f663d5c7a407964a13dc0cf9d977207/grh-compact/src/New/Seesaw/NamesAtSplitLevel.lean\#L383-L383}{{\leanfont IsSplitLevel.\allowbreak{}disjoint\_star}}}{NamesAtSplitLevel.lean}{383--383}
\crow{natCastmem}{\href{https://github.com/lsernau/grh-lean/blob/72fead102f663d5c7a407964a13dc0cf9d977207/grh-compact/src/New/Seesaw/NamesAtSplitLevel.lean\#L385-L385}{{\leanfont IsSplitLevel.\allowbreak{}natCast\_mem}}}{NamesAtSplitLevel.lean}{385--385}
\crow{memintegralAboveofinclGLeq}{\href{https://github.com/lsernau/grh-lean/blob/72fead102f663d5c7a407964a13dc0cf9d977207/grh-compact/src/New/Seesaw/NamesAtSplitLevel.lean\#L387-L390}{{\leanfont IsSplitLevel.\allowbreak{}mem\_integralAbove\_of\_inclGL\_eq}}}{NamesAtSplitLevel.lean}{387--390}
\crow{reductioneqiffidemmulmem}{\href{https://github.com/lsernau/grh-lean/blob/72fead102f663d5c7a407964a13dc0cf9d977207/grh-compact/src/New/Seesaw/NamesAtSplitLevel.lean\#L395-L396}{{\leanfont IsSplitLevel.\allowbreak{}reduction\_eq\_iff\_idem\_mul\_mem}}}{NamesAtSplitLevel.lean}{395--396}
\crow{existsembedGLscalareqscalarFin}{\href{https://github.com/lsernau/grh-lean/blob/72fead102f663d5c7a407964a13dc0cf9d977207/grh-compact/src/New/Seesaw/NamesAtSplitLevel.lean\#L402-L408}{{\leanfont IsSplitLevel.\allowbreak{}exists\_embedGL\_scalar\_eq\_scalarFin}}}{NamesAtSplitLevel.lean}{402--408}
\crow{name}{\href{https://github.com/lsernau/grh-lean/blob/72fead102f663d5c7a407964a13dc0cf9d977207/grh-compact/src/New/Seesaw/NamesAtSplitLevel.lean\#L414-L416}{{\leanfont IsSplitLevel.\allowbreak{}name}}}{NamesAtSplitLevel.lean}{414--416}
\crow{liftFneqresidueDatum}{\href{https://github.com/lsernau/grh-lean/blob/72fead102f663d5c7a407964a13dc0cf9d977207/grh-compact/src/New/Seesaw/NamesAtSplitLevel.lean\#L419-L421}{{\leanfont IsSplitLevel.\allowbreak{}liftFn\_eq\_residueDatum}}}{NamesAtSplitLevel.lean}{419--421}
\crow{isSchwartzBruhatliftFn}{\href{https://github.com/lsernau/grh-lean/blob/72fead102f663d5c7a407964a13dc0cf9d977207/grh-compact/src/New/Seesaw/NamesAtSplitLevel.lean\#L423-L425}{{\leanfont IsSplitLevel.\allowbreak{}isSchwartzBruhat\_liftFn}}}{NamesAtSplitLevel.lean}{423--425}
\crow{pairWeilFinitename}{\href{https://github.com/lsernau/grh-lean/blob/72fead102f663d5c7a407964a13dc0cf9d977207/grh-compact/src/New/Seesaw/NamesAtSplitLevel.lean\#L430-L436}{{\leanfont IsSplitLevel.\allowbreak{}pairWeilFinite\_name}}}{NamesAtSplitLevel.lean}{430--436}
\crow{detCharAtinclGL}{\href{https://github.com/lsernau/grh-lean/blob/72fead102f663d5c7a407964a13dc0cf9d977207/grh-compact/src/New/Seesaw/NamesAtSplitLevel.lean\#L441-L445}{{\leanfont IsSplitLevel.\allowbreak{}detCharAt\_inclGL}}}{NamesAtSplitLevel.lean}{441--445}
\crow{normdetCharAtdiagNanti}{\href{https://github.com/lsernau/grh-lean/blob/72fead102f663d5c7a407964a13dc0cf9d977207/grh-compact/src/New/Seesaw/NamesAtSplitLevel.lean\#L451-L455}{{\leanfont IsSplitLevel.\allowbreak{}norm\_detCharAt\_diagN\_anti}}}{NamesAtSplitLevel.lean}{451--455}
\crow{diagNeqsdiffmul}{\href{https://github.com/lsernau/grh-lean/blob/72fead102f663d5c7a407964a13dc0cf9d977207/grh-compact/src/New/Seesaw/NamesAtSplitLevel.lean\#L460-L461}{{\leanfont IsSplitLevel.\allowbreak{}diagN\_eq\_sdiff\_mul}}}{NamesAtSplitLevel.lean}{460--461}
\crow{existssubresidueOverNmemiff}{\href{https://github.com/lsernau/grh-lean/blob/72fead102f663d5c7a407964a13dc0cf9d977207/grh-compact/src/New/Seesaw/NamesAtSplitLevel.lean\#L467-L472}{{\leanfont IsSplitLevel.\allowbreak{}exists\_sub\_residueOverN\_mem\_iff}}}{NamesAtSplitLevel.lean}{467--472}
\crow{translateSumeqsumname}{\href{https://github.com/lsernau/grh-lean/blob/72fead102f663d5c7a407964a13dc0cf9d977207/grh-compact/src/New/Seesaw/NamesAtSplitLevel.lean\#L479-L521}{{\leanfont IsSplitLevel.\allowbreak{}translateSum\_eq\_sum\_name}}}{NamesAtSplitLevel.lean}{479--521}}
% the notations of the defined functions: one slot for each explicit argument of the Lean declaration
\newcommand{\nonind}[1]{\Sigma_{\jmath}(#1)}
\newcommand{\dpt}[2]{\Delta_{#2}(#1)}
\newcommand{\cemb}[2]{\operatorname{emb}^{\Delta}_{#1,#2}}
\newcommand{\nrep}[2]{\rho_{#1,#2}}
\newcommand{\splitlevel}[2]{#1\text{ is a split level for }#2}
\newcommand{\rN}[1]{[#1/N]}
\newcommand{\diagN}[2]{\operatorname{dg}_{#1}(#2)}
\newcommand{\psiN}[1]{\psi_N(#1)}
\newcommand{\ram}[1]{\operatorname{Ram}(#1)}
\newcommand{\bcol}[1]{\mathbf{b}_{#1}}
\newcommand{\cchar}[1]{\mathbf{c}_{#1}}
\newcommand{\law}[4]{#3\text{ is a lift law}\allowbreak\ \text{from }#1\allowbreak\ \text{to }#2\allowbreak\ \text{with constant }#4}
\newcommand{\chain}[2]{\operatorname{ch}(#1;#2)}
\newcommand{\chaintwo}[2]{\operatorname{ch}_2(#1;#2)}
\newcommand{\twoact}[3]{\omega^{(2)}_{#1,#3}(#2)}
\newcommand{\transum}[5]{\mathsf{T}_{#4,#3}(#1,#2,#5)}
\newcommand{\projfin}[2]{\operatorname{pr}^{#2}_{#1}}
\newcommand{\detchar}[2]{\operatorname{dc}_{#1,#2}}
\newcommand{\UnitMembers}{312}
\newcommand{\UnitInNote}{78}
\newcommand{\UnitProofs}{83}
\newcommand{\UnitSteps}{503}
\newcommand{\UnitAllSteps}{550}


\linespread{1.08}                       % 11 pt on 14.7 pt
\setlength{\parindent}{0pt}
\setlength{\parskip}{5.5pt plus 1.5pt minus 0.5pt}
\setlength{\abovedisplayskip}{7.5pt plus 2pt minus 3pt}
\setlength{\belowdisplayskip}{7.5pt plus 2pt minus 3pt}
\setlength{\abovedisplayshortskip}{3pt plus 2pt}
\setlength{\belowdisplayshortskip}{6pt plus 2pt minus 2pt}
\emergencystretch=1.5em
\clubpenalty=10000 \widowpenalty=10000 \displaywidowpenalty=10000
\relpenalty=7000 \binoppenalty=9000
\raggedbottom
\pagestyle{fancy}\fancyhf{}\renewcommand{\headrulewidth}{0pt}
\fancyfoot[C]{\sffamily\fontsize{7.6}{9}\selectfont\color{inkfaint}\thepage\,/\,\pageref*{LastPage}}
\setlist{topsep=2pt,partopsep=0pt,parsep=0pt,itemsep=4pt,leftmargin=2.1em,labelsep=0.55em}

\newcommand{\Q}{\mathbb{Q}}
\newcommand{\Z}{\mathbb{Z}}
\newcommand{\F}{\mathbb{F}}
\newcommand{\len}{\operatorname{len}}
\newcommand{\cok}{\operatorname{cok}}
\newcommand{\End}{\operatorname{End}}
\newcommand{\Hom}{\operatorname{Hom}}
\newcommand{\Gal}{\operatorname{Gal}}
\newcommand{\diag}{\operatorname{diag}}
\newcommand{\kw}[1]{\emph{#1}}
\newcommand{\m}[1]{\mbox{$#1$}}        % a short formula that must not be broken across lines
\newcommand{\head}[1]{\par\addvspace{17pt plus 4pt minus 2pt}\Needspace{6\baselineskip}%
  \noindent{\fontsize{13.2}{16}\selectfont\bfseries\color{ink}#1}\par\nopagebreak\vspace{-1pt}}
\newcommand{\sect}[2]{\refstepcounter{section}\label{#2}\par\addvspace{19pt plus 4pt minus 2pt}\Needspace{7\baselineskip}%
  \noindent\llap{{\fontsize{13.2}{16}\selectfont\bfseries\color{inkfaint}\thesection}\hspace{9pt}}%
  {\fontsize{13.2}{16}\selectfont\bfseries\color{ink}#1}\par\nopagebreak\vspace{-1pt}}
\newcommand{\subhead}[1]{\par\addvspace{9pt plus 2pt}\Needspace{4\baselineskip}\noindent\textbf{#1}\enspace\ignorespaces}
% where a statement lives in the Lean archive
\newcommand{\leantag}{{\sffamily\bfseries\fontsize{6.3}{8}\selectfont\color{leancol}\addfontfeatures{LetterSpace=7}LEAN}}
\newcommand{\sep}{\unskip\nobreak\hspace{5pt}{\color{inkfaint}\textperiodcentered}\hspace{5pt plus 2pt}\ignorespaces}
\newcommand{\leanline}[1]{\par\nopagebreak\vspace{-2.5pt}{\fontsize{8.9}{12.4}\selectfont\color{inksoft}\raggedright
  \noindent\llap{\leantag\hspace{9pt}}#1\par}\vspace{1pt}}
% a displayed statement in a tinted box: #1 colour, #2 text
\newcommand{\statement}[2]{\par\addvspace{7pt}\begingroup\setlength{\fboxsep}{8.5pt}\setlength{\fboxrule}{0.6pt}%
  \noindent\fcolorbox{#1!30!white}{#1!3!white}{\parbox{\dimexpr\linewidth-2\fboxsep-2\fboxrule\relax}{#2}}%
  \endgroup\par\addvspace{7pt}}
\newcommand{\nt}[2]{\par\noindent\hangindent=1.1em\hangafter=1 #1\enspace{\color{inksoft}#2}\par\vspace{1pt}}
\newcommand{\legdot}[1]{\tikz[baseline=-0.62ex]\fill[#1] (0,0) circle (2.5pt);\hspace{3pt}}
\newcommand{\crow}[4]{\csname CS@#1\endcsname & #2 & \textcolor{inksoft}{#3} & \textcolor{inksoft}{#4}\\}
\newcommand{\leanquoteline}[2]{\par\noindent\hangindent=\dimexpr#1\fontdimen2\font+2.2em\relax\hangafter=1
  \hspace*{#1\fontdimen2\font}#2}
\newcommand{\quoted}[1]{\par\vspace{3pt}{\ttfamily\fontsize{7.5}{10.4}\selectfont\color{leancol}\raggedright
  \setlength{\parskip}{0pt}\leanquote{#1}\par}\vspace{4pt}}
\newcommand{\bx}[3]{{\sffamily\bfseries\fontsize{7.3}{9.2}\selectfont\color{#1!82!black}#2}\\[1.5pt]#3}
\newcommand{\cs}[2]{\expandafter\newcommand\csname CS@#1\endcsname{#2}}

% theorems
\newtheoremstyle{plainx}{7pt plus 2pt}{7pt plus 2pt}{\itshape}{}{\bfseries}{.}{0.5em}{}
\newtheoremstyle{uprightx}{7pt plus 2pt}{7pt plus 2pt}{}{}{\bfseries}{.}{0.5em}{}
\theoremstyle{plainx}
\newtheorem{theorem}{Theorem}[section]
\newtheorem{lemma}[theorem]{Lemma}
\newtheorem{proposition}[theorem]{Proposition}
\newtheorem{corollary}[theorem]{Corollary}
\theoremstyle{uprightx}
\newtheorem{definition}[theorem]{Definition}
\newtheorem{setting}[theorem]{Setting}          % hypotheses shared by several statements; each of them cites it by number
\newtheorem{remark}[theorem]{Remark}
\newtheorem{example}[theorem]{Example}
\renewcommand{\qedsymbol}{{\setlength{\fboxsep}{0pt}\setlength{\fboxrule}{0.5pt}\fbox{\rule{0pt}{0.46em}\hspace{0.46em}}}}
% what holds throughout the note: the ambient objects only (implicit and instance arguments), never a hypothesis
\newenvironment{conventions}{\par\addvspace{5pt}\noindent\textbf{Conventions.}\enspace\ignorespaces}{\par\addvspace{3pt}}
\makeatletter
\renewenvironment{proof}[1][Proof]{\par\pushQED{\qed}\normalfont\topsep0pt\partopsep0pt
  \trivlist\item[\hskip\labelsep\itshape #1.]\ignorespaces}{\popQED\endtrivlist\@endpefalse}
\makeatother

% where each declaration of the concordance is treated: \cs{key}{place}, the place naming a \label or a lettered item
\cs{IsSplitLevel}{Definition~\ref{d:split}}
\cs{residueOverN}{Definition~\ref{d:split}}
\cs{residueOverNapply}{Definition~\ref{d:split}}
\cs{residueEquiv}{Definition~\ref{d:red}}
\cs{reduction}{Definition~\ref{d:red}}
\cs{inclGL}{Definition~\ref{d:red}}
\cs{reductionGL}{Definition~\ref{d:red}}
\cs{liftGL}{Definition~\ref{d:red}}
\cs{lastBasisVector}{Definition~\ref{d:red}}
\cs{psiN}{Definition~\ref{d:split}}
\cs{diagN}{Definition~\ref{d:diag}}
\cs{diagNInv}{Definition~\ref{d:diag}}
\cs{mirabolic}{Definition~\ref{d:diag}}
\cs{ramified}{Definition~\ref{d:diag}}
\cs{liftFn}{Definition~\ref{d:fn}}
\cs{lastRowDatum}{Definition~\ref{d:fn}}
\cs{lastColDatum}{Definition~\ref{d:fn}}
\cs{charDatum}{Definition~\ref{d:fn}}
\cs{existsGLlastcolandrow}{Lemma~\ref{l:kofy}}
\cs{fourierOp}{Definition~\ref{d:F}}
\cs{LiftLaw}{Definition~\ref{d:law}}
\cs{chain}{Definition~\ref{d:chain}}
\cs{chaina}{Definition~\ref{d:chain}}
\cs{twoKernelAction}{Definition~\ref{d:chain}}
\cs{IsJointlyInvariant}{Definition~\ref{d:joint}}
\cs{embedGL}{Definition~\ref{d:emb}}
\cs{detCharAt}{Definition~\ref{d:emb}}
\cs{name}{Definition~\ref{d:names}}
\cs{projOnFinPart}{Definition~\ref{d:names}}
\cs{translateSum}{Definition~\ref{d:names}}
\cs{nonInducedPlaces}{Definition~\ref{d:data}}
\cs{diagPointN}{Definition~\ref{d:data}}
\cs{conjEmbedGL}{Definition~\ref{d:data}}
\cs{nonInducedRep}{Definition~\ref{d:data}}
\cs{idemmulresidueOverN}{Lemma~\ref{l:elem}}
\cs{reductioneqiff}{Lemma~\ref{l:elem}}
\cs{disjointstar}{Lemma~\ref{l:elem}}
\cs{natCastmem}{Lemma~\ref{l:elem}}
\cs{reductioneqiffidemmulmem}{Lemma~\ref{l:elem}}
\cs{diagNeqsdiffmul}{Lemma~\ref{l:elem}}
\cs{existssubresidueOverNmemiff}{Lemma~\ref{l:resid}}
\cs{lastRowDatumconju}{Lemma~\ref{l:fnfacts}}
\cs{existsinclGLeq}{Lemma~\ref{l:fnfacts}}
\cs{lastColDatumnezero}{Lemma~\ref{l:fnfacts}}
\cs{liftFncompeqsumsmul}{Lemma~\ref{l:fnfacts}}
\cs{mirabolicfixfourierOp}{Lemma~\ref{l:mirfix}}
\cs{eqzeroofsmulinvariant}{Lemma~\ref{l:lawfacts}}
\cs{inclGLlastColDatum}{Lemma~\ref{l:lawfacts}}
\cs{ucharDatumeqfourierOp}{Proposition~\ref{p:slide}}
\cs{compUpTopairseqsmul}{Theorem A}
\cs{chainequivariant}{Lemma~\ref{l:cheq}}
\cs{chainequivarianta}{Lemma~\ref{l:cheq}}
\cs{existslinearchaineqpure}{Lemma~\ref{l:cheq}}
\cs{sumchainmemnsmul}{Theorem C}
\cs{atPlacesembedGL}{Lemma~\ref{l:embfacts}}
\cs{embedGLmemsupportedOn}{Lemma~\ref{l:embfacts}}
\cs{memintegralAboveofinclGLeq}{Lemma~\ref{l:embfacts}}
\cs{existsembedGLscalareqscalarFin}{Lemma~\ref{l:embfacts}}
\cs{scalarFinatPlace}{Lemma~\ref{l:embfacts}}
\cs{pairWeilFinitescalarFin}{Lemma~\ref{l:names}}
\cs{pairWeilFinitename}{Lemma~\ref{l:names}}
\cs{liftFneqresidueDatum}{Lemma~\ref{l:names}}
\cs{isSchwartzBruhatliftFn}{Lemma~\ref{l:names}}
\cs{detCharAtinclGL}{Lemma~\ref{l:detchar}}
\cs{normdetCharAtdiagNanti}{Lemma~\ref{l:detchar}}
\cs{translateSumeqsumname}{Theorem B}
\cs{isSplitLevelnonInducedPlaces}{Lemma~\ref{l:nonind}}
\cs{adeleExtensionmapdiagN}{Lemma~\ref{l:nonind}}
\cs{lastColDatumconjnezeroiff}{Lemma~\ref{l:colmir}}
\cs{atPlacesconjmemmirabolic}{Lemma~\ref{l:colmir}}
\cs{memramifiedofmemramifiedColumns}{Lemma~\ref{l:ramcol}}
\cs{fourierOpmemcolumn}{Lemma~\ref{l:ramcol}}
\cs{restrictToYofdiag}{Lemma~\ref{l:restr}}
\cs{restrictToY}{Lemma~\ref{l:restr}}
\cs{weylNormSqrestrictToYle}{Lemma~\ref{l:restr}}
\cs{existsindexDvdfourierOppow}{Theorem D}
\cs{charUnramifiedAtrankTwistofNmem}{Lemma~\ref{l:chars}}
\cs{spaceCharrankTwist}{Lemma~\ref{l:chars}}

% ------------------------------------------------------------------ macros of this note that are not notations of
% defined functions (those are made by the build, in data.tex, from the spec)
\newcommand{\CC}{\mathbb{C}}
\newcommand{\ZZ}{\mathbb{Z}}
\newcommand{\QQ}{\mathbb{Q}}
\newcommand{\NN}{\mathbb{N}}
\newcommand{\Af}{\mathbb{A}_f}
\newcommand{\Oh}{\hat{\mathcal{O}}}
\newcommand{\GL}{\operatorname{GL}}
\newcommand{\cj}[1]{\overline{#1}}
\newcommand{\ES}[1]{E_{#1}}                    % the product of the completions at the places of a finite set
\newcommand{\OS}[1]{O_{#1}}                    % its integral elements
\newcommand{\RS}[1]{(\ZZ/N)^{#1}}              % the residue ring
\newcommand{\st}{\ast}                         % the last index
\newcommand{\Ip}[1]{#1^{+}}                    % an index set with the last index added
\newcommand{\red}{\operatorname{red}}          % reduction modulo N
\newcommand{\secN}{\operatorname{sec}}         % the chosen section of the reduction of matrices
\newcommand{\uz}{u_0}                          % diag(1,...,1,1/N) at every place of the split level
\newcommand{\Mir}[1]{\operatorname{Mir}(#1)}   % the mirabolic group
\newcommand{\nat}[1]{#1^{\natural}}            % a function of residue matrices, extended by zero
\newcommand{\rrow}{\mathbf{r}}                 % the indicator of the last row
\newcommand{\ebas}{\mathbf{e}}                 % the last basis vector
\newcommand{\Fop}[1]{\mathsf{F}(#1)}           % the Gauss-sum operator of a representation
\newcommand{\nm}[2]{\operatorname{nm}_{#1}(#2)}% the operator named by a function, at a set of places
\newcommand{\emb}[2]{\operatorname{emb}_{#1}(#2)}
\newcommand{\Gf}[1]{G_f(#1)}                   % the finite adelic unitary group of a standard space
\newcommand{\Kz}[1]{K_0(#1)}
\newcommand{\Wk}[1]{W_{\kappa}(#1)}            % the space of a weight
\newcommand{\Alg}[2]{\mathcal{A}(#1;#2)}       % algebraic forms on an index set, with a last entry
\newcommand{\cls}[1]{\operatorname{cl}(#1)}    % the number of classes of a level
\newcommand{\jm}{\jmath}
\newcommand{\Oj}{O_{\jmath}}
\newcommand{\ind}[1]{\mathfrak{p}_{\jmath}(#1)}% the prime induced by an infinite place and the p-adic embedding
\newcommand{\evt}[1]{\operatorname{ev}_{\jmath,#1}}
\newcommand{\gauge}[2]{\operatorname{gauge}_{\jmath}(#1;#2)}
\newcommand{\Int}[2]{\operatorname{Int}_{\jmath}(#1;#2)}
\newcommand{\ND}{N(D)}
\newcommand{\piN}{\pi_N(D)}
\newcommand{\PiD}{\Pi(D)}
\newcommand{\Col}{\mathrm{Col}(D)}
\newcommand{\Colram}{\mathrm{Col}^{\mathrm{ram}}(D)}
\newcommand{\Kcol}[1]{K^{\mathrm{col}}_{#1}(D)}
\newcommand{\Stg}[1]{\mathrm{St}_{#1}(D)}
\newcommand{\XI}[1]{\underline{X}_{#1}(D)}
\newcommand{\XIo}[1]{\underline{X}^{\circ}_{#1}(D)}
\newcommand{\resY}{\operatorname{res}_{D}}
\newcommand{\bdU}{\operatorname{bd}_{D}}
\newcommand{\brE}[1]{\operatorname{br}_{D}(#1)}
\newcommand{\wnorm}[1]{\lVert #1\rVert^{2}_{\mathrm{W}}}
\newcommand{\inbox}[1]{\par\addvspace{6pt}\begingroup\setlength{\fboxsep}{7pt}\setlength{\fboxrule}{0.5pt}%
  \noindent\fcolorbox{rulegrey}{white}{\parbox{\dimexpr\linewidth-2\fboxsep-2\fboxrule\relax}{\fontsize{9.8}{13.4}\selectfont #1}}%
  \endgroup\par\addvspace{6pt}}

\begin{document}
\color{ink}

{\raggedright
{\sffamily\bfseries\fontsize{7.2}{9}\selectfont\color{leancol}\addfontfeatures{LetterSpace=8}READBACK NOTE 23\par}
\vspace{5pt}
{\fontsize{21.5}{25.5}\selectfont\bfseries Chains of definite theta lifts and the Gauss-sum operator at the level prime\par}
\vspace{5pt}
{\fontsize{12.4}{16}\selectfont\itshape\color{inksoft} One unit of the Lean artifact \ArtifactName, read back into
mathematics\par}}
\vspace{8pt}
{\color{rulegrey}\hrule height 0.6pt}
\vspace{7pt}
\thispagestyle{empty}

\textbf{Summary.} At the level prime $N$, which splits completely in the CM field, a theta lift between definite
unitary groups is governed by functions of residue matrices modulo~$N$. This note isolates the formal law that
one lift obeys there, defines the \kw{Gauss-sum operator} \m{\mathsf{F}}, a finite Fourier transform in the last
row followed by \m{\operatorname{diag}(1,\dots,1,N^{-1})}, and proves that \m{\mathsf{F}} slides through a chain of
lifts (Theorem~A). Theorem~B rewrites the sums of translates with denominator~$N$ that the seesaw identity produces
as data of exactly this kind. Theorem~C turns a jointly invariant combination of chains into a plain sum over
classes, and Theorem~D bounds, at an embedding above~$N$, what powers of \m{\mathsf{F}} cost in integrality.

\statement{leancol}{\fontsize{9.6}{13.2}\selectfont\textbf{What this note is.} The artifact \ArtifactName{} is a
body of Lean code whose final theorems state the Riemann hypothesis and its analogue for Dirichlet
$L$-functions. Its authors report that all files compile and that Lean's kernel accepts these theorems from its
standard axioms; nothing was compiled for this note, and no mathematician has read the argument through. The
headline is therefore a \emph{claim}, and this note does not say that it is proved: it says what one part of the
formal text asserts and how it argues. The part is \UnitInNote{} of the \UnitMembers{} declarations of the unit
\emph{Seesaw}: its files on names at a split level and on the characters of the lifts, which serve Theorem~A of
the write-up (short forms). The other declarations of the unit are treated in further notes of the unit (the index
sets and characters of the data; two consecutive lifts; the seesaw identity). Every statement below was compared, argument by argument, with its Lean text, and every proof is
written out, except where an \emph{Input} from another unit or a quoted result of the background library is named
with its owner. The Lean text read is \PinText; Appendix~A is the concordance of \ConcDecls{} declarations.}

\sect{The statements}{s:intro}

The theorems are stated in the settings that follow. The notions they use are defined in §\ref{s:setting}; the
place of each definition is given in the setting.

\begin{setting}[a split level]\label{s:split}
$E$ is a CM field with a CM type and Weil data, $N$ is a prime number, and $\Sigma$ is a finite set of finite
places of $E$ that is a split level for~$N$ (Definition~\ref{d:split}): $\Sigma$ is disjoint from its complex
conjugate, each of its places has norm~$N$ and uniformiser~$N$, and the places of $E$ dividing $N$ are those of
$\Sigma$ and their conjugates. We write \m{\ES{\Sigma}=\prod_{w\in\Sigma}E_w}, $\OS{\Sigma}$ for its ring of integral
elements and \m{\red\colon\OS{\Sigma}\to\RS{\Sigma}} for the reduction modulo~$N$ (Definition~\ref{d:red}). For a
finite index set $I$ we put \m{\Ip{I}=I\sqcup\{\st\}}, with $\st$ the \kw{last index}, and $\ebas$ is the last basis
vector of \m{(\RS{\Sigma})^{\Ip{I}}}. Functions $f$ of matrices over $\RS{\Sigma}$ give functions $\nat{f}$ of
matrices over $\ES{\Sigma}$; among them are the indicators $\bcol{y}$ of a last column and the characters
$\cchar{y}$ of a last row (Definition~\ref{d:fn}). For a complex representation $\rho$ of
\m{\GL_{\Ip{I}}(\ES{\Sigma})}, $\ram{\rho}$ is its set of ramified vectors (Definition~\ref{d:diag}) and $\Fop{\rho}$
its Gauss-sum operator (Definition~\ref{d:F}); $\uz$ is the matrix \m{\operatorname{diag}(1,\dots,1,N^{-1})} at
every place of $\Sigma$ (Definition~\ref{d:diag}).
\end{setting}

\statement{arith}{\phantomsection\label{t:A}\textbf{Theorem A} (the Gauss-sum operator slides through a chain of
lift laws)\textbf{.}\enspace In Setting~\ref{s:split}, let $I_0,I_1,I_2,\dots$ be finite index sets, let $\rho_j$ be
a complex representation of \m{\GL_{\Ip{I_j}}(\ES{\Sigma})} for each \m{j\ge0}, and let $L_j$ and \m{\varkappa_j\in\CC}
be such that, for every~$j$, $\law{\rho_j}{\rho_{j+1}}{L_j}{\varkappa_j}$ (Definition~\ref{d:law}). Let
\m{\varphi\in\ram{\rho_0}}, let \m{d_0,d_1,\dots} be complex numbers and let \m{p\ge0}. For \m{i\ge0} put
\[
A_i\;=\;d_i\cdot\rho_{2i+2}(\uz)\circ\sum_{y}L_{2i+1}[\cchar{y}]\circ L_{2i}[\bcol{y}]\;\colon\;\rho_{2i}\to\rho_{2i+2},
\]
the sum over all \m{y\in(\RS{\Sigma})^{\Ip{I_{2i+1}}}}, and let $o_i$ be the number of those $y$ that are non-zero at
every place: for every \m{w\in\Sigma} some entry of $y$ has non-zero component at~$w$. Then
\[
A_{p-1}\circ\dots\circ A_0\,(\varphi)\;=\;\Bigl(\prod_{i<p}d_i\,o_i\prod_{l<2i+2}\varkappa_l\Bigr)\cdot
L_{2p-1}[\bcol{\ebas}]\circ\dots\circ L_0[\bcol{\ebas}]\bigl(\Fop{\rho_0}^{\,p}\varphi\bigr).
\]}

\statement{arith}{\phantomsection\label{t:B}\textbf{Theorem B} (translates with denominator $N$ are named
data)\textbf{.}\enspace In Setting~\ref{s:split}, let $\Sigma'$ be a second finite set of finite places of $E$ that is
a split level for~$N$, let $I$, $I'$, $I''$ be finite index sets, and put \m{\mathcal{H}=\operatorname{Hom}(X_{\Ip{I}},X_{I'})} and
\m{\mathcal{H}'=\operatorname{Hom}(X_{I'},X_{\Ip{I''}})}, where $X_J$ is the standard definite hermitian space on an
index set~$J$, and let \m{v\in\mathcal{W}_f(\mathcal{H})} and \m{v'\in\mathcal{W}_f(\mathcal{H}')} be elements of the
finite parts of their Weil representations (Definition~\ref{d:names}). Then
\begin{align*}
&\transum{I}{I'}{\Sigma'}{N}{I''}\bigl(\projfin{\mathcal{H}}{\Sigma'}v\otimes\projfin{\mathcal{H}'}{\Sigma'}v'\bigr)\\
&\qquad=\;\sum_{z}\;\nm{\Sigma}{S\mapsto\bcol{z}(S\cdot\diagN{I}{\Sigma'})}\,v\;\otimes\;
\nm{\Sigma}{S'\mapsto\cchar{z}(\diagN{I''}{\Sigma\smallsetminus\Sigma'}\cdot S')}\,v',
\end{align*}
the sum over all \m{z\in(\RS{\Sigma})^{I'}}. Here \m{\transum{I}{I'}{\Sigma'}{N}{I''}} is the sum of the pairs of
translations by vectors with denominator~$N$ at $\Sigma'$, \m{\projfin{\mathcal{H}}{\Sigma'}} the projection of \m{\mathcal{W}_f(\mathcal{H})} to the vectors that
are vacua above \m{\Sigma'\cup\cj{\Sigma'}}, \m{\nm{\Sigma}{F}} the operator named at $\Sigma$ by a function~$F$ of
matrices over $\ES{\Sigma}$ (Definition~\ref{d:names}), and \m{\diagN{I}{\Sigma'}} the diagonal matrix with last
entry $N$ at the places of \m{\Sigma\cap\Sigma'} (Definition~\ref{d:diag}); the matrices $\diagN{I}{\Sigma'}$ and
$\diagN{I''}{\Sigma\smallsetminus\Sigma'}$ are formed for the split level~$\Sigma$.}

\begin{setting}[a chain of kernels]\label{s:chain}
$E$ is a CM field with a CM type. For \m{j\ge0}: $I_j$ is a finite index set and \m{\Gf{I_j}} the group of finite
adelic points of the unitary group of the standard definite hermitian space on~$I_j$; $\kappa$ is a weight and
$e_j$ a last entry, so that \m{\Alg{I_j}{e_j}} is the space of algebraic modular forms of weight \m{(\kappa,e_j)}
on \m{\Gf{I_j}}, with values in the space \m{\Wk{I_j}} of the weight; $K_j$ is a level (a compact open subgroup of
the stabiliser \m{\Kz{I_j}} of the standard lattice), $\mathcal{S}_j$ is a complex vector space, and
\m{\mathcal{K}_j} is a theta kernel from $I_j$ to $I_{j+1}$ with weights \m{(\kappa,e_j)}, \m{(\kappa,e_{j+1})},
level $K_j$ and data in $\mathcal{S}_j$ (Definition~\ref{d:kernel}), with its representation
\m{\omega_j} of \m{\Gf{I_j}\times\Gf{I_{j+1}}} on $\mathcal{S}_j$ and its jet matrix $\mathcal{J}_j$. The chain
\m{\chain{\mathcal{K}}{\Phi}_{j}\colon\Alg{I_0}{e_0}\to\Alg{I_j}{e_j}} of the first $j$ lifts with data
\m{\Phi=(\Phi_j)_j}, \m{\Phi_j\in\mathcal{S}_j}, is defined in Definition~\ref{d:chain}.
\end{setting}

\statement{arith}{\phantomsection\label{t:C}\textbf{Theorem C} (a jointly invariant combination of chains is a sum
over classes)\textbf{.}\enspace In Setting~\ref{s:chain}, let $M$ be a complex vector space, $\mathcal{X}$ a set,
\m{m\ge0} a natural number, \m{\mathcal{X}_0\subseteq\mathcal{X}} a finite subset, \m{L_x\colon\Wk{I_m}\to M} a
linear map and \m{\Phi_x=(\Phi_{x,j})_{j}} a family of data, for each \m{x\in\mathcal{X}}. Let
\m{\varphi\in\Alg{I_0}{e_0}}, let \m{P_j(g,u)} be a property of a point \m{g\in\Gf{I_j}} and a vector
\m{u\in\Wk{I_j}} for each~$j$, let \m{h\in\Gf{I_m}} and let \m{A\subseteq M}. Assume:
\begin{enumerate}[label=(\alph*)]
\item the family \m{(L_x,\Phi_x)_{x\in\mathcal{X}_0}} is jointly invariant of length $m$ for $\mathcal{K}$
  (Definition~\ref{d:joint});
\item for every \m{j<m} and every \m{g\in\Gf{I_j}}, the group of rational points in \m{gK_jg^{-1}} is trivial;
\item $\varphi$ is fixed by $K_0$ under right translation;
\item for all \m{x\in\mathcal{X}_0}, \m{j<m}, points \m{g\in\Gf{I_j}}, \m{g'\in\Gf{I_{j+1}}} and vectors~$u$:
  \m{P_j(g,u)} implies \m{P_{j+1}\bigl(g',\mathcal{J}_j(\omega_j(g,g')\Phi_{x,j})\,u\bigr)};
\item \m{P_0(g,\varphi(g))} for every \m{g\in\Gf{I_0}};
\item for every linear map \m{\beta\colon\mathcal{S}_{m-1}\otimes\dots\otimes\mathcal{S}_0\to\Wk{I_m}}: if
  \m{P_m\bigl(h,\beta(\Phi_{x,m-1}\otimes\dots\otimes\Phi_{x,0})\bigr)} for all \m{x\in\mathcal{X}_0}, then
  \m{\sum_{x\in\mathcal{X}_0}L_x\bigl(\beta(\Phi_{x,m-1}\otimes\dots\otimes\Phi_{x,0})\bigr)\in A}.
\end{enumerate}
Then \m{\sum_{x\in\mathcal{X}_0}L_x\bigl(\chain{\mathcal{K}}{\Phi_x}_{m}(\varphi)(h)\bigr)} is a sum of
\m{\prod_{j<m}\cls{K_j}} elements of~$A$, where \m{\cls{K_j}} is the number of double cosets of \m{\Gf{I_j}} modulo the
rational points on the left and $K_j$ on the right.}

\begin{setting}[the data at the level prime]\label{s:data}
\m{P=(N_0,d)} are parameters and $D$ is admissible data for~$P$ (Chapter~III), with components: the scales $x$,
$x_1$, $b$, the exponent $\mu$, the number of lifts $m$, the rank $r$, the thick set, the primes $g_0$, $g_1$ and
the integer $\pi_g$. Attached to $D$ are the CM field \m{E(D)} with its number \m{f(D)} of infinite places, the
level prime \m{\ND} and the integer \m{\piN} of \m{E(D)} with \m{\piN\cj{\piN}=\ND}, the set \m{\PiD} of places
dividing $\piN$, the number \m{\varrho(D)} with \m{m=2\varrho(D)+2}, the index sets
\m{\XI{j}=\XIo{j}\sqcup\{\st\}} of the standard definite spaces \m{X_j(D)}, the spaces \m{\Stg{j}} of algebraic
forms on \m{X_j(D)} (stage~$j$), the level \m{\Kcol{j}} of stage $j$, the space \m{\Col\subseteq\Stg{0}} of columns
and its subspace \m{\Colram} of ramified columns (Definition~\ref{d:columns}, which recalls them from the notes on
the data and on the columns). Moreover \m{\jm\colon\CC\to\CC_{\ND}} is a ring homomorphism, $\Oj$ is the ring of the
complex numbers $z$ with \m{|\jm(z)|\le1}, \m{\nonind{D}} is the set of non-induced places above \m{\ND}, a split
level for \m{\ND}, \m{\dpt{D}{j}} is the diagonal point of stage~$j$ and \m{\nrep{D}{j}} the representation of
\m{\GL_{\XI{j}}(\ES{\nonind{D}})} on \m{\Stg{j}} (Definition~\ref{d:data}); \m{\Int{\mathcal{V}}{\mathsf{d}}} is the
$\Oj$-module of the forms of a space $\mathcal{V}$ that are $\jm$-integral at a point~$\mathsf{d}$
(Definition~\ref{d:gauge}).
\end{setting}

\statement{arith}{\phantomsection\label{t:D}\textbf{Theorem D} (powers of the Gauss-sum operator on integral
columns)\textbf{.}\enspace In Setting~\ref{s:data}, with the structure of a split level for \m{\ND} on \m{\nonind{D}}
given by Lemma~\ref{l:nonind}, there is a natural number $e$ with
\[e\;\le\;4\,m\cdot f(D)\cdot b\cdot\dim_{\CC}\Col\]
such that for every \m{n\ge0} and all \m{\Phi_1,\dots,\Phi_n\in\Int{\Col}{1}} there is an \m{n\times n} matrix $u$
over $\Oj$ whose determinant divides \m{\ND^{\,e}} in $\Oj$ and for which every \m{\psi_i=\sum_{k}u_{ik}\Phi_k}
satisfies
\[\Fop{\nrep{D}{0}}^{\,\varrho(D)+1}\,\psi_i\;\in\;\Int{\Col}{\dpt{D}{0}}.\]}

\sect{The setting}{s:setting}
\begin{conventions}\phantomsection\label{c:conventions}
$E$ is a CM field with a fixed CM type; \m{x\mapsto\cj{x}} is its complex conjugation. $\Af$ is the ring of finite
adeles of $E$ and $\Oh$ its ring of integers. Finite places of $E$ are written $w$, $v$; $\cj{w}$ is the
conjugate place, $E_w$ the completion with ring of integers $O_w$, residue field $\mathbb{F}_w$ and norm
\m{\mathrm{N}w=\#\mathbb{F}_w}. Whenever the additive character $\psi_0$ or a Weil representation occurs, $E$ carries
Weil data: an element $\mathbf{j}$ of $E$ with \m{\cj{\mathbf{j}}=-\mathbf{j}}, a finite idele $\alpha$ and a
splitting character $\upsilon$; then \m{\psi_0(t)=\psi_f(\mathbf{j}\alpha\cj{\alpha}\,t)}, with $\psi_f$ the standard
additive character of the finite adeles, is an additive character of $\Af$ that is trivial on~$\Oh$ and whose
values are roots of unity. An infinite place $\tau$ of $E$ is identified with the embedding \m{E\to\CC} of the CM
type above it. $N$ is a prime number; $\Sigma$ and $\Sigma'$ are finite sets of finite places of~$E$. That such a
set is a split level is a hypothesis; it is stated in each statement or in the setting it cites, and Lean carries
it as an instance argument. Index sets $I$, $I'$, $I''$, $I_j$, $J$ are finite index sets, and
\m{\Ip{I}=I\sqcup\{\st\}}. $R$ is a commutative ring. Representations are on complex vector spaces. For a finite
index set $I$, $X_I$ is the standard definite hermitian space \m{E^{I}} with the form \m{\sum_i\cj{x_i}y_i},
\m{\Gf{I}} the group of finite adelic points of its unitary group, and \m{\Kz{I}=\Gf{I}\cap\GL_I(\Oh)}. In
Setting~\ref{s:chain}: the $\mathcal{S}_j$ are the data spaces, complex vector spaces; the $K_j$ are levels;
$\kappa$ is the weight and the $e_j$ are the last entries; $\Phi$ denotes data; $M$ is a complex vector space (so
also a rational one, which the class integral needs). In Setting~\ref{s:data}: \m{P=(N_0,d)} are parameters, $D$
is admissible data, $\ell$ is a prime number and \m{\jm\colon\CC\to\CC_{\ell}} a ring homomorphism, with
\m{\ell=\ND} where the statement says so.
\end{conventions}

\subhead{Split levels.}
\begin{definition}[split level; residues over $N$]\label{d:split}
Let $\Sigma$ be a finite set of finite places of $E$ and $N$ a natural number. We say that
\m{\splitlevel{\Sigma}{N}} if: (a)~\m{\cj{w}\notin\Sigma} for every \m{w\in\Sigma}; (b)~\m{\mathrm{N}w=N} for every
\m{w\in\Sigma}; (c)~$N$ is a uniformiser of $E_w$ for every \m{w\in\Sigma}; (d)~a finite place $w$ of $E$ divides
$N$ if and only if \m{w\in\Sigma} or \m{\cj{w}\in\Sigma}. For any $\Sigma$ and $N$ put
\m{\ES{\Sigma}=\prod_{w\in\Sigma}E_w} and let \m{\OS{\Sigma}=\prod_{w\in\Sigma}O_w} be its ring of integral elements.
For \m{N>0} and \m{y\in\RS{\Sigma}} let \m{\tilde y_w\in\{0,\dots,N-1\}} be the least non-negative representative of
$y_w$; then
\m{\rN{y}\in\Af} is the finite adele with component \m{\tilde y_w/N} at each \m{w\in\Sigma} and $0$ at the other
places, and, for $E$ with Weil data,
\[\psiN{y}\;=\;\psi_0\bigl(-\rN{y}\bigr).\]
\end{definition}
\leanline{\lean{IsSplitLevel}\sep \lean{residueOverN}\sep \lean{residueOverNapply} (the formula for the components,
a restatement of the definition)\sep \lean{psiN}}

\begin{definition}[reduction modulo $N$]\label{d:red}
Let \m{\splitlevel{\Sigma}{N}}, with $N$ prime; we say for short that $\Sigma$ is a split level.
For \m{w\in\Sigma} the residue field $\mathbb{F}_w$ has $N$ elements by~(b), so there is exactly one ring isomorphism
\m{\mathbb{F}_w\to\ZZ/N}. The \kw{reduction} \m{\red\colon\OS{\Sigma}\to\RS{\Sigma}} is the ring homomorphism whose
component at $w$ is the residue map \m{O_w\to\mathbb{F}_w} followed by this isomorphism. For a finite index set $I$
we regard \m{\GL_I(\OS{\Sigma})} as a subgroup of \m{\GL_I(\ES{\Sigma})}, write again
\m{\red\colon\GL_I(\OS{\Sigma})\to\GL_I(\RS{\Sigma})} for the entrywise reduction, and fix a map
\m{\secN\colon\GL_I(\RS{\Sigma})\to\GL_I(\OS{\Sigma})} with \m{\red\circ\secN=\mathrm{id}}. It exists: an entry
\m{y\in\RS{\Sigma}} lifts to \m{(\tilde y_w)_{w}\in\OS{\Sigma}}, and a matrix over $\OS{\Sigma}$ whose reduction is
invertible is invertible, because an element of $\OS{\Sigma}$ with invertible reduction is a unit. The
\kw{last basis vector} \m{\ebas\in(\RS{\Sigma})^{\Ip{I}}} is the function \m{(i,w)\mapsto1} if \m{i=\st} and
\m{(i,w)\mapsto0} otherwise, for \m{i\in\Ip{I}} and \m{w\in\Sigma}.
\end{definition}
\leanline{\lean{residueEquiv}\sep \lean{reduction}\sep \lean{inclGL}\sep \lean{reductionGL}\sep \lean{liftGL}\sep
\lean{lastBasisVector}}

\begin{definition}[diagonal elements, the mirabolic group, ramified vectors]\label{d:diag}
Let $\Sigma$ be a split level for the prime~$N$, $I$ a finite index set and $\Sigma'$ a finite set of finite
places. Then \m{\diagN{I}{\Sigma'}\in\GL_{\Ip{I}}(\ES{\Sigma})} is the diagonal matrix with entries $1$ at the
indices of $I$ and, at the last index, the element of $\ES{\Sigma}$ whose component is $N$ at the
\m{w\in\Sigma\cap\Sigma'} and $1$ at the other \m{w\in\Sigma}. We put
\[\uz\;=\;\diagN{I}{\Sigma}^{-1},\]
the matrix \m{\operatorname{diag}(1,\dots,1,N^{-1})} at every place of~$\Sigma$ (we write $\uz$ in every rank). The
\kw{mirabolic group} \m{\Mir{I}\subseteq\GL_{\Ip{I}}(\ES{\Sigma})} consists of the \m{k\in\GL_{\Ip{I}}(\OS{\Sigma})}
with \m{\red(k)\,\ebas=\ebas}, that is, whose last column is congruent to $\ebas$ modulo~$N$. For a complex
representation $\rho$ of \m{\GL_{\Ip{I}}(\ES{\Sigma})}, a vector $\varphi$ of $\rho$ is \kw{ramified} if for every
\m{w\in\Sigma} there is a unit $x$ of $\OS{\Sigma}$ with \m{\red(x)_v=1} for all \m{v\in\Sigma}, \m{v\ne w}, such
that $\varphi$ lies in the image of \m{1-\rho(x\cdot1)}, where \m{x\cdot1} is the scalar matrix; \m{\ram{\rho}} is
the set of ramified vectors.
\end{definition}
\leanline{\lean{diagN}\sep \lean{diagNInv}\sep \lean{mirabolic}\sep \lean{ramified}}

\begin{definition}[functions of residue matrices]\label{d:fn}
Let $\Sigma$ be a split level for the prime~$N$ and $I$, $I'$ finite index sets. For a complex function $f$ of
\m{I'\times I} matrices over $\RS{\Sigma}$, \m{\nat{f}} is the function of \m{I'\times I} matrices $S$ over
$\ES{\Sigma}$ with \m{\nat{f}(S)=f(\red S)} if all entries of $S$ lie in $\OS{\Sigma}$, and \m{\nat{f}(S)=0}
otherwise; \m{f\mapsto\nat{f}} is $\CC$-linear. Three families are used:
\begin{enumerate}[label=(\roman*)]
\item $\rrow$, on \m{\Ip{I'}\times\Ip{I}} matrices: the indicator of the integral $S$ whose last row is congruent
  to $\ebas$;
\item \m{\bcol{y}}, for \m{y\in(\RS{\Sigma})^{I'}}, on \m{I'\times\Ip{I}} matrices: the indicator of the integral
  $S$ whose last column is congruent to~$y$;
\item \m{\cchar{y}}, for \m{y\in(\RS{\Sigma})^{I}}, on \m{\Ip{I'}\times I} matrices, for $E$ with Weil data:
  \m{\cchar{y}(S)=\psiN{\textstyle\sum_{i\in I}r_iy_i}} for integral $S$ with last row \m{r=(r_i)_{i\in I}} modulo~$N$,
  and \m{\cchar{y}(S)=0} for $S$ not integral.
\end{enumerate}
\end{definition}
\leanline{\lean{liftFn}\sep \lean{lastRowDatum}\sep \lean{lastColDatum}\sep \lean{charDatum}}

\begin{lemma}\label{l:kofy}
Let $I$ be a finite index set and $R$ a commutative ring. There is a map $k$ from \m{R^{\Ip{I}}} to
\m{\GL_{\Ip{I}}(R)} such that every vector $y$ whose last entry \m{y_{\st}} is a unit of $R$ is both the last
column and the last row of \m{k(y)}.
\end{lemma}
\begin{proof}
Let $y$ have \m{y_{\st}\in R^{\times}} and let $C$ be the identity matrix with its last column replaced by~$y$. Its
determinant is \m{y_{\st}}, and its last row is \m{(0,\dots,0,y_{\st})}. Hence
\m{(CC^{\mathsf{T}})_{i\st}=\sum_{a}C_{ia}C_{\st a}=y_i\,y_{\st}}, the matrix \m{CC^{\mathsf{T}}} is symmetric, and
\m{k(y)=y_{\st}^{-1}CC^{\mathsf{T}}} is invertible with last column and last row~$y$. For the other $y$ put
\m{k(y)=1}.
\end{proof}
\leanline{\lean{existsGLlastcolandrow}\sep In the Lean statement the variable $y$ of the surrounding section is
listed as an argument and not used: the statement quantifies over its own~$y$.}

\begin{definition}[the Gauss-sum operator]\label{d:F}
Let $E$ have Weil data, let $\Sigma$ be a split level for the prime~$N$, $I$ a finite index set and $\rho$ a complex
representation of \m{\GL_{\Ip{I}}(\ES{\Sigma})}. With $k$ the map of Lemma~\ref{l:kofy} for \m{R=\RS{\Sigma}} (the
one chosen by the formal text) and \m{\tilde k(y)=\secN(k(y))},
\[
\Fop{\rho}\;=\;\rho(\uz)\circ\sum_{y}\psiN{y_{\st}}\cdot\rho\bigl(\tilde k(y)\bigr),
\]
the sum over the \m{y\in(\RS{\Sigma})^{\Ip{I}}} whose last entry \m{y_{\st}} is a unit of $\RS{\Sigma}$. It is an
endomorphism of the space of~$\rho$: a finite Fourier transform in the last row and column, followed by~$\uz$.
\end{definition}
\leanline{\lean{fourierOp}}

\begin{definition}[lift law]\label{d:law}
Let $E$ have Weil data, let $\Sigma$ be a split level for the prime~$N$, $I$ and $I'$ finite index sets, $\rho$
and $\rho'$ complex representations of \m{\GL_{\Ip{I}}(\ES{\Sigma})} and \m{\GL_{\Ip{I'}}(\ES{\Sigma})}, $L$ a
$\CC$-linear map from the complex functions of \m{\Ip{I'}\times\Ip{I}} matrices over $\ES{\Sigma}$ to the linear maps
\m{\rho\to\rho'} (we write \m{L[F]} for its value at a function~$F$), and \m{\varkappa\in\CC}. We say that
$\law{\rho}{\rho'}{L}{\varkappa}$ if
\begin{enumerate}[label=(\roman*)]
\item (\kw{integral equivariance}) for all \m{a\in\GL_{\Ip{I}}(\OS{\Sigma})}, \m{b\in\GL_{\Ip{I'}}(\OS{\Sigma})}, every
  function $f$ of \m{\Ip{I'}\times\Ip{I}} matrices over $\RS{\Sigma}$ and every vector $\varphi$ of $\rho$:
  \[\rho'(b)\Bigl(L\bigl[\nat{(T\mapsto f(\red(b)\,T\,\red(a)))}\bigr]\bigl(\rho(a)\varphi\bigr)\Bigr)\;=\;L[\nat{f}]\,\varphi;\]
\item (\kw{the law of $\uz$}) \m{\rho'(\uz)\bigl(L[\rrow]\,\varphi\bigr)=\varkappa\cdot L[\bcol{\ebas}]\bigl(\rho(\uz)\varphi\bigr)}
  for every vector $\varphi$ of~$\rho$.
\end{enumerate}
\end{definition}
\leanline{\lean{LiftLaw}}

\subhead{Kernels and chains.}
\begin{definition}[theta kernels and their lifts; recalled from the background library]\label{d:kernel}
In Setting~\ref{s:chain}, a \kw{theta kernel} $\mathcal{K}$ from $I$ to $I'$, with weights \m{(\kappa,e)},
\m{(\kappa,e')}, level \m{K\subseteq\Gf{I}} and data in a complex vector space $\mathcal{S}$, consists of a
representation $\omega$ of \m{\Gf{I}\times\Gf{I'}} on $\mathcal{S}$ in which every vector has an open stabiliser, and
a linear map $\mathcal{J}$, the \kw{jet matrix}, from $\mathcal{S}$ to the linear maps \m{\Wk{I}\to\Wk{I'}}, such
that \m{\mathcal{J}(\omega(\gamma,\delta)\Phi)(\gamma u)=\delta\,\mathcal{J}(\Phi)(u)} for all rational points
$\gamma$, $\delta$ of the two unitary groups, acting on \m{\Wk{I}}, \m{\Wk{I'}} through the weights. Its
\kw{lift} with datum \m{\Phi\in\mathcal{S}} is the linear map \m{\Alg{I}{e}\to\Alg{I'}{e'}},
\[
\mathcal{K}[\Phi](\varphi)(h)\;=\;\int_{\Gamma\backslash\Gf{I}}\mathcal{J}\bigl(\omega(g,h)\Phi\bigr)\bigl(\varphi(g)\bigr)\,dg ,
\]
with $\Gamma$ the group of rational points and the measure that gives $K$ volume~$1$; the integrand is right
invariant under an open subgroup, so the integral is a finite sum over classes, divided by an index. The
lift is equivariant: \m{\mathcal{K}[\omega(a,b)\Phi](a\cdot\varphi)=b\cdot\mathcal{K}[\Phi](\varphi)} for
\m{a\in\Gf{I}}, \m{b\in\Gf{I'}}, where \m{(a\cdot\varphi)(g)=\varphi(ga)} is right translation.
\end{definition}

\begin{definition}[chains of lifts]\label{d:chain}
In Setting~\ref{s:chain}, let \m{\mathcal{K}=(\mathcal{K}_j)_{j}} be the kernels and \m{\Phi=(\Phi_j)_j} a family
of data, \m{\Phi_j\in\mathcal{S}_j}. Then \m{\chain{\mathcal{K}}{\Phi}} is the family of linear maps
\[
\chain{\mathcal{K}}{\Phi}_{j}\;=\;\mathcal{K}_{j-1}[\Phi_{j-1}]\circ\dots\circ\mathcal{K}_0[\Phi_0]\;\colon\;\Alg{I_0}{e_0}\to\Alg{I_j}{e_j},
\]
the identity for \m{j=0}. For a family \m{\Phi=(\Phi_i)_i} with
\m{\Phi_i\in\mathcal{S}_{2i}\otimes\mathcal{S}_{2i+1}}, \m{\chaintwo{\mathcal{K}}{\Phi}} is the family
\m{\chaintwo{\mathcal{K}}{\Phi}_{p}\colon\Alg{I_0}{e_0}\to\Alg{I_{2p}}{e_{2p}}} of the composites of the first $p$
\kw{pairs of lifts}, the $i$-th pair being the linear map that sends a pure tensor \m{s\otimes s'} to
\m{\mathcal{K}_{2i+1}[s']\circ\mathcal{K}_{2i}[s]}, evaluated at~$\Phi_i$. For a family \m{a=(a_j)_j} of points
\m{a_j\in\Gf{I_j}} and \m{i\ge0},
\[
\twoact{\mathcal{K}}{a}{i}\;=\;\omega_{2i}(a_{2i},a_{2i+1})\otimes\omega_{2i+1}(a_{2i+1},a_{2i+2}),
\]
an endomorphism of \m{\mathcal{S}_{2i}\otimes\mathcal{S}_{2i+1}}.
\end{definition}
\leanline{\lean{chain}\sep \lean{chaina}\sep \lean{twoKernelAction}}

\begin{definition}[joint invariance]\label{d:joint}
In Setting~\ref{s:chain}, let $M$ be a complex vector space, $\mathcal{X}$ a set, \m{m\ge0},
\m{\mathcal{X}_0\subseteq\mathcal{X}} finite, \m{L_x\colon\Wk{I_m}\to M} linear and \m{\Phi_x=(\Phi_{x,j})_j} data, for
\m{x\in\mathcal{X}}. The family \m{(L_x,\Phi_x)_{x\in\mathcal{X}_0}} is \kw{jointly invariant of length $m$ for
$\mathcal{K}$} if for every family \m{a=(a_j)_j} with \m{a_j\in K_j} for all $j$ and \m{a_j=1} for \m{j\ge m}, and
every linear map \m{\beta\colon\mathcal{S}_{m-1}\otimes\dots\otimes\mathcal{S}_0\to\Wk{I_m}},
\[
\sum_{x\in\mathcal{X}_0}L_x\Bigl(\beta\bigl(\textstyle\bigotimes_{j<m}\omega_j(a_j,a_{j+1})\Phi_{x,j}\bigr)\Bigr)
\;=\;\sum_{x\in\mathcal{X}_0}L_x\Bigl(\beta\bigl(\textstyle\bigotimes_{j<m}\Phi_{x,j}\bigr)\Bigr).
\]
In words: the tensor \m{\sum_{x}L_x\otimes\Phi_{x,m-1}\otimes\dots\otimes\Phi_{x,0}} is invariant under the joint
action of \m{\prod_{j<m}K_j}.
\end{definition}
\leanline{\lean{IsJointlyInvariant}\sep In the Lean definition the family $a$ of the surrounding section is listed
as an argument and not used: the definition quantifies over its own family~$a$ of elements of the levels.}

\subhead{Split places inside the adelic group; names.}
\begin{definition}[embedding at a split level; determinant characters]\label{d:emb}
Let $\Sigma$ be a finite set of finite places of $E$ with \m{\cj{w}\notin\Sigma} for \m{w\in\Sigma} (in the formal
text: a split level for some~$N$) and $I$ a finite index set. For \m{g\in\GL_I(\ES{\Sigma})},
\m{\emb{\Sigma}{g}\in\Gf{I}} is the element whose component is $g_w$ at each \m{w\in\Sigma}, the inverse of the
conjugate transpose of $g_w$ at $\cj{w}$, and $1$ at the other places; \m{g\mapsto\emb{\Sigma}{g}} is a group
homomorphism. For a Hecke character $\eta$ of $E$, with local components $\eta_w$, and \m{k\in\NN},
\m{\detchar{\eta}{k}} is the character
\[
\detchar{\eta}{k}(g)\;=\;\prod_{w\in\Sigma}\eta_w(\det g_w)\,\lvert\det g_w\rvert_w^{-k/2}
\]
of \m{\GL_I(\ES{\Sigma})} with values in \m{\CC^{\times}}; here \m{\lvert\,\cdot\,\rvert_w} is the normalised absolute
value of~$E_w$.
\end{definition}
\leanline{\lean{embedGL}\sep \lean{detCharAt}}

\begin{definition}[names, projections, sums of translates; the first two recalled]\label{d:names}
Let $E$ have Weil data and let $V$ be a hermitian space over $E$ of the kind used in this proof (in this note:
\m{V=\operatorname{Hom}(X_I,X_{I'})} for finite index sets $I$, $I'$, whose finite adelic vectors are the
\m{I'\times I} matrices over~$\Af$). Its Weil representation has a \kw{finite part} \m{\mathcal{W}_f(V)}, spanned by
the translates \m{\varrho_f(t)\,\mathrm{vac}} of the vacuum by the finite Heisenberg operators. For a set $\Sigma'$
of finite places, \m{\projfin{V}{\Sigma'}} is the restriction to \m{\mathcal{W}_f(V)} of the projection
\m{P^{\Sigma'}}, the mean of \m{\varrho_f(t)} over the integral vectors $t$ supported above
\m{\Sigma'\cup\cj{\Sigma'}}; its image consists of the vectors that are vacua above these places. For a split level
$\Sigma$ for the prime $N$ and a function $F$ of \m{I'\times I} matrices over $\ES{\Sigma}$, \m{\nm{\Sigma}{F}} is the
endomorphism of \m{\mathcal{W}_f(\operatorname{Hom}(X_I,X_{I'}))} named by~$F$: the operator that the background
library attaches to the function \m{t\mapsto F(t_{\Sigma})} of finite adelic matrices, \m{t_\Sigma} being the
components of $t$ at $\Sigma$; the map \m{(F,v)\mapsto\nm{\Sigma}{F}\,v} is bilinear. Finally, for finite index sets
$I$, $I'$, $I''$, a finite set $\Sigma$ of finite places and a prime~$N$,
\[
\transum{I}{I'}{\Sigma}{N}{I''}\;=\;\sum_{y}\;\varrho_f\bigl(c(y)\bigr)\otimes\varrho_f\bigl(r(y)\bigr),
\]
the sum over \m{y\in(\RS{\Sigma})^{I'}}: an endomorphism of
\m{\mathcal{W}_f(\mathcal{H})\otimes\mathcal{W}_f(\mathcal{H}')} for the spaces
\m{\mathcal{H}=\operatorname{Hom}(X_{\Ip{I}},X_{I'})} and \m{\mathcal{H}'=\operatorname{Hom}(X_{I'},X_{\Ip{I''}})}, where
\m{c(y)} is the \m{I'\times\Ip{I}} matrix with last column \m{(-\rN{y_i})_{i\in I'}} and zero elsewhere, and
\m{r(y)} the \m{\Ip{I''}\times I'} matrix with last row \m{(-\cj{\rN{y_i}})_{i\in I'}} and zero elsewhere.
\end{definition}
\leanline{\lean{name}\sep \lean{projOnFinPart}\sep \lean{translateSum}}

\subhead{The data.}
\begin{definition}[columns and stages; recalled from other notes]\label{d:columns}
In Setting~\ref{s:data}: \m{\Stg{j}} is the space of algebraic modular forms on \m{\Gf{\XI{j}}} of the weight
\m{(\kappa_D,e_{D,j})} that Chapter~V attaches to stage~$j$. The level \m{\Kcol{j}} consists of the
\m{k\in\Kz{\XI{j}}} that reduce to $1$ at the chosen places of $D$ and whose last row is congruent to
\m{(0,\dots,0,1)} modulo $\piN$; equivalently (Input~2) $k$ lies in the chosen level and, at every \m{w\in\PiD},
every entry of the last row of \m{k_w-1} is divisible by a uniformiser. \m{\Col} is the space of the forms of
\m{\Stg{0}} fixed by \m{\Kcol{0}}, and \m{\Colram} is the image of \m{\prod_{w\in\PiD}(1-z_w)} on $\Col$, $z_w$ being
right translation by the scalar matrix of a fixed integral norm-one adele $\nu_w$, equal to $1$ away from $w$ and
$\cj{w}$, whose component at $w$ generates the units of the residue field (Input~3). The \kw{restriction}
\m{\resY} maps \m{\Stg{m}} to algebraic forms on the definite space \m{Y_m(D)} of rank one less:
\[
\resY(\varphi)(h)\;=\;\Bigl(\prod_{\tau}\tau\bigl(\cj{\piN}\bigr)^{-\brE{\tau}}\Bigr)\cdot B_D\,\varphi\bigl(\bdU(h)\bigr),
\]
the product over the infinite places $\tau$ of \m{E(D)}, each with its embedding of the CM type; here
\m{\bdU(h)=\operatorname{diag}(h,1)}, $B_D$ is the integer branching matrix of the columns and
\m{\brE{\tau}\in\ZZ} the branching exponent (Input~4).
\end{definition}

\begin{definition}[gauges and integral forms; recalled from the background library]\label{d:gauge}
Let $I$ be a finite index set, $\ell$ a prime and \m{\jm\colon\CC\to\CC_\ell} a ring homomorphism. For an infinite
place $\tau$ of $E$, \m{\ind{\tau}} is the prime of $E$ above $\ell$ induced by the embedding
\m{\jm\circ\tau\colon E\to\CC_\ell}, and \m{\evt{\tau}\colon\Af\to\CC_\ell} is the component at \m{\cj{\ind{\tau}}}
followed by the continuous extension of \m{x\mapsto\jm(\tau(\cj{x}))}. The \kw{gauge} of a vector
\m{u\in\Wk{I}} at \m{g\in\Gf{I}} is
\[
\gauge{g}{u}\;=\;\bigl\lVert\,\varrho_{\kappa,e}\bigl((\evt{\tau}(g^{-1}))_{\tau}\bigr)\,\jm(u)\,\bigr\rVert,
\]
the largest $\ell$-adic absolute value of a coordinate, where \m{\varrho_{\kappa,e}} is the representation of
\m{\prod_\tau\GL_I} of weight \m{(\kappa,e)} and $\jm$ is applied to the coordinates of~$u$. An algebraic form
$\varphi$ \kw{has gauge at most $C$ at a point} \m{\mathsf{d}\in\Gf{I}} if \m{\gauge{g\mathsf{d}}{\varphi(g)}\le C}
for every \m{g\in\Gf{I}}. For a complex space $\mathcal{V}$ of algebraic forms, \m{\Int{\mathcal{V}}{\mathsf{d}}} is
the $\Oj$-module of the \m{\varphi\in\mathcal{V}} that have gauge at most $1$ at~$\mathsf{d}$; for
\m{\mathsf{d}=1} these are the \kw{$\jm$-integral} forms. For a form $\varphi$ and a ring homomorphism
\m{\sigma\colon\CC\to\CC}, \m{\wnorm{\sigma\circ\varphi(g)}} is the squared norm, for the Gram matrix of the weight, of
the vector with coordinates \m{\sigma(\varphi(g)_i)}.
\end{definition}

\begin{definition}[the non-induced places and the representation at the level prime]\label{d:data}
In Setting~\ref{s:data}, but for any prime $\ell$ and ring homomorphism \m{\jm\colon\CC\to\CC_\ell},
\[
\nonind{D}\;=\;\bigl\{\cj{\ind{\tau}}\;:\;\tau\text{ an infinite place of }E(D)\bigr\},
\]
the set of the complex conjugates of the induced primes; it depends on $\ell$ and~$\jm$. Now let \m{\ell=\ND} and
let \m{\nonind{D}} be a split level for $\ND$ (which holds, Lemma~\ref{l:nonind}). For \m{j\ge0},
\[
\dpt{D}{j}\;=\;\emb{\nonind{D}}{\diagN{\XIo{j}}{\PiD}}\;\in\;\Gf{\XI{j}},
\]
the point supported above $\ND$ with component \m{\operatorname{diag}(1,\dots,1,\ND)} at the \m{w\in\nonind{D}} that
divide $\piN$ and $1$ at the other \m{w\in\nonind{D}}; \m{\cemb{D}{j}} is the homomorphism
\[
\cemb{D}{j}\colon\GL_{\XI{j}}(\ES{\nonind{D}})\to\Gf{\XI{j}},\qquad g\mapsto\dpt{D}{j}\cdot\emb{\nonind{D}}{g}\cdot\dpt{D}{j}^{-1};
\]
and \m{\nrep{D}{j}} is the representation of \m{\GL_{\XI{j}}(\ES{\nonind{D}})} on \m{\Stg{j}} in which $g$ acts by
right translation by \m{\cemb{D}{j}(g)}.
\end{definition}
\leanline{\lean{nonInducedPlaces}\sep \lean{diagPointN}\sep \lean{conjEmbedGL}\sep \lean{nonInducedRep}}

\subhead{What is taken for granted.}
\inbox{\textbf{Inputs from other units} (as the unit's dossier reads them).\par
\textbf{1} (unit \emph{Columns}). The set \m{\PiD} of the places dividing $\piN$ is a split level for~$\ND$.\par
\textbf{2} (unit \emph{Columns}). \m{k\in\Kcol{j}} if and only if $k$ lies in the chosen level and, at every
\m{w\in\PiD}, every entry of the last row of \m{k_w-1} has valuation at most \m{e^{-1}}.\par
\textbf{3} (unit \emph{Columns}). For each \m{w\in\PiD}, every ramified column $\varphi$ has the form
\m{\psi-z_w\psi} for some column~$\psi$; the adele $\nu_w$ is integral, of norm one, and equal to $1$ away from
$w$ and~$\cj{w}$.\par
\textbf{4} (unit \emph{Columns}, from the background library). The branching matrix $B_D$ has integer entries;
it intertwines the action of \m{\operatorname{diag}(g,s)}, in the weight of stage~$m$, with a power of $s$ times the
action of $g$ in the weight of \m{Y_m(D)}, the exponent at $\tau$ being the one that enters \m{\brE{\tau}}; and it
satisfies a size bound in the Gram norms of the two weights. The reference for this branching from rank
\m{n+1} to rank $n$ is to be confirmed.\par
\textbf{5} (the note on the index sets and characters of the data, this unit; unit \emph{Columns}; Chapter~III).
The level prime \m{\ND} is congruent to $1$ modulo the conductor of \m{E(D)}; so it splits completely in \m{E(D)},
each place above it has norm \m{\ND} and uniformiser \m{\ND}, none is fixed by complex conjugation, and the group of
infinite places permutes the places above a place of $K$ simply transitively (Chapter~III). A place dividing
\m{\ND} is not a chosen place of~$D$, and \m{\lvert\XI{j}\rvert=n_j(D)} is even (unit \emph{Columns}).}

\inbox{\textbf{Quoted from the background library} (\emph{Lib}). The catalogue of the write-up lists the following, as far
as it contains them, as elementary facts or as definitions, without a reference to the literature; where a
reference would be expected we write that it is to be confirmed.\par
\textbf{Q1} (class integrals; the catalogue's file on class means). If all the groups of rational points in the
conjugates of $K$ are trivial and $F$ is right $K$-invariant with values in a set $A$, then the integral of $F$ over
the classes is a sum of elements of $A$, one for each double coset; and the class integral commutes with finite
sums of additive maps.\par
\textbf{Q2} (names at split places; the catalogue's files on the mixed model). For $\Sigma$ disjoint from its
conjugate: the indicator of the matrices integral at $\Sigma$ names the projection \m{P^{\Sigma}}; the name of its
Schrödinger translate by a vector $t$ is the Heisenberg translate by the part of $t$ above \m{\Sigma\cup\cj\Sigma},
composed with \m{P^{\Sigma}}; and for \m{g}, \m{h} supported on \m{\Sigma\cup\cj{\Sigma}} the pair Weil operator
with characters $a$, $b$ sends the operator named by $\Psi$ to the one named by \m{\Psi(h^{-1}\,\cdot\,g)}, times the
values of the space characters at the determinants at $\Sigma$ and powers of their absolute values. Reference for
the Schrödinger model at split places: to be confirmed.\par
\textbf{Q3} (gauges). The gauge at $g$ does not change when $g$ is multiplied on the right by an element that is
integral at every place above~$\ell$; the gauge of a sum is at most the largest gauge of a term; an integer matrix,
and a Schur matrix of integral matrices, do not increase $\ell$-adic norms; some rational point moves any given point,
and a given conjugate of it, to points integral above~$\ell$ (weak approximation).\par
\textbf{Q4} (dimension of forms with free classes). If all the groups of rational points in the conjugates of a
level are trivial, the space of forms of that level has dimension the number of classes times the dimension of the
weight. The catalogue cites B.~Gross, \emph{Algebraic modular forms}, Israel J. Math. 113 (1999), without a place:
reference to be confirmed. For the columns the groups are trivial (unit \emph{Columns}).\par
\textbf{Q5} (characters). The splitting character $\upsilon$ is unramified at every finite place; a character of
the norm-one adeles with a weight and a tame part that is trivial at a place $w$ and at $\cj{w}$ has an unramified
base change at~$w$; for an index set of even cardinality the space character of the rank twist of $\nu$ is the base
change of~$\nu$.}

\sect{Split levels and the Gauss-sum operator}{s:proofsA}

\begin{lemma}[elementary facts]\label{l:elem}
Let $\Sigma$ be a finite set of finite places of $E$, $N$ a natural number, \m{y\in\RS{\Sigma}}, and let
\m{\mathbf{1}_{\Sigma}\in\Af} be the adele equal to $1$ at the places of $\Sigma$ and $0$ elsewhere.
\begin{enumerate}[label=(\roman*)]
\item If $N$ is prime, \m{\mathbf{1}_{\Sigma}\cdot\rN{y}=\rN{y}}.
\item Let \m{\splitlevel{\Sigma}{N}} with $N$ prime, \m{x\in\OS{\Sigma}} and \m{c\in\RS{\Sigma}}. Then \m{\red(x)=c} if and
  only if \m{(x_w-\tilde c_w)/N\in O_w} for every \m{w\in\Sigma}.
\item If \m{\splitlevel{\Sigma}{N}}, then $\Sigma$ and its conjugate \m{\cj{\Sigma}} are disjoint.
\item If \m{\splitlevel{\Sigma}{N}}, then every \m{w\in\Sigma} divides~$N$.
\item Let \m{\splitlevel{\Sigma}{N}} with $N$ prime, \m{x\in\OS{\Sigma}}, \m{c\in\RS{\Sigma}} and \m{u\in\Af} such that
  $x$ is the family of the components of \m{N u} at $\Sigma$. Then \m{\red(x)=c} if and only if
  \m{\mathbf{1}_{\Sigma}\cdot(u-\rN{c})\in\Oh}.
\item Let \m{\splitlevel{\Sigma}{N}} with $N$ prime, $I$ a finite index set and $\Sigma'$ a finite set of finite
  places. Then \m{\diagN{I}{\Sigma}=\diagN{I}{\Sigma\smallsetminus\Sigma'}\cdot\diagN{I}{\Sigma'}}.
\end{enumerate}
\end{lemma}
\begin{proof}
(i) The adele \m{\rN{y}} vanishes outside~$\Sigma$. (ii) The component of \m{\red(x)} at $w$ is $c_w$ exactly when
\m{x_w-\tilde c_w} lies in the maximal ideal of $O_w$, which is \m{N\,O_w} by (c) of Definition~\ref{d:split}.
(iii) and (iv) are (a) and one direction of (d). (v) is (ii), since \m{(x_w-\tilde c_w)/N} is the component at $w$ of
\m{u-\rN{c}} and membership in $\Oh$ of an adele supported on $\Sigma$ is tested at the places of~$\Sigma$. (vi) Both
sides are diagonal with last entry $N$ at every place of~$\Sigma$.
\end{proof}
\leanline{\lean{idemmulresidueOverN}\sep \lean{reductioneqiff}\sep \lean{disjointstar}\sep \lean{natCastmem}\sep
\lean{reductioneqiffidemmulmem}\sep \lean{diagNeqsdiffmul}\sep An unnamed instance records that $\red$ is a local
homomorphism: an element of $\OS{\Sigma}$ with invertible reduction is a unit.}

\begin{lemma}\label{l:resid}
Let \m{\splitlevel{\Sigma}{N}} with $N$ prime, and let \m{z\in\Af}. There is \m{b\in\RS{\Sigma}} such that for every
\m{y\in\RS{\Sigma}}:
\[\mathbf{1}_{\Sigma}\cdot(z-\rN{y})\in\Oh\quad\Longleftrightarrow\quad\mathbf{1}_{\Sigma}\cdot(Nz)\in\Oh\ \text{ and }\ b=y.\]
\end{lemma}
\begin{proof}
If \m{\mathbf{1}_{\Sigma}(z-\rN{y})} is integral, so is
\m{\mathbf{1}_{\Sigma}Nz=N\cdot\mathbf{1}_{\Sigma}(z-\rN{y})+\mathbf{1}_\Sigma(\tilde y_w)_w}. So if \m{\mathbf{1}_{\Sigma}Nz} is not integral, both sides are false for every $y$ and every $b$ serves. If
\m{\mathbf{1}_{\Sigma}Nz\in\Oh}, let \m{x\in\OS{\Sigma}} be the family of the components of \m{Nz} at
$\Sigma$ and take for $b$ its reduction, \m{b=\red(x)}. By Lemma~\ref{l:elem}\,(v), \m{\mathbf{1}_{\Sigma}(z-\rN{y})\in\Oh} if and only if
\m{\red(x)=y}, that is \m{b=y}.
\end{proof}
\leanline{\lean{existssubresidueOverNmemiff}\sep The text reading of the census lists a character~$b$ of the
surrounding section as an argument; Lean does not take it: the $b$ of the statement is the element it provides.}

\begin{lemma}[functions of residue matrices]\label{l:fnfacts}
Let $\Sigma$ be a split level for the prime~$N$ and $I$, $I'$ finite index sets.
\begin{enumerate}[label=(\roman*)]
\item For every \m{\Ip{I'}\times\Ip{I}} matrix $S$ over $\ES{\Sigma}$, with $\uz$ in ranks \m{\Ip{I'}} and \m{\Ip{I}}:
  \m{\rrow(\uz^{-1}S\,\uz)=\bcol{\ebas}(S)}.
\item Let $f$ be a function of \m{I\times I} matrices over $\RS{\Sigma}$ and \m{g\in\GL_I(\ES{\Sigma})} with
  \m{\nat{f}(g)\ne0} and \m{\nat{f}(g^{-1})\ne0}. Then \m{g=k} for some \m{k\in\GL_I(\OS{\Sigma})} with
  \m{f(\red k)\ne0}.
\item If \m{g\in\Mir{I}}, then \m{\bcol{\ebas}(g)\ne0}.
\item Let $\rho$, $\rho'$ be complex representations of \m{\GL_{\Ip{I}}(\ES{\Sigma})}, \m{\GL_{\Ip{I'}}(\ES{\Sigma})},
  $L$ a linear map from functions of \m{\Ip{I'}\times\Ip{I}} matrices over $\ES{\Sigma}$ to the linear maps
  \m{\rho\to\rho'}, $\varphi$ a vector of~$\rho$, $J$ a finite set, $r$ a map from the
  \m{\Ip{I'}\times\Ip{I}} matrices over $\RS{\Sigma}$ to $J$ and \m{g\colon J\to\CC}. Then
  \m{L[\nat{(g\circ r)}]\,\varphi=\sum_{y\in J}g(y)\cdot L[\nat{(T\mapsto[r(T)=y])}]\,\varphi},
  where \m{[r(T)=y]} is $1$ if \m{r(T)=y} and $0$ otherwise.
\end{enumerate}
\end{lemma}
\begin{proof}
(i) The matrix \m{\uz^{-1}S\uz} has the entries of $S$, except that the last row is multiplied by $N$ and the last
column divided by~$N$ (the corner entry is unchanged). It is integral with last row \m{\equiv\ebas} if and only if
the entries of $S$ outside the last column are integral, \m{S_{\st\st}\equiv1} and \m{S_{i\st}/N} is integral for
\m{i\ne\st}: that is, $S$ is integral with last column \m{\equiv\ebas}, because $N$ is a uniformiser at every place
of~$\Sigma$. (ii) Both $g$ and $g^{-1}$ have entries in $\OS{\Sigma}$, since $\nat{f}$ vanishes on non-integral
matrices. (iii) is the definition of \m{\Mir{I}}. (iv) \m{g\circ r=\sum_{y}g(y)\,[r=y]}, and \m{f\mapsto L[\nat{f}]\varphi}
is linear.
\end{proof}
\leanline{\lean{lastRowDatumconju}\sep \lean{existsinclGLeq}\sep \lean{lastColDatumnezero}\sep
\lean{liftFncompeqsumsmul}\sep In the last one the census lists the variable $y$ of the section, which Lean does not
take.}

\begin{lemma}[the mirabolic group and the Gauss-sum operator]\label{l:mirfix}
Let $E$ have Weil data, let $\Sigma$ be a split level for the prime~$N$, $I$ a finite index set, $\rho$ a complex
representation of \m{\GL_{\Ip{I}}(\ES{\Sigma})}, $\varphi$ a vector of~$\rho$, $T$ a linear endomorphism of the space
of~$\rho$ and \m{g\in\Mir{I}}. Assume that \m{\rho(a)\bigl(T(\rho(b)\varphi)\bigr)=\varphi} for all
\m{a,b\in\GL_{\Ip{I}}(\ES{\Sigma})} with \m{a\,g\,b\in\Mir{I}}. Then \m{T\bigl(\Fop{\rho}\varphi\bigr)=\Fop{\rho}\varphi}.
\end{lemma}
\begin{proof}
As $g$ and $g^{-1}$ lie in \m{\Mir{I}}, Lemma~\ref{l:fnfacts}\,(iii) and (i) give \m{\rrow(\uz^{-1}g\uz)\ne0} and
\m{\rrow(\uz^{-1}g^{-1}\uz)\ne0}; by (ii), \m{\uz^{-1}g\,\uz=k} with \m{k\in\GL_{\Ip{I}}(\OS{\Sigma})} and the last
row of \m{\bar k=\red(k)} equal to~$\ebas$. Multiplication by $\bar k$ permutes the vectors $y$ with
\m{y_{\st}} a unit and preserves \m{y_{\st}}, because \m{(\bar ky)_{\st}} is the product of the last row of $\bar k$
with~$y$. So, with \m{\tilde k(y)} as in Definition~\ref{d:F},
\[
\Fop{\rho}\varphi=\sum_{y}\psiN{y_{\st}}\;\rho\bigl(\uz\,\tilde k(\bar ky)\bigr)\varphi ,
\qquad
T\bigl(\Fop{\rho}\varphi\bigr)=\sum_{y}\psiN{y_{\st}}\;T\bigl(\rho(\uz\,\tilde k(y))\varphi\bigr),
\]
and it suffices to show \m{T(\rho(b)\varphi)=\rho(a)^{-1}\varphi} for \m{b=\uz\tilde k(y)} and
\m{a=\tilde k(\bar ky)^{-1}\uz^{-1}}. This is the hypothesis, once \m{a\,g\,b=\tilde k(\bar ky)^{-1}\,k\,\tilde k(y)} is
seen to lie in \m{\Mir{I}}: it is integral and invertible, and its reduction sends $\ebas$ first to $y$ (the last
column of \m{k(y)}), then to \m{\bar ky}, then back to $\ebas$ (as \m{k(\bar ky)\ebas=\bar ky}).
\end{proof}
\leanline{\lean{mirabolicfixfourierOp}\sep Lean takes $N$, its primality and the split level as arguments here;
the text reading of the census misses them because the statement does not write their names.}

\sect{Lift laws: proof of Theorem A}{s:proofsB}

\begin{lemma}\label{l:lawfacts}
In Setting~\ref{s:split}, let $I$, $I'$ be finite index sets, $\rho$, $\rho'$ complex representations of
\m{\GL_{\Ip{I}}(\ES{\Sigma})}, \m{\GL_{\Ip{I'}}(\ES{\Sigma})}, and let $\law{\rho}{\rho'}{L}{\varkappa}$; let $\varphi$ be a
vector of~$\rho$.
\begin{enumerate}[label=(\roman*)]
\item Let \m{\varphi\in\ram{\rho}}, \m{w\in\Sigma}, and let $f$ be a function of \m{\Ip{I'}\times\Ip{I}} matrices
  over $\RS{\Sigma}$ such that \m{f(c\,T)=f(T)} for every $T$ and every \m{c\in\RS{\Sigma}} with the property: for
  all \m{t,y\in\RS{\Sigma}} with \m{y_w=0}, \m{c\,t=y} if and only if \m{t=y}. Then \m{L[\nat{f}]\,\varphi=0}.
\item For \m{k\in\GL_{\Ip{I'}}(\OS{\Sigma})} and \m{y\in(\RS{\Sigma})^{\Ip{I'}}}:
  \m{\rho'(k)\bigl(L[\bcol{y}]\,\varphi\bigr)=L[\bcol{\red(k)y}]\,\varphi}.
\end{enumerate}
The property of $c$ in (i) holds exactly for the $c$ with \m{c_v=1} for \m{v\ne w} and \m{c_w\ne0}.
\end{lemma}
\begin{proof}
(i) As $\varphi$ is ramified there are a unit $x$ of $\OS{\Sigma}$ with \m{\red(x)_v=1} for \m{v\ne w} and a vector
$\psi$ with \m{\varphi=\psi-\rho(x\cdot1)\psi}. The element \m{c=\red(x)} has the property: for \m{y_w=0} the
equation \m{ct=y} says \m{t_v=y_v} for \m{v\ne w} and \m{c_wt_w=0}, and $c_w$ is a unit. Integral equivariance, with
\m{a=x\cdot1} and \m{b=1}, gives \m{L[\nat{(T\mapsto f(cT))}](\rho(x\cdot1)\psi)=L[\nat{f}]\psi}, and \m{f(cT)=f(T)}.
Hence \m{L[\nat{f}]\varphi=L[\nat{f}]\psi-L[\nat{f}]\psi=0}. (ii) Apply integral equivariance with \m{a=1},
\m{b=k} to the indicator $f$ of the matrices with last column \m{\red(k)y}: then
\m{f(\red(k)T)} is the indicator of the matrices with last column~$y$.
\end{proof}
\leanline{\lean{eqzeroofsmulinvariant}\sep \lean{inclGLlastColDatum}\sep In the first the census lists the section
variable~$y$, which Lean does not take; in the second Lean also takes $N$, its primality and the split level.}

\begin{proposition}[one law and the Gauss-sum operator]\label{p:slide}
In Setting~\ref{s:split}, let $I$, $I'$ be finite index sets, $\rho$, $\rho'$ complex representations of
\m{\GL_{\Ip{I}}(\ES{\Sigma})}, \m{\GL_{\Ip{I'}}(\ES{\Sigma})}, let $\law{\rho}{\rho'}{L}{\varkappa}$ and
\m{\varphi\in\ram{\rho}}. Put \m{C=\rho'(\uz)\bigl(L[\cchar{\ebas}]\,\varphi\bigr)}. Then
\[
C\;=\;\Fop{\rho'}\bigl(L[\bcol{\ebas}]\,\varphi\bigr)\qquad\text{and}\qquad C\;=\;\varkappa\cdot L[\bcol{\ebas}]\bigl(\Fop{\rho}\varphi\bigr).
\]
\end{proposition}
\begin{proof}
For \m{t\in\RS{\Sigma}} let $f_t$ be the indicator of the residue matrices $T$ with corner entry \m{T_{\st\st}=t}; then
\m{\cchar{\ebas}=\sum_t\psiN{t}\,\nat{f_t}} by Lemma~\ref{l:fnfacts}\,(iv). If $t$ is not a unit, say \m{t_w=0}, then
$f_t$ is invariant under the scalars $c$ of Lemma~\ref{l:lawfacts}\,(i), so \m{L[\nat{f_t}]\varphi=0}. Hence
\begin{equation}\label{e:trunc}
L[\cchar{\ebas}]\,\varphi\;=\;L[\nat{h}]\,\varphi,\qquad h(T)=\psiN{T_{\st\st}}\ \text{if }T_{\st\st}\text{ is a unit},\quad h(T)=0\ \text{otherwise.}
\end{equation}
\emph{First identity.} By Lemma~\ref{l:lawfacts}\,(ii) and \m{k(y)\ebas=y},
\[
\Fop{\rho'}\bigl(L[\bcol{\ebas}]\varphi\bigr)=\rho'(\uz)\sum_{y}\psiN{y_{\st}}\,L[\bcol{y}]\,\varphi ,
\]
the sum over the \m{y\in(\RS{\Sigma})^{\Ip{I'}}} with \m{y_{\st}} a unit. A matrix with last column $y$ has corner
entry \m{y_{\st}}, so \m{\sum_y\psiN{y_{\st}}\bcol{y}=\nat{h}}, and the right side is $C$ by~\eqref{e:trunc}.

\emph{Second identity.} By the law of $\uz$,
\m{\varkappa\,L[\bcol{\ebas}](\Fop{\rho}\varphi)=\rho'(\uz)\sum_y\psiN{y_{\st}}\,L[\rrow]\bigl(\rho(\tilde k(y))\varphi\bigr)},
the sum over the \m{y\in(\RS{\Sigma})^{\Ip{I}}} with \m{y_{\st}} a unit. Integral equivariance with
\m{a=\tilde k(y)}, \m{b=1}, applied to the indicator $f$ of the matrices with last row~$y$, gives
\m{L[\rrow](\rho(\tilde k(y))\varphi)=L[\nat{f}]\varphi}: indeed \m{T\,k(y)} has last row $y$ if and only if $T$
has last row \m{y\,k(y)^{-1}=\ebas}, the last row of \m{k(y)} being~$y$. Summing over $y$ with the weights
\m{\psiN{y_{\st}}} gives again \m{\nat{h}}, and the right side is $C$ by~\eqref{e:trunc}.
\end{proof}
\leanline{\lean{ucharDatumeqfourierOp}\sep Lean also takes $N$, its primality and the split level.}

\begin{proof}[Proof of Theorem A]
Write \m{F=\Fop{\rho_0}} and, for \m{q,j\ge0} and a vector $\psi$ of $\rho_0$,
\[B_j^{q}(\psi)\;=\;L_{j-1}[\bcol{\ebas}]\circ\dots\circ L_0[\bcol{\ebas}]\,(F^{q}\psi),\]
a vector of~$\rho_j$.

\emph{(1) The vectors \m{B_j^{q}(\varphi)} are ramified.} Let \m{w\in\Sigma} and
\m{\varphi=\psi-\rho_0(x\cdot1)\psi} with $x$ as in Definition~\ref{d:diag}. The scalar \m{x\cdot1} is central, so
\m{\rho_0(x\cdot1)} commutes with~$F$. Integral equivariance of $L_j$ with \m{a=(x\cdot1)^{-1}}, \m{b=x\cdot1} gives
\m{\rho_{j+1}(x\cdot1)\circ L_j[\nat{f}]=L_j[\nat{f}]\circ\rho_j(x\cdot1)} for every~$f$, because
\m{f(\red(x)T\red(x)^{-1})=f(T)}. By induction on $j$,
\m{B_j^{q}(\varphi)=B_j^{q}(\psi)-\rho_j(x\cdot1)B_j^{q}(\psi)}.

\emph{(2)} \m{\Fop{\rho_j}\bigl(B_j^{q}(\varphi)\bigr)=\bigl(\prod_{l<j}\varkappa_l\bigr)B_j^{q+1}(\varphi)}. For \m{j=0}
this is the definition. For the step, the two identities of Proposition~\ref{p:slide} for the law $L_j$ and the
ramified vector \m{B_j^{q}(\varphi)} give
\m{\Fop{\rho_{j+1}}\bigl(L_j[\bcol{\ebas}]\,B_j^{q}(\varphi)\bigr)=\varkappa_j\,L_j[\bcol{\ebas}]\bigl(\Fop{\rho_j}B_j^{q}(\varphi)\bigr)}.

\emph{(3) One pair.} Let \m{\psi=B_{2q}^{q}(\varphi)}, a ramified vector of~$\rho_{2q}$. In
\m{\sum_yL_{2q+1}[\cchar{y}]\,L_{2q}[\bcol{y}]\,\psi}, a term with $y$ zero at some place $w$ vanishes: then
\m{\bcol{y}} is invariant under the scalars of Lemma~\ref{l:lawfacts}\,(i), since for such $c$ a column \m{c\,t}
equals $y$ if and only if $t$ does. If $y$ is non-zero at every place, then at each place some entry of $y$ is a
unit, so $y$ is the last column of an invertible matrix over each factor \m{\ZZ/N}, and these assemble to
\m{k\in\GL_{\Ip{I_{2q+1}}}(\RS{\Sigma})} with \m{k\,\ebas=y}. With \m{\tilde k=\secN(k)}, Lemma~\ref{l:lawfacts}\,(ii)
gives \m{L_{2q}[\bcol{y}]\psi=\rho_{2q+1}(\tilde k)\,L_{2q}[\bcol{\ebas}]\psi}, and integral equivariance of
\m{L_{2q+1}} with \m{a=\tilde k}, \m{b=1} gives \m{L_{2q+1}[\cchar{y}]\circ\rho_{2q+1}(\tilde k)=L_{2q+1}[\cchar{\ebas}]},
because the product of a row $r$ with \m{k\ebas=y} is the product of \m{rk} with~$\ebas$. So every such term equals
the term of \m{y=\ebas}, and
\begin{align*}
A_q\,\psi\;&=\;d_q\,o_q\cdot\rho_{2q+2}(\uz)\,L_{2q+1}[\cchar{\ebas}]\bigl(B_{2q+1}^{q}(\varphi)\bigr)
\;=\;d_q\,o_q\cdot\Fop{\rho_{2q+2}}\bigl(B_{2q+2}^{q}(\varphi)\bigr)\\
&=\;d_q\,o_q\prod_{l<2q+2}\varkappa_l\cdot B_{2q+2}^{q+1}(\varphi),
\end{align*}
by the first identity of Proposition~\ref{p:slide} for \m{L_{2q+1}} and the ramified vector
\m{B_{2q+1}^{q}(\varphi)}, and by~(2).

\emph{(4)} The theorem follows by induction on~$p$: it is trivial for \m{p=0}, and if
\m{A_{q-1}\circ\dots\circ A_0(\varphi)} is the stated multiple of \m{B_{2q}^{q}(\varphi)}, then (3) gives the
statement for \m{q+1}.
\end{proof}
\leanline{\lean{compUpTopairseqsmul}\sep The census lists the section variable~$y$, which Lean does not take.}

\sect{Chains: proof of Theorem C}{s:proofsC}

\begin{lemma}[equivariance and linearity of chains]\label{l:cheq}
In Setting~\ref{s:chain}, let \m{\mathcal{K}=(\mathcal{K}_j)_j} be the kernels.
\begin{enumerate}[label=(\roman*)]
\item Let \m{\varphi\in\Alg{I_0}{e_0}} be fixed by $K_0$, let \m{a=(a_j)_j} with \m{a_j\in K_j} for all~$j$,
  \m{\Phi=(\Phi_j)_j} data and \m{j\ge0}. Then
  \m{\chain{\mathcal{K}}{(\omega_i(a_i,a_{i+1})\Phi_i)_i}_{j}(\varphi)=a_j\cdot\chain{\mathcal{K}}{\Phi}_{j}(\varphi)}.
\item Let \m{a=(a_j)_j} be any family of points \m{a_j\in\Gf{I_j}}, \m{t=(t_i)_i} with
  \m{t_i\in\mathcal{S}_{2i}\otimes\mathcal{S}_{2i+1}}, \m{p\ge0} and \m{\varphi\in\Alg{I_0}{e_0}}. Then
  \m{\chaintwo{\mathcal{K}}{(\twoact{\mathcal{K}}{a}{i}\,t_i)_i}_{p}(a_0\cdot\varphi)=a_{2p}\cdot\chaintwo{\mathcal{K}}{t}_{p}(\varphi)}.
\item For \m{j\ge0}, \m{\varphi\in\Alg{I_0}{e_0}} and \m{g\in\Gf{I_j}} there is a linear map $\beta$ on
  \m{\mathcal{S}_{j-1}\otimes\dots\otimes\mathcal{S}_0} with
  \m{\chain{\mathcal{K}}{\Phi}_{j}(\varphi)(g)=\beta(\Phi_{j-1}\otimes\dots\otimes\Phi_0)} for all data~$\Phi$.
\end{enumerate}
\end{lemma}
\begin{proof}
(i) and (ii) follow by induction on the number of lifts from the equivariance of one lift
(Definition~\ref{d:kernel}): in (i) the form $\varphi$ is fixed by \m{a_0\in K_0}; in (ii) one first treats pure
tensors \m{t_p=s\otimes s'}, for which the pair is \m{\mathcal{K}_{2p+1}[s']\circ\mathcal{K}_{2p}[s]}, and extends by
linearity. (iii) A lift is linear in its datum, so the value of the composite at $\varphi$ and $g$ is multilinear in
\m{(\Phi_{j-1},\dots,\Phi_0)}.
\end{proof}
\leanline{\lean{chainequivariant}\sep \lean{chainequivarianta}\sep \lean{existslinearchaineqpure}\sep In the last
one the census lists the section variable $\Phi$, which Lean does not take.}

\begin{proof}[Proof of Theorem C]
By induction on $m$, for all $M$, $L$, $h$, $A$ and~$P$. Let \m{\beta_j(g)} be the linear map of
Lemma~\ref{l:cheq}\,(iii) for $\varphi$ and~$g$, and write \m{\Phi_x^{<j}=\Phi_{x,j-1}\otimes\dots\otimes\Phi_{x,0}}.

\emph{\m{m=0}.} The chain is the identity and \m{\varphi(h)=\beta_0(h)(1)}. By (e), \m{P_0(h,\varphi(h))}, so (f)
with \m{\beta=\beta_0(h)} gives \m{\sum_xL_x(\varphi(h))\in A}; the empty product of class numbers is~$1$.

\emph{From $j$ to \m{m=j+1}.} By the definition of the lift,
\begin{align*}
\sum_{x\in\mathcal{X}_0}L_x\bigl(\chain{\mathcal{K}}{\Phi_x}_{j+1}(\varphi)(h)\bigr)&=\int_{\Gamma\backslash\Gf{I_j}}G(g)\,dg,\\
G(g)&=\sum_{x\in\mathcal{X}_0}L_x\Bigl(\mathcal{J}_j\bigl(\omega_j(g,h)\Phi_{x,j}\bigr)\bigl(\chain{\mathcal{K}}{\Phi_x}_{j}(\varphi)(g)\bigr)\Bigr),
\end{align*}
where Q1 allows the finite sum to pass under the class integral. By Q1 and (b) it suffices to show that $G$ is
right $K_j$-invariant and that every \m{G(g)} is a sum of \m{\prod_{i<j}\cls{K_i}} elements of~$A$.

\emph{Values.} Fix $g$ and apply the induction hypothesis to the same $\Phi$ and $\varphi$, the point $g$, the set
$A$ and the linear maps \m{L'_x(u)=L_x\bigl(\mathcal{J}_j(\omega_j(g,h)\Phi_{x,j})u\bigr)} on \m{\Wk{I_j}}. For a
linear $\beta'$ on \m{\mathcal{S}_{j-1}\otimes\dots\otimes\mathcal{S}_0} let \m{T_{\beta'}} be the linear map on
\m{\mathcal{S}_j\otimes\dots\otimes\mathcal{S}_0} with
\m{T_{\beta'}(s\otimes t)=\mathcal{J}_j(\omega_j(g,h)s)\,\beta'(t)}, so that
\m{L'_x(\beta'(\Phi_x^{<j}))=L_x(T_{\beta'}(\Phi_x^{<j+1}))}. Joint invariance of length $j$ for the $L'_x$ is
hypothesis (a) for the families $a$ with \m{a_i=1} for \m{i\ge j} and the maps \m{T_{\beta'}}. Hypothesis (f) for
the $L'_x$ and the point $g$: if \m{P_j(g,\beta'(\Phi_x^{<j}))} for all $x$, then (d) gives
\m{P_{j+1}(h,T_{\beta'}(\Phi_x^{<j+1}))} for all $x$, and (f) for $\beta=T_{\beta'}$ gives
\m{\sum_xL'_x(\beta'(\Phi_x^{<j}))\in A}. The other hypotheses are inherited. So \m{G(g)=\sum_xL'_x(\chain{\mathcal{K}}{\Phi_x}_{j}(\varphi)(g))} is a sum of
\m{\prod_{i<j}\cls{K_i}} elements of~$A$.

\emph{Invariance.} Let \m{a\in K_j} and let $a'$ be the family with $a$ at the index $j$ and $1$ elsewhere. By
Lemma~\ref{l:cheq}\,(i),
\begin{align*}
\chain{\mathcal{K}}{\Phi_x}_{j}(\varphi)(ga)&=\beta_j(g)\Bigl(\textstyle\bigotimes_{i<j}\omega_i(a'_i,a'_{i+1})\Phi_{x,i}\Bigr),\\
\omega_j(ga,h)\Phi_{x,j}&=\omega_j(g,h)\,\omega_j(a'_j,a'_{j+1})\Phi_{x,j}.
\end{align*}
Hence
\[G(ga)=\sum_{x\in\mathcal{X}_0}L_x\Bigl(T_{\beta_j(g)}\bigl(\textstyle\bigotimes_{i\le j}\omega_i(a'_i,a'_{i+1})\Phi_{x,i}\bigr)\Bigr),\]
which by (a) for the family $a'$ equals \m{\sum_xL_x(T_{\beta_j(g)}(\Phi_x^{<j+1}))=G(g)}.
\end{proof}
\leanline{\lean{sumchainmemnsmul}}

\sect{Names at a split level: proof of Theorem B}{s:proofsD}

\begin{lemma}[the embedding at a split level]\label{l:embfacts}
Let \m{\splitlevel{\Sigma}{N}}, $I$ a finite index set, and for \m{g\in\Gf{I}} let \m{g_\Sigma\in\GL_I(\ES{\Sigma})}
be the family of its components at~$\Sigma$.
\begin{enumerate}[label=(\roman*)]
\item For \m{g\in\GL_I(\ES{\Sigma})}: \m{\emb{\Sigma}{g}_{\Sigma}=g}, and \m{\emb{\Sigma}{g}} has component $1$ at
  every place outside \m{\Sigma\cup\cj{\Sigma}}.
\item Let \m{g\in\Gf{I}} and \m{a\in\GL_I(\OS{\Sigma})} with \m{a=g_\Sigma}. Then the component of $g$ at every
  place dividing $N$ is integral and invertible.
\item Let $N$ be prime, \m{w\in\Sigma}, and let $\nu$ be a finite adele with \m{\nu\cj{\nu}=1} that is integral and
  has \m{\nu_v=1} at every place $v$ with \m{v\ne w} and \m{\cj{v}\ne w}. Then there is a unit $x$ of $\OS{\Sigma}$
  with \m{\red(x)_v=1} for all \m{v\in\Sigma}, \m{v\ne w}, such that \m{\emb{\Sigma}{x\cdot1}} is the scalar matrix
  of~$\nu$ in \m{\Gf{I}}.
\end{enumerate}
\end{lemma}
\begin{proof}
(i) holds by construction. (ii) A place $w$ dividing $N$ lies in $\Sigma$ or in \m{\cj\Sigma}, by (d) of
Definition~\ref{d:split}. In the first case \m{g_w=a_w\in\GL_I(O_w)}; in the second $g_w$ is the inverse of the
conjugate transpose of the component at~\m{\cj{w}}, which is integral and invertible. (iii) Let $x$ be the family of
the components of $\nu$ at~$\Sigma$; it is integral, and so is the family of the components of \m{\cj\nu}, which is
its inverse because \m{\nu\cj\nu=1}. For \m{v\in\Sigma}, \m{v\ne w}, also \m{\cj{v}\ne w} by (a), so \m{\nu_v=1}. The
two elements \m{\emb{\Sigma}{x\cdot1}} and the scalar matrix of $\nu$ are unitary, have component $1$ outside
\m{\Sigma\cup\cj\Sigma} and the same components at~$\Sigma$; so they are equal.
\end{proof}
\leanline{\lean{atPlacesembedGL}\sep \lean{embedGLmemsupportedOn}\sep \lean{memintegralAboveofinclGLeq}\sep
\lean{existsembedGLscalareqscalarFin}\sep \lean{scalarFinatPlace} (the component at a place $w$ of the scalar matrix
of $\nu$ is the scalar matrix of~$\nu_w$; a rewriting rule, used silently)}

\begin{lemma}[Weil operators at a split level]\label{l:names}
Let $E$ have Weil data, $I$, $I'$ finite index sets, $a$, $b$ characters of the norm-one adeles of $E$ trivial on the
rational points, and let \m{\omega_{a,b}} be the pair Weil representation of \m{\Gf{I}\times\Gf{I'}} on
\m{\mathcal{W}_f(\operatorname{Hom}(X_I,X_{I'}))} with characters $a$,~$b$.
\begin{enumerate}[label=(\roman*)]
\item For finite adeles $\nu$, $\nu'$ of norm one, with scalar matrices \m{s(\nu)\in\Gf{I}}, \m{s(\nu')\in\Gf{I'}}:
  \m{\omega_{a,b}(s(\nu),s(\nu'))=a(\det s(\nu'))\,b(\det s(\nu))^{-1}\cdot\omega_{1,1}(1,s(\nu'/\nu))}.
\item Let $\Sigma$ be a split level for the prime~$N$, \m{g\in\GL_I(\ES{\Sigma})}, \m{k\in\GL_{I'}(\ES{\Sigma})}, $f$ a
  function of \m{I'\times I} matrices over $\RS{\Sigma}$ and \m{v\in\mathcal{W}_f(\operatorname{Hom}(X_I,X_{I'}))}. Then
  \[
  \omega_{a,b}\bigl(\emb{\Sigma}{g},\emb{\Sigma}{k}\bigr)\,\nm{\Sigma}{\nat{f}}\,v=
  \detchar{\eta_b}{|I'|}(g)^{-1}\,\detchar{\eta_a}{|I|}(k)\cdot\nm{\Sigma}{S\mapsto\nat{f}(k^{-1}Sg)}\,v,
  \]
  where \m{\eta_a(x)=a(x/\cj{x})\,\upsilon(x)^{-|I|}} and \m{\eta_b(x)=b(x/\cj{x})\,\upsilon(x)^{-|I'|}} are the space
  characters.
\item In the situation of (ii), the function \m{t\mapsto\nat{f}(t_\Sigma)} of finite adelic matrices is the function
  that takes the value of $f$ at the family of the reductions of $t$ modulo the places of $\Sigma$ if $t$ is
  integral there, and $0$ otherwise; it is a Schwartz--Bruhat function at~$\Sigma$, so that it has a name.
\end{enumerate}
\end{lemma}
\begin{proof}
(i) A pair of scalars acts on \m{\operatorname{Hom}(X_I,X_{I'})} through the scalar \m{\nu'/\nu}; the characters
contribute \m{a(\det s(\nu'))\,b(\det s(\nu))^{-1}} by the definition of the twisted representation. (ii) is Q2:
\m{\emb{\Sigma}{g}} and \m{\emb{\Sigma}{k}} are supported on \m{\Sigma\cup\cj\Sigma}, with components $g$, $k$
at~$\Sigma$ (Lemma~\ref{l:embfacts}\,(i)), and $\Sigma$ is disjoint from its conjugate; the scalar of Q2 is the
stated product of the two determinant characters. (iii) is the definition of $\nat{f}$, read through the
isomorphisms \m{\mathbb{F}_w\cong\ZZ/N}.
\end{proof}
\leanline{\lean{pairWeilFinitescalarFin}\sep \lean{pairWeilFinitename}\sep \lean{liftFneqresidueDatum}\sep
\lean{isSchwartzBruhatliftFn} (an instance)}

\begin{lemma}[determinant characters]\label{l:detchar}
Let $I$, $I'$ be finite index sets and $\Sigma$ a finite set of finite places of~$E$.
\begin{enumerate}[label=(\roman*)]
\item Let $E$ have Weil data and let $a$ be a character of the norm-one adeles of $E$, trivial on the rational
  points, whose base change \m{x\mapsto a(x/\cj{x})} is unramified at every \m{w\in\Sigma}. Then for all
  \m{i,k\in\NN} and \m{g\in\GL_I(\OS{\Sigma})}, with \m{\eta(x)=a(x/\cj{x})\upsilon(x)^{-i}}: \m{\detchar{\eta}{k}(g)=1}.
\item Let \m{\splitlevel{\Sigma}{N}} with $N$ prime, $\Sigma'$ a finite set of finite places,
  \m{\jm\colon\CC\to\CC_N} a ring homomorphism, $\eta$ a Hecke character and \m{k\le k'} natural numbers. Then
  \m{\lvert\jm(\detchar{\eta}{k'}(\diagN{I}{\Sigma'}))\rvert\le\lvert\jm(\detchar{\eta}{k}(\diagN{I'}{\Sigma'}))\rvert}.
\end{enumerate}
\end{lemma}
\begin{proof}
(i) \m{\det g} is a unit of $\OS{\Sigma}$, so \m{\lvert\det g_w\rvert_w=1}, and \m{\eta_w(\det g_w)=1} because both
the base change of $a$ and $\upsilon$ are unramified at~$w$ (Q5). (ii) The determinant of \m{\diagN{I}{\Sigma'}}
does not depend on the rank: its component is $N$ at the places of \m{\Sigma\cap\Sigma'} and $1$ elsewhere. So the two
values differ by the factor \m{\prod_{w\in\Sigma\cap\Sigma'}\lvert N\rvert_w^{-(k'-k)/2}=N^{(k'-k)\,\#(\Sigma\cap\Sigma')/2}},
because \m{\lvert N\rvert_w=N^{-1}} for a place of norm $N$ with uniformiser~$N$. The square of this factor is a
non-negative power of the integer $N$, of $N$-adic absolute value at most~$1$.
\end{proof}
\leanline{\lean{detCharAtinclGL}\sep \lean{normdetCharAtdiagNanti}}

\begin{proof}[Proof of Theorem B]
Write \m{\mathcal{H}=\operatorname{Hom}(X_{\Ip{I}},X_{I'})} and \m{\mathcal{H}'=\operatorname{Hom}(X_{I'},X_{\Ip{I''}})},
and let \m{\epsilon=\mathbf{1}_{\Sigma}}, \m{\epsilon'=\mathbf{1}_{\Sigma'}}. Both $\Sigma$ and $\Sigma'$ are split
levels for~$N$; hence \m{\Sigma\cup\cj\Sigma=\Sigma'\cup\cj{\Sigma'}} is the set of the places dividing~$N$, and
\m{\epsilon+\cj{\epsilon}=\epsilon'+\cj{\epsilon'}}, \m{\epsilon'\cj{\epsilon'}=0}.

\emph{Step 1: both sides are names at $\Sigma$.} The projection \m{P^{\Sigma'}} depends only on
\m{\Sigma'\cup\cj{\Sigma'}}, so it equals \m{P^{\Sigma}}, which is named at $\Sigma$ by the indicator
\m{\mathbf{1}_{\mathrm{int}}} of the matrices integral at~$\Sigma$ (Q2). The vectors \m{c(y)}, \m{r(y)} of
Definition~\ref{d:names} are supported above \m{\Sigma'\cup\cj{\Sigma'}}; by Q2 the operator
\m{\varrho_f(c(y))P^{\Sigma}} is named by the Schrödinger translate \m{c(y)\star\mathbf{1}_{\mathrm{int}}}, where
\[
(t\star\Psi)(X)\;=\;\psi_0\bigl(\tfrac12\textstyle\sum\cj{t}\cdot\epsilon\,(2X+t)\bigr)\,\Psi(X+t),
\]
the sum over the entries; likewise for \m{r(y)}. So the left side of the theorem is the image of
\m{Z=\sum_y\,(c(y)\star\mathbf{1}_{\mathrm{int}})\otimes(r(y)\star\mathbf{1}_{\mathrm{int}})} under the linear map
\m{\Psi\otimes\Psi'\mapsto\nm{\Sigma}{\Psi}v\otimes\nm{\Sigma}{\Psi'}v'}, and the right side is the image of
\m{Z'=\sum_z\bigl(S\mapsto\bcol{z}(S\diagN{I}{\Sigma'})\bigr)\otimes\bigl(S'\mapsto\cchar{z}(\diagN{I''}{\Sigma\smallsetminus\Sigma'}S')\bigr)}.
(The formal text first embeds the tensor product of the finite parts into that of the ambient spaces, which is
injective over a field.) A tensor of two functions is determined by its values at pairs of points, so it suffices
to prove \m{Z(T,R)=Z'(T,R)} for all finite adelic matrices $T$ of size \m{I'\times\Ip{I}} and $R$ of size
\m{\Ip{I''}\times I'}.

\emph{Step 2: the left side at \m{(T,R)}.} Put \m{t_i=T_{i\st}} and \m{s_i=R_{\st i}} for \m{i\in I'}, and
\m{x_i=\rN{y_i}}, an adele supported on~$\Sigma'$; note \m{\cj{x_i}\,x_i=0}. Then
\[
Z(T,R)=\sum_{y}\;\psi_0\Bigl(-\sum_{i}\bigl(\cj{x_i}\,\epsilon\,t_i+x_i\,\epsilon\,s_i\bigr)\Bigr)\cdot
\bigl[\epsilon\,(T+c(y))\text{ and }\epsilon\,(R+r(y))\text{ integral}\bigr].
\]
The integrality conditions concern the entries of $T$ outside the last column and of $R$ outside the last row,
which do not involve~$y$, and the entries \m{\epsilon(t_i-x_i)} and \m{\epsilon(s_i-\cj{x_i})}. Let
\m{\mathsf{p}=1+\epsilon'(N-1)}, the adele equal to $N$ on $\Sigma'$ and $1$ elsewhere, so that \m{N+1-\mathsf{p}} is $1$ on
$\Sigma'$ and $N$ elsewhere. Applying Lemma~\ref{l:resid} for the split level $\Sigma'$ to the adele
\m{\epsilon t_i+\cj{\epsilon s_i}}, and using that an adele $u$ has \m{\epsilon u} and \m{\epsilon\cj{u}} integral if
and only if \m{\epsilon'u} and \m{\epsilon'\cj{u}} are, one finds \m{b_i\in\RS{\Sigma'}} such that for every
\m{y_i}:
\[
\epsilon(t_i-x_i)\ \text{and}\ \epsilon(s_i-\cj{x_i})\ \text{integral}\iff
\epsilon\,t_i\mathsf{p}\ \text{and}\ \epsilon\,(N+1-\mathsf{p})s_i\ \text{integral, and}\ y_i=b_i .
\]
So every term with \m{y\ne b} vanishes. The term \m{y=b} can vanish as well: it is non-zero exactly when
\m{T\cdot\operatorname{diag}(1,\dots,1,\mathsf{p})} and \m{\operatorname{diag}(1,\dots,1,N+1-\mathsf{p})\cdot R} are
integral at~$\Sigma$ (the formal proof records this by a case split after discarding the other terms).

\emph{Step 3: the right side at \m{(T,R)}.} At the places of $\Sigma$ the matrix \m{\diagN{I}{\Sigma'}} is
\m{\operatorname{diag}(1,\dots,1,\mathsf{p})} and \m{\diagN{I''}{\Sigma\smallsetminus\Sigma'}} is
\m{\operatorname{diag}(1,\dots,1,N+1-\mathsf{p})}. So \m{Z'(T,R)} vanishes unless the same two matrices are integral
at~$\Sigma$; then exactly one $z$ contributes, the reduction of the last column \m{(t_i\mathsf{p})_i}, and
\m{Z'(T,R)=\psiN{\zeta}} with \m{\zeta\in\RS{\Sigma}} the reduction of
\m{\sum_i(N+1-\mathsf{p})s_i\cdot t_i\mathsf{p}=N\sum_is_it_i} at~$\Sigma$ (note
\m{(N+1-\mathsf{p})\mathsf{p}=N}).

\emph{Step 4: the two characters agree.} It remains to see that, for \m{y=b},
\[\psi_0\Bigl(-\sum_i(\cj{x_i}\,\epsilon\,t_i+x_i\,\epsilon\,s_i)\Bigr)\;=\;\psi_0\bigl(-\rN{\zeta}\bigr).\] The character $\psi_0$ is
trivial on $\Oh$, so it suffices that the two arguments differ by an integral adele. By Step~2 the adeles \m{\epsilon(t_i-x_i)} and
\m{\epsilon(s_i-\cj{x_i})} are integral, and by Lemma~\ref{l:elem}\,(v), applied to \m{u=\sum_is_it_i},
\m{\epsilon\,(\sum_is_it_i-\rN{\zeta})} is integral. Expanding the product
\m{\epsilon\sum_i(t_i-x_i)(s_i-\cj{x_i})}, which is integral, and using \m{x_i\cj{x_i}=0} gives the claim. This is the
computation that the formal proof closes with its last three lines.
\end{proof}
\leanline{\lean{translateSumeqsumname}\sep The formal statement takes the split level structure of $\Sigma'$ as an
explicit hypothesis and that of $\Sigma$ as an instance.}

\sect{The level prime of the data: proof of Theorem D}{s:proofsE}

\begin{lemma}[the non-induced places]\label{l:nonind}
In Setting~\ref{s:data}, with \m{\ell=\ND} and \m{\jm\colon\CC\to\CC_{\ND}}:
\begin{enumerate}[label=(\roman*)]
\item \m{\nonind{D}} is a split level for \m{\ND}; in symbols, \m{\splitlevel{\nonind{D}}{\ND}}.
\item Let \m{\nonind{D}} carry this structure, let \m{j\ge0}, $\Sigma'$ a finite set of finite places of \m{E(D)}
  and $\tau$ an infinite place of \m{E(D)}. Then \m{\evt{\tau}}, applied to the entries of
  \m{\emb{\nonind{D}}{\diagN{\XIo{j}}{\Sigma'}}}, gives \m{\operatorname{diag}(1,\dots,1,\ND)} if
  \m{\cj{\ind{\tau}}\in\Sigma'} and the identity matrix otherwise.
\end{enumerate}
\end{lemma}
\begin{proof}
(i) By Input~5, \m{\ND} splits completely in \m{E(D)}: every place above \m{\ND} has norm \m{\ND} and uniformiser
\m{\ND}, the group of infinite places, which is the Galois group of \m{E(D)} over $K$, permutes the places above
\m{\ND} that lie over a fixed place of $K$ simply transitively, and no place above \m{\ND} is fixed by complex
conjugation. The induced places \m{\ind{\tau}} form one such orbit; so no induced place is the conjugate of another,
and every place above \m{\ND} is an induced place or the conjugate of one. These are (a)--(d) of
Definition~\ref{d:split} for the set of the conjugates. (ii) \m{\evt{\tau}} reads the component at
\m{\cj{\ind{\tau}}\in\nonind{D}}, where the embedding places the component of the diagonal matrix
(Lemma~\ref{l:embfacts}\,(i)); that component has last entry \m{\ND} if \m{\cj{\ind{\tau}}\in\Sigma'} and $1$
otherwise, and \m{\evt{\tau}} fixes rational integers.
\end{proof}
\leanline{\lean{isSplitLevelnonInducedPlaces}\sep \lean{adeleExtensionmapdiagN}}

\begin{lemma}[the column level is a conjugated mirabolic group]\label{l:colmir}
In Setting~\ref{s:data}, with \m{\ell=\ND}, let \m{\Sigma=\nonind{D}} with its structure of a split level for
\m{\ND}, let \m{j\ge0}, and for \m{k\in\Gf{\XI{j}}} put \m{k^{\Delta}=\bigl(\dpt{D}{j}^{-1}\,k\,\dpt{D}{j}\bigr)_{\Sigma}},
the family of the components at~$\Sigma$.
\begin{enumerate}[label=(\roman*)]
\item For \m{k\in\Gf{\XI{j}}}: \m{\bcol{\ebas}(k^{\Delta})\ne0} if and only if for every \m{w\in\Sigma}: $k_w$ is
  integral; if \m{\cj{w}\in\PiD}, every entry of the last row of \m{(k^{-1})_{\cj{w}}-1} is divisible by a
  uniformiser; and if \m{w\in\PiD}, every entry of the last row of \m{k_w-1} is divisible by a uniformiser.
\item For \m{k\in\Kcol{j}}: \m{k^{\Delta}\in\Mir{\XIo{j}}}.
\end{enumerate}
\end{lemma}
\begin{proof}
(i) Let \m{w\in\Sigma}. If \m{w\in\PiD}, the component of \m{\dpt{D}{j}} at $w$ is
\m{\delta=\operatorname{diag}(1,\dots,1,\ND)}, and \m{\delta^{-1}k_w\delta} has the entries of $k_w$ with the last
row divided by \m{\ND} and the last column multiplied by \m{\ND}. It is integral with last column congruent to the
last basis vector exactly when $k_w$ is integral and the last row of \m{k_w-1} is divisible by~\m{\ND}, a
uniformiser. If \m{w\notin\PiD}, then \m{\cj{w}\in\PiD} by Input~1 (as $w$ divides \m{\ND}), the component of
\m{\dpt{D}{j}} at $w$ is $1$, and the condition is that $k_w$ is integral with last column congruent to the last
basis vector. Since $k$ is unitary, $k_w$ is the conjugate transpose of \m{(k^{-1})_{\cj{w}}}, so the last column of
\m{k_w-1} is the conjugate of the last row of \m{(k^{-1})_{\cj{w}}-1}. (ii) By Input~2, $k$ and $k^{-1}$ are
integral with last rows congruent to \m{(0,\dots,0,1)} at the places of~\m{\PiD}; so (i) gives
\m{\bcol{\ebas}(k^{\Delta})\ne0} and \m{\bcol{\ebas}((k^{\Delta})^{-1})\ne0}, and Lemma~\ref{l:fnfacts}\,(ii) shows
that \m{k^{\Delta}} is integral and invertible with last column congruent to~$\ebas$.
\end{proof}
\leanline{\lean{lastColDatumconjnezeroiff}\sep \lean{atPlacesconjmemmirabolic}}

\begin{lemma}[ramified columns; the operator preserves columns]\label{l:ramcol}
In Setting~\ref{s:data}, with \m{\ell=\ND} and \m{\nonind{D}} with its structure of a split level for \m{\ND}:
\begin{enumerate}[label=(\roman*)]
\item if \m{\varphi\in\Colram}, then \m{\varphi\in\ram{\nrep{D}{0}}};
\item if \m{\varphi\in\Col}, then \m{\Fop{\nrep{D}{0}}\varphi\in\Col}.
\end{enumerate}
\end{lemma}
\begin{proof}
Write \m{\Sigma=\nonind{D}} and \m{\rho=\nrep{D}{0}}.
(i) Let \m{w\in\Sigma}. It divides \m{\ND}, so $w$ or $\cj{w}$ is a place $\mathfrak{p}$ of \m{\PiD} (Input~1). By
Input~3, \m{\varphi=\psi-z_{\mathfrak{p}}\psi} for a column~$\psi$, and \m{\nu_{\mathfrak{p}}} is integral, of norm
one and equal to $1$ away from $w$ and~$\cj{w}$. By Lemma~\ref{l:embfacts}\,(iii) there is a unit $x$ of
$\OS{\Sigma}$, congruent to $1$ at the places \m{v\ne w} of $\Sigma$, with \m{\emb{\Sigma}{x\cdot1}} the scalar
matrix of~\m{\nu_{\mathfrak{p}}}. Scalar matrices are central, so \m{\cemb{D}{0}(x\cdot1)} is this scalar matrix
too, \m{\rho(x\cdot1)=z_{\mathfrak{p}}}, and $\varphi$ lies in the image of \m{1-\rho(x\cdot1)}.

(ii) Let \m{k\in\Kcol{0}} and let $T$ be right translation by~$k$; we show \m{T\,\Fop{\rho}\varphi=\Fop{\rho}\varphi}
by Lemma~\ref{l:mirfix} with \m{g=k^{\Delta}\in\Mir{\XIo{0}}} (Lemma~\ref{l:colmir}\,(ii)). Let
\m{a,b\in\GL_{\XI{0}}(\ES{\Sigma})} with \m{a\,g\,b\in\Mir{\XIo{0}}}. The operator \m{\rho(a)\,T\,\rho(b)} is right
translation by $\cemb{D}{0}(a)\,k\,\cemb{D}{0}(b)=u\,k$ with
\[u\;=\;\cemb{D}{0}(a)\,k\,\cemb{D}{0}(b)\,k^{-1}.\] The element $u$ has component $1$ outside
\m{\Sigma\cup\cj{\Sigma}}, and \m{u^{\Delta}=(a\,g\,b)\,g^{-1}\in\Mir{\XIo{0}}}. By Lemma~\ref{l:fnfacts}\,(iii) and
Lemma~\ref{l:colmir}\,(i), applied to $u$ and to $u^{-1}$, the element $u$ is integral at the places above \m{\ND}
with last rows congruent to \m{(0,\dots,0,1)} at the places of~\m{\PiD}. It is supported at places above \m{\ND},
none of which is a chosen place (Input~5), so it lies in the chosen level, and \m{u\in\Kcol{0}} by Input~2. Hence
\m{u\,k\in\Kcol{0}} fixes the column~$\varphi$, which is the hypothesis of Lemma~\ref{l:mirfix}.
\end{proof}
\leanline{\lean{memramifiedofmemramifiedColumns}\sep \lean{fourierOpmemcolumn}}

\begin{lemma}[restriction to \m{Y_m(D)} keeps gauges and sizes]\label{l:restr}
In Setting~\ref{s:data}, but for any prime $\ell$ with a ring homomorphism \m{\jm\colon\CC\to\CC_\ell}:
\begin{enumerate}[label=(\roman*)]
\item Let \m{\mathsf{d}\in\Gf{\XI{m}}} be such that for every infinite place $\tau$ the matrix
  \m{\evt{\tau}(\mathsf{d})} is \m{\operatorname{diag}(1,\dots,1,\ND)} if \m{\cj{\ind{\tau}}\in\PiD} and the identity
  otherwise. Let $C$ be a real number and \m{\varphi\in\Stg{m}} of gauge at most $C$ at~$\mathsf{d}$. Then
  \m{\resY(\varphi)} has gauge at most $C$ at~$1$.
\item If \m{\ell\ne\ND}, $C$ is real and \m{\varphi\in\Stg{m}} has gauge at most $C$ at~$1$, then \m{\resY(\varphi)}
  has gauge at most $C$ at~$1$.
\item For every ring homomorphism \m{\sigma\colon\CC\to\CC}, every \m{\varphi\in\Stg{m}} and every point $h$ of the
  finite adelic unitary group of \m{Y_m(D)}:
  \[\wnorm{\sigma\circ\resY(\varphi)(h)}\;\le\;\exp\Bigl(x^{70}-\sum_{\tau}\brE{\tau}\cdot\log\ND\Bigr)\cdot\wnorm{\sigma\circ\varphi(\bdU(h))}.\]
\end{enumerate}
\end{lemma}
\begin{proof}
(i) Fix $h$ and put \m{g=\bdU(h)}. By Definition~\ref{d:columns},
\m{\resY(\varphi)(h)=s\cdot B_D\,\varphi(g)} with \m{s=\prod_\tau\tau(\cj{\piN})^{-\brE{\tau}}}. By Input~4, $B_D$
carries the action of
\[\evt{\tau}(g\mathsf{d})^{-1}\;=\;\evt{\tau}(\mathsf{d})^{-1}\,\evt{\tau}(g)^{-1}\]
in the weight of stage $m$ to the action of \m{\evt{\tau}(h)^{-1}} in the weight of \m{Y_m(D)}, times the scalar
\m{\prod_\tau\lambda_\tau^{-\brE{\tau}}}, where \m{\lambda_\tau\in\{\ND,1\}} is the last entry of
\m{\evt{\tau}(\mathsf{d})}. On the other hand \m{\lvert\jm(\tau(\cj{\piN}))\rvert=\lvert\lambda_\tau\rvert}: the
left side is the absolute value of \m{\piN} at the place \m{\cj{\ind{\tau}}}, which is that of \m{\ND} if this
place divides \m{\piN} and $1$ otherwise. So the two scalars have the same $\ell$-adic size, and the gauge of
\m{\resY(\varphi)(h)} at $h$ is the norm of $B_D$ applied to the vector whose norm is the gauge of \m{\varphi(g)} at
\m{g\mathsf{d}}; the integer matrix $B_D$ does not increase norms (Q3). (ii) is (i) for \m{\mathsf{d}=1}: for
\m{\ell\ne\ND} no place above $\ell$ divides~\m{\piN}. (iii) The scalar $s$ contributes
\m{\lvert\sigma(s)\rvert^2=\prod_\tau\lvert\sigma\tau(\cj{\piN})\rvert^{-2\brE{\tau}}}, and
\m{\lvert\sigma\tau(\cj{\piN})\rvert^{2}=\sigma\tau(\cj{\piN}\,\piN)=\ND}, because \m{\sigma\circ\tau} is an
embedding of the CM field \m{E(D)} and so commutes with complex conjugation. The matrix $B_D$ contributes at most
\m{\exp(x^{70})} by its size bound (Input~4) and the comparisons of sizes of the data (Chapter~III).
\end{proof}
\leanline{\lean{restrictToYofdiag}\sep \lean{restrictToY}\sep \lean{weylNormSqrestrictToYle}\sep In the first two
Lean takes the prime $\ell$, its primality and the homomorphism $\jm$ as instance arguments.}

\begin{proof}[Proof of Theorem D]
Write \m{\Sigma=\nonind{D}}, \m{\rho=\nrep{D}{0}}, \m{F=\Fop{\rho}}, \m{\Delta=\dpt{D}{0}}, and for a finite set
$\Sigma'$ of places \m{\mathsf{d}(\Sigma')=\emb{\Sigma}{\diagN{\XIo{0}}{\Sigma'}}^{-1}}; so
\m{\Delta\,\mathsf{d}(\PiD)=1}. For an index $a$ of \m{\XI{0}} let \m{t(a)} be the sum, over the classes of
\m{\Kcol{0}}, the coordinates $i$ of the weight and the infinite places $\tau$, of the number of times $a$ occurs in
the tableau that names the coordinate $i$ at~$\tau$.

\emph{Step 1.} For every $\Sigma'$: every finite family in \m{\Int{\Col}{\Delta\mathsf{d}(\Sigma')}} is carried into
\m{\Int{\Col}{\Delta}} by a square matrix over $\Oj$ whose determinant divides \m{\ND^{\,t(\st)}}. Indeed, a column is
determined by its values at representatives of the classes, which are arbitrary vectors of the weight because the
rational points act freely (Q4); adjust each representative by a rational point so that it, and its conjugate by
\m{\mathsf{d}(\Sigma')}, are integral above \m{\ND} (Q3). In these coordinates, by Lemma~\ref{l:nonind}\,(ii), a
column that is integral at \m{\Delta\mathsf{d}(\Sigma')} becomes integral at $\Delta$ after its coordinate $i$ at a
class is multiplied by \m{\ND} to the number of occurrences of the last index in the tableaux of~$i$: the action of
\m{\operatorname{diag}(1,\dots,1,\ND)} in the weight is diagonal with these powers of \m{\ND}, up to an integral
Schur matrix (Q3). The determinant of this diagonal map is \m{\ND^{\,t(\st)}}.

\emph{Step 2.} For \m{q\ge0}, every finite family in \m{\Int{\Col}{1}} is carried by a matrix over $\Oj$ of
determinant dividing \m{\ND^{\,t(\st)(q+1)}} into the forms $\psi$ with \m{F^{q}\psi\in\Int{\Col}{\Delta}}. For
\m{q=0} this is Step~1 with \m{\Sigma'=\PiD}. For the step it suffices, by Step~1 with \m{\Sigma'=\Sigma} and by
composing matrices, that $F$ maps \m{\Int{\Col}{\Delta}} into \m{\Int{\Col}{\Delta\mathsf{d}(\Sigma)}}. Let $\psi$ be
in the former. Then \m{F\psi} is a column (Lemma~\ref{l:ramcol}\,(ii)), and it is a sum of terms
\m{\psiN{y_{\st}}\cdot\rho(\uz\,\tilde k(y))\psi}; the gauge of a sum is at most the largest gauge of a term, and
\m{\lvert\jm(\psiN{y_{\st}})\rvert=1} since \m{\psiN{y_{\st}}} is a root of unity. The value of
\m{\rho(\uz\tilde k(y))\psi} at $g$ is \m{\psi(g')} with
\m{g'=g\,\Delta\,\mathsf{d}(\Sigma)\,\emb{\Sigma}{\tilde k(y)}\,\Delta^{-1}}, and
\m{g'\Delta=g\Delta\mathsf{d}(\Sigma)\cdot\emb{\Sigma}{\tilde k(y)}}, where the last factor is integral above \m{\ND}
(Lemma~\ref{l:embfacts}\,(ii)). So the gauge of this value at \m{g\Delta\mathsf{d}(\Sigma)} is its gauge at
\m{g'\Delta} (Q3), at most~$1$.

\emph{Step 3.} Take \m{q=\varrho(D)+1} and \m{e=t(\st)\,(\varrho(D)+2)}. To bound $e$: \m{t(a)=t(b)} for all
indices $a$,~$b$, by the symmetry of the weight under permutations of the coordinates; and
\m{\sum_at(a)} is the number of classes times the number of coordinates times the sum over $\tau$ of the number of
boxes of the weight at~$\tau$, which is at most \m{a(D)} for each~$\tau$ (unit \emph{Columns}). The number of
classes times the number of coordinates is \m{\dim_{\CC}\Col} (Q4). Hence
\m{t(\st)\le\dim_{\CC}\Col\cdot f(D)\,a(D)/n_0(D)}, and \m{e\le4\,m\,f(D)\,b\,\dim_{\CC}\Col} is one of the comparisons
of sizes of the data.
\end{proof}
\leanline{\lean{existsindexDvdfourierOppow}\sep The split level is an explicit argument of the Lean statement, with
Lemma~\ref{l:nonind}\,(i) as its default value.}

\begin{lemma}[the characters of the chain above the level prime]\label{l:chars}
In Setting~\ref{s:data}, let $\vartheta$ be a tame character of $D$ and, for \m{i\ge0}, let
\m{\mu_{D,i}(\vartheta^{i})} be the character of lift~$i$ (the product of the base character
\m{\mu^{0}_{D,i}} and \m{\vartheta^{i}}).
\begin{enumerate}[label=(\roman*)]
\item For \m{j\ge0}, a finite index set $I$ and a finite place $w$ of \m{E(D)} dividing \m{\ND}: the base change of
  the rank twist \m{\mu_{D,j}(\vartheta^{j})\,\upsilon^{\lfloor|I|/2\rfloor}} is unramified at~$w$.
\item For \m{j,i\ge0}: the space character, in rank \m{|\XI{j}|}, of the rank twist of
  \m{\mu_{D,i}(\vartheta^{i})} for \m{\XI{j}} is the base change \m{x\mapsto\mu_{D,i}(\vartheta^{i})(x/\cj{x})}.
\end{enumerate}
\end{lemma}
\begin{proof}
(i) The tame part of \m{\mu_{D,j}(\vartheta^{j})} is trivial at $w$ and at $\cj{w}$: these places are not chosen
(Input~5) and do not belong to the set \m{Q(\vartheta)}, whose places have norm prime to \m{\ND}; and $\upsilon$ is
unramified. Apply Q5. (ii) is Q5, as \m{|\XI{j}|=n_j(D)} is even (Input~5).
\end{proof}
\leanline{\lean{charUnramifiedAtrankTwistofNmem}\sep \lean{spaceCharrankTwist}}

\sect{Where the results are used}{s:used}
Theorem~A, Theorem~B, Lemma~\ref{l:names}\,(ii), Lemma~\ref{l:detchar} and Lemma~\ref{l:ramcol}\,(i) are used
together in the unit \emph{Columns}, in the file that trades denominators at the level prime for a column operator:
there the chain of pairs of lifts whose data are the translates of Theorem~B is rewritten, by Theorem~B and then
Theorem~A with \m{\rho_j=\nrep{D}{j}} and the laws given by the lifts of the definite chain, as a constant times the
chain with the simplest data applied to \m{\Fop{\nrep{D}{0}}^{\varrho(D)+1}}. Theorem~D and
Lemma~\ref{l:restr} enter the local statement at the level prime in the same unit. Theorem~C is used in
\emph{Columns} (the chain of a good datum is a plain sum over classes, which gives its integrality) and in the unit
\emph{Exchange}, together with Lemma~\ref{l:cheq}\,(i),\,(iii). Lemma~\ref{l:cheq}\,(ii), Lemma~\ref{l:names}\,(i)
and Lemma~\ref{l:chars}\,(i) are used in the note on the seesaw identity, to show that a central element at a
place dividing \m{\piN} acts trivially on a chain of pairs. The definitions of chains, of joint invariance and of
\m{\transum{I}{I'}{\Sigma}{N}{I''}} are part of the definition of a good datum (unit \emph{Columns}), and
\m{\twoact{\mathcal{K}}{a}{i}} occurs in the proof of the seesaw identity.

\sect{Notes on the formal text}{s:notes}
\subhead{Arguments that the census lists and Lean does not take.} The census of the workspace reads the arguments
of a declaration from its text. In the files treated here a section declares variables (a vector~$y$ of residues, a
family $a$ of points, data $\Phi$, a character~$b$); the text reading lists such a variable as an argument whenever
its name occurs in a statement, also when the occurrence is bound there. This concerns the vector $y$ in the
Gauss-sum operator, in the section of the split level, in Lemma~\ref{l:kofy}, Lemma~\ref{l:fnfacts}\,(iv),
Lemma~\ref{l:lawfacts}\,(i) and Theorem~A; the family $a$ in the definition of joint invariance; the data $\Phi$ in
Lemma~\ref{l:cheq}\,(iii); the character $b$ in Lemma~\ref{l:resid}. Nothing was compiled for this note, so Lean's
own answer is not recorded; the statements above follow the Lean text.

\subhead{Arguments that Lean takes and the census misses.} A lift law is defined for a split level~$\Sigma$ for a
prime~$N$, and so are the mirabolic group, the function $\nat{f}$ and the Gauss-sum operator; the statements of
Definition~\ref{d:law}, Lemma~\ref{l:mirfix}, Lemma~\ref{l:lawfacts}\,(ii) and Proposition~\ref{p:slide} therefore
take $N$, its primality and the split level, although their text does not name them. The embedding
\m{\emb{\Sigma}{g}} uses only that $\Sigma$ is disjoint from its conjugate, but takes a split level. The set
\m{\nonind{D}} and Lemma~\ref{l:restr}\,(i) take a prime $\ell$ and a homomorphism $\jm$ that their names do not show.

\subhead{A boundary case.} The adele \m{\rN{y}} is defined in the formal text for every natural number~$N$, with
no hypothesis on $\Sigma$. The description by the representative in \m{\{0,\dots,N-1\}} is the case \m{N>0}; for
\m{N=0} division by zero returns zero in Lean, so the adele is zero. Every use in this note has $N$ prime.

\subhead{Choices.} The section \m{\secN}, the map $k$ of Lemma~\ref{l:kofy} and hence the matrices \m{\tilde k(y)} in
the Gauss-sum operator are chosen once (by the axiom of choice). The proofs use only that \m{\red\circ\secN} is the
identity and that \m{k(y)} has last row and last column~$y$; a different choice changes \m{\Fop{\rho}} on vectors
that are not fixed by the principal congruence subgroup of level~$N$.

\subhead{Statements that are more special or more general than their use.} Lemma~\ref{l:detchar}\,(ii) compares
two ranks $I$,~$I'$ although the value does not depend on the rank. In Theorem~C the properties $P_j$ are arbitrary
and serve as an induction hypothesis carried along the chain; hypothesis~(b) (free classes) is what turns class
integrals into sums. In Theorem~D the exponent \m{\varrho(D)+1} is the number of pairs of lifts; the proof gives the
statement for every power~$q$ with the exponent \m{t(\st)(q+1)}. Lemma~\ref{l:names}\,(ii) is stated in the dossier
with the function \m{S\mapsto f(k^{-1}Sg)} of residue matrices; the Lean text has the function
\m{S\mapsto\nat{f}(k^{-1}Sg)} of matrices over $\ES{\Sigma}$, which is not of the form $\nat{f'}$ when $g$ or $k$ is
not integral, and this is the form used above.

\subhead{Members that act implicitly.} Some members of this note are rewriting rules that no proof cites by name:
the formula for the components of \m{\rN{y}} (Definition~\ref{d:split}) and the components of a scalar matrix
(Lemma~\ref{l:embfacts}). An unnamed instance (the reduction is a local homomorphism) is mentioned after
Lemma~\ref{l:elem}.

\subhead{How this note was written.} From the unit's dossier: the statements in words and the structured proofs of
the members of the two files (for the members that have a proof tree). Both Lean files were read completely
in the source at the pinned commit, and every statement above was compared with its Lean text argument by argument;
this is how the difference in Lemma~\ref{l:names}\,(ii) was found. The dossier's reader marks the claims for the long
proofs of Theorem~B and of Theorem~D as read from lemma names. Both proofs were read line by line in the Lean
source for this note; the computation with the two additive characters in Step~4 of Theorem~B and the three steps
of Theorem~D are written from that reading. The proofs of Theorem~A, Proposition~\ref{p:slide} and
Lemma~\ref{l:mirfix} were checked by hand as written here; they are shorter than the formal ones, which carry out
the same re-indexings through rewriting. Of the background library, the definitions of kernels and lifts, of gauges
and integral forms, of the Schrödinger action and of names were read in the source; their properties Q1--Q5 are
quoted, not read. The inputs from the unit \emph{Columns} are stated as the dossier's imports read them. The
program check of the unit's dossier covered \UnitSteps{} leaf steps of \UnitProofs{} proofs.

\head{Appendix A. Concordance with the Lean artifact}
{\fontsize{8.6}{11.6}\selectfont\renewcommand{\leanfont}{\ttfamily\fontsize{7.7}{10}\selectfont\color{leancol}}
\setlength{\LTleft}{0pt}\setlength{\LTright}{\fill}\setlength{\tabcolsep}{0pt}
\begin{longtable}{@{}p{27mm}>{\raggedright\arraybackslash}p{62mm}p{40mm}r@{}}
\textbf{In this note} & \textbf{Lean declaration} & \textbf{File} & \textbf{Lines}\\\midrule\endhead
\ConcordanceRows
\end{longtable}}

Legend: in the Lean text the set $\Sigma$ of places is {\leanfont P} and $\Sigma'$ is {\leanfont S}; the index
sets $I$, $I'$, $I''$ are {\leanfont n}, {\leanfont n'}, {\leanfont m}, and \m{\Ip{I}} is the sum of {\leanfont n}
and a one-point type; the representations $\rho$, $\rho'$ carry the Greek letter pi and the constant $\varkappa$
of a law the letter gamma; $\nat{f}$ is {\leanfont liftFn f}; $\rrow$, $\bcol{y}$, $\cchar{y}$ are
{\leanfont lastRowDatum}, {\leanfont lastColDatum y}, {\leanfont charDatum y}; $\uz$ is {\leanfont diagNInv}; the
ring homomorphism $\jm$ is the field of the class {\leanfont PadicPlace}.

Two statements as they stand in the artifact, the class of split levels and the Gauss-sum operator:
\quoted{IsSplitLevel}
\quoted{fourierOp}

\head{Notation}
\begingroup\fontsize{8.6}{10.9}\selectfont\raggedright
\setlength{\columnsep}{9mm}\setlength{\parskip}{0pt}
\begin{multicols}{2}
\nt{$E$, $K$}{a CM field with a CM type; in Setting~\ref{s:data}, \m{E=E(D)} over the imaginary quadratic field}
\nt{$\Af$, $\Oh$}{finite adeles of $E$; their integers}
\nt{$w$, $v$, $\cj{w}$}{finite places of $E$; the conjugate place}
\nt{$E_w$, $O_w$, $\mathbb{F}_w$, \m{\mathrm{N}w}}{completion, its integers, residue field, norm}
\nt{$\mathbf{j}$, $\alpha$, $\upsilon$}{the Weil data: an element, an idele, the splitting character}
\nt{$\psi_0$, $\psi_f$}{the standard finite additive characters}
\nt{$N$}{a prime number (the level prime \m{\ND} for the data)}
\nt{$\Sigma$, $\Sigma'$}{finite sets of finite places; split levels}
\nt{$\ES{\Sigma}$, $\OS{\Sigma}$, $\RS{\Sigma}$}{product of completions; integral elements; residue ring}
\nt{$\rN{y}$, $\psiN{y}$, $\tilde y$}{the adele \m{\tilde y/N} on $\Sigma$; its character value; least representative}
\nt{$\mathbf{1}_{\Sigma}$, $\epsilon$, $\epsilon'$}{the adele that is $1$ on a set of places; those of $\Sigma$, $\Sigma'$}
\nt{$I$, $I'$, $I''$, $I_j$, $J$}{finite index sets; \m{\Ip{I}=I\sqcup\{\st\}}, $\st$ the last index}
\nt{$\red$, $\secN$}{reduction modulo $N$; a chosen section on matrices}
\nt{$\ebas$}{the last basis vector}
\nt{$\diagN{I}{\Sigma'}$, $\uz$}{diagonal matrix with last entry $N$ at \m{\Sigma\cap\Sigma'}; the inverse for
  \m{\Sigma'=\Sigma}}
\nt{$\Mir{I}$, $\ram{\rho}$}{the mirabolic group; ramified vectors}
\nt{$\rho$, $\rho'$, $\rho_j$}{complex representations of \m{\GL_{\Ip{I}}(\ES{\Sigma})}}
\nt{$\varphi$, $\psi$}{vectors of a representation; algebraic forms}
\nt{$f$, $\nat{f}$, $h$, $f_t$}{functions of residue matrices; extension by zero}
\nt{$\rrow$, $\bcol{y}$, $\cchar{y}$}{indicator of the last row; of a last column; character of a last row}
\nt{$k(y)$, $\tilde k(y)$, $C$}{the matrix of Lemma~\ref{l:kofy}; its lift; an auxiliary matrix or a real bound}
\nt{$\Fop{\rho}$, $F$}{the Gauss-sum operator}
\nt{$L$, $L_j$, $\varkappa$, $\varkappa_j$}{lift laws and their constants}
\nt{$A_i$, $d_i$, $o_i$, $p$, $q$, $l$}{the pair operators, scalars, counts; numbers of pairs; an index}
\nt{$B_j^{q}$, $T$, $S$, $S'$, $R$, $r$, $t$, $c$, $x$, $y$, $z$, $u$, $a$, $b$, $g$, $k$}{auxiliary operators,
  matrices, residues, units and group elements, as introduced where used}
\nt{$\Gf{I}$, $\Kz{I}$, $X_I$}{finite adelic unitary group of the standard space $X_I$; stabiliser of its lattice}
\nt{$\kappa$, $e_j$, $\Wk{I}$, $\Alg{I}{e}$}{weight, last entries, space of the weight, algebraic forms}
\nt{$\mathcal{K}_j$, $\omega_j$, $\mathcal{J}_j$, $\mathcal{S}_j$, $K_j$}{kernels: representation, jet matrix, data
  space, level}
\nt{$\Phi$, $\Phi_x$, $\chain{\mathcal{K}}{\Phi}$, $\chaintwo{\mathcal{K}}{\Phi}$}{data; chains of lifts and of pairs}
\nt{$\twoact{\mathcal{K}}{a}{i}$}{the joint action on the data of a pair}
\nt{$M$, $\mathcal{X}$, $\mathcal{X}_0$, $L_x$, $A$, $P_j$, $\beta$, $G$, $\Gamma$}{in Theorem~C: target space,
  index set, linear maps, a subset, properties, a linear map, the integrand, rational points}
\nt{$\gamma$, $\delta$}{rational points, in the definition of a kernel; $\delta$ also a diagonal matrix}
\nt{$\cls{K}$, $m$, $n$, $i$, $j$}{number of classes; natural numbers}
\nt{$\emb{\Sigma}{g}$, $g_\Sigma$}{embedding at a split level; components at $\Sigma$}
\nt{$\eta$, $\detchar{\eta}{k}$, $\nu$, $\nu'$}{a Hecke character; the determinant character; norm-one adeles}
\nt{$V$, $\mathcal{H}$, $\mathcal{W}_f(V)$, $\varrho_f$}{hermitian spaces; finite part; finite Heisenberg operators}
\nt{$\projfin{V}{\Sigma'}$, $\nm{\Sigma}{F}$}{projection to vacua above \m{\Sigma'\cup\cj{\Sigma'}}; named operator}
\nt{$\transum{I}{I'}{\Sigma}{N}{I''}$}{the sum of pairs of translates with denominator $N$}
\nt{$\mathsf{p}$, $\zeta$, $Z$, $\Psi$, $X$}{in the proof of Theorem~B: an adele, a residue, tensors of functions,
  functions, a matrix}
\nt{$P=(N_0,d)$, $D$}{parameters; admissible data}
\nt{$x$, $x_1$, $b$, $\mu$, $m$, $r$, $g_0$, $g_1$, $\pi_g$}{components of $D$ (also local letters, as said)}
\nt{$\ND$, $\piN$, $\PiD$}{the level prime; \m{\piN\cj{\piN}=\ND}; the places dividing $\piN$}
\nt{$f(D)$, $a(D)$, $n_j(D)$, $\varrho(D)$}{number of infinite places; a size; ranks; \m{m=2\varrho(D)+2}}
\nt{$\XI{j}$, $\XIo{j}$, $X_j(D)$, $Y_m(D)$}{index sets; the definite spaces}
\nt{$\Stg{j}$, $\Kcol{j}$, $\Col$, $\Colram$}{forms of stage $j$; its level; columns; ramified columns}
\nt{$\nu_w$, $z_w$, $\mathfrak{p}$}{the central adele at a place of \m{\PiD}; its translation; such a place}
\nt{$\resY$, $\bdU$, $B_D$, $\brE{\tau}$, $s$, $\lambda_\tau$}{restriction; \m{h\mapsto\operatorname{diag}(h,1)};
  branching matrix; exponents; scalars}
\nt{$\tau$, $\sigma$}{an infinite place with its embedding; a ring homomorphism \m{\CC\to\CC}}
\nt{$\ell$, $\jm$, $\Oj$}{a prime; a ring homomorphism \m{\CC\to\CC_\ell}; its valuation ring in $\CC$}
\nt{$\ind{\tau}$, $\evt{\tau}$}{the induced prime; evaluation of adeles at its conjugate}
\nt{$\gauge{g}{u}$, $\Int{\mathcal{V}}{\mathsf{d}}$, $\wnorm{\cdot}$}{gauge; integral forms at a point; Gram norm}
\nt{$\nonind{D}$, $\dpt{D}{j}$, $\Delta$, $\delta$}{non-induced places; the diagonal point; a diagonal matrix}
\nt{$\cemb{D}{j}$, $\nrep{D}{j}$, $k^{\Delta}$}{conjugated embedding; the representation on stage $j$; conjugated
  components}
\nt{$\mathsf{d}$, $t(a)$, $e$}{a point; tableau counts; the exponent of Theorem~D}
\nt{$\vartheta$, $Q(\vartheta)$, $\mu_{D,i}(\vartheta^{i})$, $\mu^{0}_{D,i}$}{a tame character, its set of places,
  the characters of the lifts}
\end{multicols}
\endgroup

\end{document}
